\chapter{Neeraj Kumar [NRJ]: Complete Exercise Question Bank}
\label{chap:nrj_all_questions}

\begin{problembox}
\textbf{Problem 1} \hfill \textbf{[Neeraj Kumar [NRJ] EXERCISE I (JEE MAIN) Q1]}
\label{q:ext-nrj_all_questions-1}

A metal rod is dipped in a solution of its ions. Its 
electrode potential is  independent of
 
 (a) temperature of the solution
 
 (b) concentration of the solution
 
 (c) area of the metal exposed
 
 (d) nature of the metal
\end{problembox}

\begin{problembox}
\textbf{Problem 2} \hfill \textbf{[Neeraj Kumar [NRJ] EXERCISE I (JEE MAIN) Q2]}
\label{q:ext-nrj_all_questions-2}

Indicator electrode is
 
 (a) SHE
 
 (b) Calomel electrode
 
 (c) Ag/AgCl electrode
 
 (d) Quinhydrone electrode
\end{problembox}

\begin{problembox}
\textbf{Problem 3} \hfill \textbf{[Neeraj Kumar [NRJ] EXERCISE I (JEE MAIN) Q3]}
\label{q:ext-nrj_all_questions-3}

The position of some metals in the electrochemical 
series in decreasing electropositive character is 
given as: Mg > Al  > Zn > Cu > Ag. What will 
happen if a copper spoon is used to stir a solution 
of aluminium nitrate?
 
 (a)  The spoon will get coated with aluminium.
 
 (b)  An alloy of copper and aluminium is formed.
 
 (c) The solution becomes blue.
 
 (d) No chemical change will take place.
\end{problembox}

\begin{problembox}
\textbf{Problem 4} \hfill \textbf{[Neeraj Kumar [NRJ] EXERCISE I (JEE MAIN) Q4]}
\label{q:ext-nrj_all_questions-4}

Four colourless salt solutions are placed in 
separate test tubes and a strip of copper is placed 
in each. Which solution   nally turns blue?
 
 (a) AgNO3 
(b) Pb(NO3)2
 
 (c) Zn(NO3)2 
(d) Cd(NO3)2
\end{problembox}

\begin{problembox}
\textbf{Problem 5} \hfill \textbf{[Neeraj Kumar [NRJ] EXERCISE I (JEE MAIN) Q5]}
\label{q:ext-nrj_all_questions-5}

The metal that cannot be obtained on reduction 
of its oxide by aluminium is
 
 (a) K 
(b) Mn
 
 (c) Cr 
(d) Fe
\end{problembox}

\begin{problembox}
\textbf{Problem 6} \hfill \textbf{[Neeraj Kumar [NRJ] EXERCISE I (JEE MAIN) Q6]}
\label{q:ext-nrj_all_questions-6}

Beryllium is placed above magnesium in the group 
II. Beryllium dust, therefore, when added to 
MgCl2 solution will
 
 (a) have no e  ect
 
 (b) precipitate Mg metal
 
 (c) precipitate MgO
 
 (d) lead to dissolution of Be metal
\end{problembox}

\begin{problembox}
\textbf{Problem 7} \hfill \textbf{[Neeraj Kumar [NRJ] EXERCISE I (JEE MAIN) Q7]}
\label{q:ext-nrj_all_questions-7}

A standard hydrogen electrode has zero electrode 
potential because
 
 (a) hydrogen is easiest to oxidize
 
 (b)  this electrode potential is assumed to be zero
 
 (c) hydrogen atom has only one electron
 
 (d) hydrogen is the lightest element
\end{problembox}

\begin{problembox}
\textbf{Problem 8} \hfill \textbf{[Neeraj Kumar [NRJ] EXERCISE I (JEE MAIN) Q8]}
\label{q:ext-nrj_all_questions-8}

The standard reduction potential values of three 
metallic cations, X, Y and Z are +0.52, -3.03 and 
-1.18 V, respectively. The order of reducing power 
of the  corresponding metals is
 
 (a) Y > Z > X
 
 (b) X > Y > Z
 
 (c) Z > Y > X
 
 (d) Z > X > Y
EXERCISE I (JEE MAIN)
Electrochemistry
CHAPTER
8
 
 
Visit www.jeebooks.in for more
\end{problembox}

\begin{problembox}
\textbf{Problem 9} \hfill \textbf{[Neeraj Kumar [NRJ] EXERCISE I Q9]}
\label{q:ext-nrj_all_questions-9}

A gas 'X' at 1 atm is bubbled through a solution 
containing a mixture of 1 M -- Y- and 1 M -- Z- at 
25$/circ$C. If the reduction potential of Z- > Y- > X, 
then,
 
 (a) Y- will oxidize X and not Z-
 
 (b) Y- will oxidize Z- and not X
 
 (c) Y- will oxidize both X and Z-
 
 (d) Y- will reduce both X and Z-
\end{problembox}

\begin{problembox}
\textbf{Problem 10} \hfill \textbf{[Neeraj Kumar [NRJ] EXERCISE I Q10]}
\label{q:ext-nrj_all_questions-10}

Standard electrode potential data are useful for 
understanding the suitability of an oxidant in a 
redox titration. Some half-cell reactions and their 
standard potentials are given below:
 
   MnO4
-(aq) + 8H+(aq) + 5e- $/to$ Mn2+ (aq) + 
4H2O(l)  E $/circ$ = 1.51 V
 
   Cr2O7
2-(aq) + 14H+ (aq) + 6e- $/to$ 2Cr3+ (aq) + 
7H2O(l)  E $/circ$ = 1.38 V
 
   Fe3+ (aq) + e- $/to$ Fe2+ (aq) E $/circ$ = 0.77 V
 
  Cl2 (g) + 2e- $/to$ 2Cl- (aq) E $/circ$ = 1.40 V
 
 Identify the only incorrect statement regarding the 
quantitative estimation of aqueous Fe(NO3)2
 
 (a) MnO4
- can be used in aqueous HCl
 
 (b) Cr2O7
2-can be used in aqueous HCl
 
 (c) MnO4
- can be used in aqueous H2SO4
 
 (d) Cr2O7
2-can be used in aqueous H2SO4
\end{problembox}

\begin{problembox}
\textbf{Problem 11} \hfill \textbf{[Neeraj Kumar [NRJ] EXERCISE I Q11]}
\label{q:ext-nrj_all_questions-11}

The decreasing order of standard electrode 
potential of Mg, K, Ba and Ca, is
 
 (a) K, Ca, Ba, Mg 
(b) Ba, Ca, K, Mg
 
 (c) Ca, Mg, K, Ba 
(d) Mg, Ca, Ba, K
\end{problembox}

\begin{problembox}
\textbf{Problem 12} \hfill \textbf{[Neeraj Kumar [NRJ] EXERCISE I Q12]}
\label{q:ext-nrj_all_questions-12}

A metal having negative reduction potential, 
when dipped in the solution of its own ions, has a 
tendency to
 
 (a) remain as metal atoms
 
 (b) become electrically positive
 
 (c) become electrically negative
 
 (d) be deposited from the solution
\end{problembox}

\begin{problembox}
\textbf{Problem 13} \hfill \textbf{[Neeraj Kumar [NRJ] EXERCISE I Q13]}
\label{q:ext-nrj_all_questions-13}

The calomel electrode is reversible with respect to
 
 (a) Hg2
2+ 
(b) H+
 
 (c) Hg2+ 
(d) Cl-
\end{problembox}

\begin{problembox}
\textbf{Problem 14} \hfill \textbf{[Neeraj Kumar [NRJ] EXERCISE I Q14]}
\label{q:ext-nrj_all_questions-14}

Which one of the following does not get oxidized 
by bromine water?
 
 (a) Fe2+ to Fe3+ 
(b) Cu+ to Cu2+
 
 (c) Mn2+ to MnO4
- 
(d) Sn2+ to Sn4+
\end{problembox}

\begin{problembox}
\textbf{Problem 15} \hfill \textbf{[Neeraj Kumar [NRJ] EXERCISE I Q15]}
\label{q:ext-nrj_all_questions-15}

In salt bridge, normally KCl is used because
 
 (a) it is a strong electrolyte.
 
 (b) it is good conductor of electricity.
 
 (c)  K+ and Cl- ions have nearly same ionic 
mobility.
 
 (d) it is an ionic compound.
\end{problembox}

\begin{problembox}
\textbf{Problem 16} \hfill \textbf{[Neeraj Kumar [NRJ] EXERCISE I Q16]}
\label{q:ext-nrj_all_questions-16}

By how much would the oxidizing power of 
MnO4
-/Mn2+ couple change if the H+ ions 
concentration is decreased 100 times at 25$/circ$C?
 
 (a) increases by 189 mV
 
 (b) decreases by 189 mV
 
 (c) will increase by 19 mV
 
 (d) will decrease by 19 mV
\end{problembox}

\begin{problembox}
\textbf{Problem 17} \hfill \textbf{[Neeraj Kumar [NRJ] EXERCISE I Q17]}
\label{q:ext-nrj_all_questions-17}

The solution of CuSO4, in which copper rod is 
immersed, is diluted to 10 times. The reduction 
electrode potential
 
 (a) increases by 0.0295 V
 
 (b) decreases by 0.0295 V
 
 (c) increases by 0.059 V
 
 (d) decreases by 0.059 V
\end{problembox}

\begin{problembox}
\textbf{Problem 18} \hfill \textbf{[Neeraj Kumar [NRJ] EXERCISE I Q18]}
\label{q:ext-nrj_all_questions-18}

The standard reduction potential of oxygen in 
acidic solution is +1.23 V. What is the standard 
reduction potential of oxygen in basic solution?
 
 (a) +0.404 V 
(b) -0.404 V
 
 (c) +2.056 V 
(d) -2.056 V
\end{problembox}

\begin{problembox}
\textbf{Problem 19} \hfill \textbf{[Neeraj Kumar [NRJ] EXERCISE I Q19]}
\label{q:ext-nrj_all_questions-19}

The standard reduction potentials of Cu2+ | Cu and 
Cu2+ | Cu+ are 0.337 V and 0.153 V, respectively. 
The standard electrode potential of Cu+ | Cu half-
cell is
 
 (a) 0.184 V 
(b) 0.827 V
 
 (c) 0.521 V 
(d) 0.490 V
\end{problembox}

\begin{problembox}
\textbf{Problem 20} \hfill \textbf{[Neeraj Kumar [NRJ] EXERCISE I Q20]}
\label{q:ext-nrj_all_questions-20}

The electrode potential of hydrogen electrode in 
neutral solution and 298 K is
 
 (a) -0.413 V 
(b) zero
 
 (c) -0.826 V 
(d) +0.413 V
\end{problembox}

\begin{problembox}
\textbf{Problem 21} \hfill \textbf{[Neeraj Kumar [NRJ] EXERCISE I Q21]}
\label{q:ext-nrj_all_questions-21}

Electrode potential will be more for hydrogen 
electrode at pH (at the same temperature)
 
 (a) 4   (b) 3   (c) 2   (d) 5
\end{problembox}

\begin{problembox}
\textbf{Problem 22} \hfill \textbf{[Neeraj Kumar [NRJ] EXERCISE I Q22]}
\label{q:ext-nrj_all_questions-22}

Saturated solution of KNO3 is used to make 'salt 
bridge' because
 
 (a)  velocity of K+ is greater than that of NO3
-
 
 (b)  velocity of NO3
- is greater than that of K+
 
 (c)  velocities of both K+ and NO3
- are nearly the 
same
 
 (d) KNO3 is highly soluble in water
 
 
Visit www.jeebooks.in for more
\end{problembox}

\begin{problembox}
\textbf{Problem 23} \hfill \textbf{[Neeraj Kumar [NRJ] EXERCISE I Q23]}
\label{q:ext-nrj_all_questions-23}

The standard reduction potentials of Pt|Cr2O7
2-, 
Cr+3; Pt|MnO4
-, Mn+2; Pt|Ce+4, Ce+3 in the 
presence of acid are 1.33 V, 1.51 V and 1.61 V, 
respectively, at 25$/circ$C. The decreasing order of 
oxidizing power is
 
 (a) Cr2O7
2- > MnO4
- > Ce+4
 
 (b) MnO4
- > Cr2O7
2-
 > Ce+4
 
 (c) Ce+4 > MnO4
- > Cr2O7
2-
 
 (d) MnO4
- > Ce+4 > Cr2O7
2-
\end{problembox}

\begin{problembox}
\textbf{Problem 24} \hfill \textbf{[Neeraj Kumar [NRJ] EXERCISE I Q24]}
\label{q:ext-nrj_all_questions-24}

The standard reduction potentials at 25$/circ$C of 
Li+|Li, Ba2+|Ba, Na+|Na and Mg2+|Mg are -3.05, 
-2.73, -2.71 and -2.37  V, respectively. Which is 
the strongest reducing agent?
 
 (a) Li 
(b) Ba
 
 (c) Na 
(d) Mg
\end{problembox}

\begin{problembox}
\textbf{Problem 25} \hfill \textbf{[Neeraj Kumar [NRJ] EXERCISE I Q25]}
\label{q:ext-nrj_all_questions-25}

Some standard electrode potentials are given:
 
  Fe2+ + 2e- $/to$ Fe;  E $/circ$ = -0.440 V
 
  Fe3+ + 3e- $/to$ Fe;  E $/circ$ = -0.036 V
 
 The standard electrode potential for: Fe3+ + e- $/to$ 
Fe2+, is
 
 (a) -0.476 V 
(b) -0.404 V
 
 (c) +0.988 V 
(d) +0.772 V
\end{problembox}

\begin{problembox}
\textbf{Problem 26} \hfill \textbf{[Neeraj Kumar [NRJ] EXERCISE I Q26]}
\label{q:ext-nrj_all_questions-26}

The dissociation constant for CH3COOH is 1.8 
$/times$ 10-5 at 298 K. The electrode potential for the 
half-cell: Pt | H2 (1  bar) | 0.5  M -- CH3COOH, 
at 298 K is (log 2 = 0.3; log 3 = 0.48; 2.303 RT/F 
= 0.06)
 
 (a) -0.3024 V 
(b) -0.1512 V
 
 (c) +0.3024 V 
(d) +0.1512 V
\end{problembox}

\begin{problembox}
\textbf{Problem 27} \hfill \textbf{[Neeraj Kumar [NRJ] EXERCISE I Q27]}
\label{q:ext-nrj_all_questions-27}

At 25$/circ$C, the solubility product of CuCl is 2.0 
$/times$ 10-7 and ECl  CuCl  Cu
-
$/circ$
|
|
 
is 0.128 V. The value of 
ECu
Cu
+
$/circ$
|
 is (log 2 = 0.3; 2.303 RT/F =  0.06)
 
 (a) +1.08 V 
(b) +0.53 V
 
 (c) +1.682 V 
(d) +0.878 V
\end{problembox}

\begin{problembox}
\textbf{Problem 28} \hfill \textbf{[Neeraj Kumar [NRJ] EXERCISE I Q28]}
\label{q:ext-nrj_all_questions-28}

At 25$/circ$C, the solubility product of Pb(OH)2 (s) is 
[Given: E
E
Pb
Pb
OH
Pb(OH) Pb
V;
V
2
2
0 13
0 55
+
-
$/circ$
$/circ$
= -
= -
|
|
|
.
.
and 2 303
0 06
.
.
RT
F
=
]
 
 (a) 1.36 $/times$ 10--5 
(b) 1.0 $/times$ 10--7
 
 (c) 1.0 $/times$ 10--14 
(d) 1.25 $/times$ 10--15
\end{problembox}

\begin{problembox}
\textbf{Problem 29} \hfill \textbf{[Neeraj Kumar [NRJ] EXERCISE I Q29]}
\label{q:ext-nrj_all_questions-29}

The standard potentials of MnO4
-  |  Mn2+ and 
MnO2  |  Mn2+ electrodes in acid solution are 
1.51 and 1.23 V, respectively. Standard electrode 
potential for the electrode, MnO4
- | MnO2 in acid 
solution is
 
 (a) +1.697 V 
(b) +5.09 V
 
 (c) +0.28 V 
(d) +1.37 V
\end{problembox}

\begin{problembox}
\textbf{Problem 30} \hfill \textbf{[Neeraj Kumar [NRJ] EXERCISE I Q30]}
\label{q:ext-nrj_all_questions-30}

The following reactions represent the reduction of 
IO3
- ion into I- ion in acidic and basic medium. 
Predict in which medium IO3
- ion will act as a 
better oxidizing agent?
 
 IO3
- + 6H+ + 6e- $/to$ I- + 3H2O; E $/circ$ = +0.907 V
 
 IO3
- + 3H2O + 6e- $/to$ I- + 6OH-; E $/circ$ = +0.260 V
 
 (a) Acid medium 
(b) Basic medium
 
 (c) Equally in both 
(d) Not predictable
Galvanic Cell
\end{problembox}

\begin{problembox}
\textbf{Problem 31} \hfill \textbf{[Neeraj Kumar [NRJ] EXERCISE I Q31]}
\label{q:ext-nrj_all_questions-31}

The correct cell diagram for the following reaction 
and E   for the cell is
 
  2AgBr(s) + H2(g) $/to$ 2Ag(s) + 2H+ + 2Br-
 
  E $/circ$AgBr|Ag|Br- = +0.10 V.
 
 (a)  (Pt) H2 | H+ || Br- | AgBr | Br2 (Pt); E $/circ$ = 0.10 V
 
 (b)  (Pt) H2 | H+ || Br- | AgBr | Br2 (Pt); E $/circ$ = -0.10 V
 
 (c)  (Pt) Br2 | AgBr | Br- || H+ | H2 (Pt); E $/circ$ = 0.10 V
 
 (d)  (Pt) Br2 | AgBr | Br- || H+ | H2 (Pt); E $/circ$ = -0.10 V
\end{problembox}

\begin{problembox}
\textbf{Problem 32} \hfill \textbf{[Neeraj Kumar [NRJ] EXERCISE I Q32]}
\label{q:ext-nrj_all_questions-32}

In an experimental set-up for the measurement of 
EMF of a half-cell using a reference electrode and 
a salt bridge, when the salt bridge is removed, the 
voltage
 
 (a) remains the same
 
 (b) increases to maximum
 
 (c) decreases half the value
 
 (d) drops to zero
\end{problembox}

\begin{problembox}
\textbf{Problem 33} \hfill \textbf{[Neeraj Kumar [NRJ] EXERCISE I Q33]}
\label{q:ext-nrj_all_questions-33}

After some time, the voltage of an electrochemical 
cell becomes zero. This is because
 
 (a)  their electrode potential becomes zero.
 
 (b)  their reduction potential become equal but 
have opposite sign.
 
 (c)  their reduction potential become equal and 
have the same sign.
 
 (d)  the ions of the electrolyte in the salt bridge 
stop moving.
 
 
Visit www.jeebooks.in for more
\end{problembox}

\begin{problembox}
\textbf{Problem 34} \hfill \textbf{[Neeraj Kumar [NRJ] EXERCISE I Q34]}
\label{q:ext-nrj_all_questions-34}

The cell reaction for the given cell is spontaneous 
if Pt, Cl2 (P1 atm) | Cl- | Cl2 (P2 atm), Pt
 
 (a) P1 > P2 
(b) P1 < P2
 
 (c) P1 = P2 
(d) P1 = 1 atm
\end{problembox}

\begin{problembox}
\textbf{Problem 35} \hfill \textbf{[Neeraj Kumar [NRJ] EXERCISE I Q35]}
\label{q:ext-nrj_all_questions-35}

Which one of the following statements is incorrect 
regarding an electrochemical cell?
 
 (a)  The electrode on which oxidation takes place 
is called anode.
 
 (b) Anode is the negative pole.
 
 (c)  The direction of the current is same as that of 
the direction of   ow of electrons.
 
 (d)  The   ow of current is partly due to   ow of 
electrons and partly due to   ow of ions.
\end{problembox}

\begin{problembox}
\textbf{Problem 36} \hfill \textbf{[Neeraj Kumar [NRJ] EXERCISE I Q36]}
\label{q:ext-nrj_all_questions-36}

When an electric current is drawn from a galvanic 
cell
 
 (a) EMF suddenly increases.
 
 (b)  EMF gradually increases and attains a 
maximum value.
 
 (c)  EMF decreases and   nally falls to zero.
 
 (d) EMF must remain constant.
\end{problembox}

\begin{problembox}
\textbf{Problem 37} \hfill \textbf{[Neeraj Kumar [NRJ] EXERCISE I Q37]}
\label{q:ext-nrj_all_questions-37}

Identi  cation of anode and cathode in an 
electrochemical cell is made by the use of
 
 (a) Galvanometer 
(b) Salt bridge
 
 (c) Voltmeter 
(d) Potentiometer
\end{problembox}

\begin{problembox}
\textbf{Problem 38} \hfill \textbf{[Neeraj Kumar [NRJ] EXERCISE I Q38]}
\label{q:ext-nrj_all_questions-38}

For the cell Zn | Zn2+ || Cu2+ | Cu, if the concentration 
of both, Zn2+ and Cu2+ ions are doubled, the EMF 
of the cell
 
 (a) doubles 
(b) reduces to half
 
 (c) remains same 
(d) becomes zero
\end{problembox}

\begin{problembox}
\textbf{Problem 39} \hfill \textbf{[Neeraj Kumar [NRJ] EXERCISE I Q39]}
\label{q:ext-nrj_all_questions-39}

The standard EMF of a galvanic cell can be 
calculated from
 
 (a) the size of the electrode
 
 (b) the pH of the solution
 
 (c) the amount of metal in the anode
 
 (d) the E $/circ$ values of the half-cells
\end{problembox}

\begin{problembox}
\textbf{Problem 40} \hfill \textbf{[Neeraj Kumar [NRJ] EXERCISE I Q40]}
\label{q:ext-nrj_all_questions-40}

The value of equilibrium constant for a feasible 
cell reaction must be
 
 (a) <1 
(b) Zero
 
 (c) = 1 
(d) >1
\end{problembox}

\begin{problembox}
\textbf{Problem 41} \hfill \textbf{[Neeraj Kumar [NRJ] EXERCISE I Q41]}
\label{q:ext-nrj_all_questions-41}

If the cell reaction is spontaneous then
 
 (a) $/Delta$G $/circ$ = -ve 
(b) E $/circ$red = -ve
 
 (c) E $/circ$red = +ve 
(d) $/Delta$G = -ve
\end{problembox}

\begin{problembox}
\textbf{Problem 42} \hfill \textbf{[Neeraj Kumar [NRJ] EXERCISE I Q42]}
\label{q:ext-nrj_all_questions-42}

Which of the following pair of metals, when 
coupled, will give maximum EMF for a voltaic cell?
 
 (a) Fe and Cu 
(b) Pb and Au
 
 (c) Cu and Au 
(d) Ca and Cu
\end{problembox}

\begin{problembox}
\textbf{Problem 43} \hfill \textbf{[Neeraj Kumar [NRJ] EXERCISE I Q43]}
\label{q:ext-nrj_all_questions-43}

When a lead storage battery is charged
 
 (a) PbO2 dissolves.
 
 (b)  the lead electrode becomes coated with lead 
sulphate.
 
 (c) sulphuric acid is regenerated.
 
 (d) the amount of acid decreases.
\end{problembox}

\begin{problembox}
\textbf{Problem 44} \hfill \textbf{[Neeraj Kumar [NRJ] EXERCISE I Q44]}
\label{q:ext-nrj_all_questions-44}

Use of lithium metal as an electrode in high energy 
density batteries is due to
 
 (a) lithium is the lightest element.
 
 (b)  lithium has the highest oxidation potential.
 
 (c) lithium is quite reactive.
 
 (d) lithium does not corrode readily.
\end{problembox}

\begin{problembox}
\textbf{Problem 45} \hfill \textbf{[Neeraj Kumar [NRJ] EXERCISE I Q45]}
\label{q:ext-nrj_all_questions-45}

A depolarizer used in dry cell batteries is
 
 (a) NH4Cl 
(b) MnO2
 
 (c) KOH 
(d) Na3PO4
\end{problembox}

\begin{problembox}
\textbf{Problem 46} \hfill \textbf{[Neeraj Kumar [NRJ] EXERCISE I Q46]}
\label{q:ext-nrj_all_questions-46}

Which is correct about fuel cells?
 
 (a)  Cell continuously run as long as fuels are 
supplied.
 
 (b)  These are more e   cient and free from 
pollution.
 
 (c)  These are used to provide power and drinking 
water to astronauts in space programme.
 
 (d) All of these
\end{problembox}

\begin{problembox}
\textbf{Problem 47} \hfill \textbf{[Neeraj Kumar [NRJ] EXERCISE I Q47]}
\label{q:ext-nrj_all_questions-47}

When a lead storage battery is discharged
 
 (a) SO2 is evolved
 
 (b) lead sulphate is consumed
 
 (c) lead is formed
 
 (d) sulphuric acid is consumed
\end{problembox}

\begin{problembox}
\textbf{Problem 48} \hfill \textbf{[Neeraj Kumar [NRJ] EXERCISE I Q48]}
\label{q:ext-nrj_all_questions-48}

For a cell reaction involving a two- electron 
change, the standard EMF of the cell is found to 
be 0.295 V at 25$/circ$C. The equilibrium constant of 
the reaction at 25$/circ$C will be
 
 (a) 1 $/times$ 1010 
(b) 1 $/times$ 10-10
 
 (c) 29.5 $/times$ 10-2 
(d) 2 $/times$ 1010
\end{problembox}

\begin{problembox}
\textbf{Problem 49} \hfill \textbf{[Neeraj Kumar [NRJ] EXERCISE I Q49]}
\label{q:ext-nrj_all_questions-49}

The ECell for Ag(s) | AgI (satd) || Ag+ (0.10  M) | Ag 
(s) is +0.413 V. What is the value of Ksp of AgI?
 
 (a) 1.0 $/times$ 10-8 
(b) 1.0 $/times$ 10-7
 
 (c) 1.0 $/times$ 10-14 
(d) 1.0 $/times$ 10-16
 
 
Visit www.jeebooks.in for more
\end{problembox}

\begin{problembox}
\textbf{Problem 50} \hfill \textbf{[Neeraj Kumar [NRJ] EXERCISE I Q50]}
\label{q:ext-nrj_all_questions-50}

Assuming that hydrogen behaves as an ideal gas, 
what is the EMF of the cell at 25$/circ$C if P1 = 600 mm 
and P2 = 420 mm: Pt | H2 (P1) | HCl | H2 (P2) | Pt? 
[Given: 2.303 RT/F = 0.06, log 7 = 0.85]
 
 (a) -0.0045 V 
(b) -0.0 V
 
 (c) +0.0045 V 
(d) +0.0015 V
\end{problembox}

\begin{problembox}
\textbf{Problem 51} \hfill \textbf{[Neeraj Kumar [NRJ] EXERCISE I Q51]}
\label{q:ext-nrj_all_questions-51}

For the electrochemical cell, M | M+ || X-  |  X, 
EM
M
+
$/circ$
| = 0.44 V and EX X
|
-
$/circ$
= 0.33 V. From this data, 
one can deduce that
 
 (a)  M + X $/to$ M+ + X- is the spontaneous reaction
 
 (b)  M+ + X- $/to$ M + X is the spontaneous reaction
 
 (c) Ecell = 0.77 V
 
 (d) Ecell = - 0.77 V
\end{problembox}

\begin{problembox}
\textbf{Problem 52} \hfill \textbf{[Neeraj Kumar [NRJ] EXERCISE I Q52]}
\label{q:ext-nrj_all_questions-52}

The EMF of the cell: Zn | Zn2+ (0.01 M) || Fe2+ 
(0.001 M) | Fe at 298 K is 0.2905 V, then the value 
of equilibrium constant for the cell reaction is
 
 (a) e
0 32
0 0295
.
.
 
(b) 10
0 32
0 0295
.
.
 
 (c) 10
0 26
0 0295
.
.
 
(d) 10
0 32
0 0591
.
.
\end{problembox}

\begin{problembox}
\textbf{Problem 53} \hfill \textbf{[Neeraj Kumar [NRJ] EXERCISE I Q53]}
\label{q:ext-nrj_all_questions-53}

For the cell reaction: 4Br- + O2 + 4H+   2Br2 + 
2H2O; E $/circ$ = 0.18 V. The value of (log KC) at 298 K 
is [2.303 RT/F = 0.06]
 
 (a) 12 
(b) 6
 
 (c) 18 
(d) 3
\end{problembox}

\begin{problembox}
\textbf{Problem 54} \hfill \textbf{[Neeraj Kumar [NRJ] EXERCISE I Q54]}
\label{q:ext-nrj_all_questions-54}

The standard EMF for the cell reaction: Zn(s) + 
Cu2+ (aq) $/to$ Zn2+ (aq) + Cu(s) is 1.10 volts at 25$/circ$C. 
The EMF of the cell reaction when 0.1 M Cu2+ 
and 0.1 M Zn2+ solutions are used at 25$/circ$C is
 
 (a) 1.10 V 
(b) 1.041 V
 
 (c) -1.10 V 
(d) -1.041 V
\end{problembox}

\begin{problembox}
\textbf{Problem 55} \hfill \textbf{[Neeraj Kumar [NRJ] EXERCISE I Q55]}
\label{q:ext-nrj_all_questions-55}

For the cell: Ni | Ni2+ || Cu2+ | Cu; E $/circ$ = 0.77 V. By 
which of the following activity, Ecell will increase?
 
 (a) On decreasing [Ni+2]
 
 (b) On decreasing [Cu+2]
 
 (c) On increasing mass of Ni electrode
 
 (d) On increasing mass of Cu electrode
\end{problembox}

\begin{problembox}
\textbf{Problem 56} \hfill \textbf{[Neeraj Kumar [NRJ] EXERCISE I Q56]}
\label{q:ext-nrj_all_questions-56}

The standard EMF of a Daniel cell at 298 K is 
E1. When the concentration of ZnSO4 is 1.0 M 
and that of CuSO4 is 0.01 M, the EMF becomes 
E2 at 298 K. The correct relationship between E1 
and E2 is
 
 (a) E1 = E2 
(b) E2 = 0
 
 (c) E1 > E2 
(d) E1 < E2
\end{problembox}

\begin{problembox}
\textbf{Problem 57} \hfill \textbf{[Neeraj Kumar [NRJ] EXERCISE I Q57]}
\label{q:ext-nrj_all_questions-57}

For a reaction A(s) + 2B+ $/to$ A2+ + B(s); KC has 
been found to be 1012. The EMF of the cell is
 
 (a) 0.354 V 
(b) 0.708 V
 
 (c) 0.534 V 
(d) 0.453 V
\end{problembox}

\begin{problembox}
\textbf{Problem 58} \hfill \textbf{[Neeraj Kumar [NRJ] EXERCISE I Q58]}
\label{q:ext-nrj_all_questions-58}

An electrochemical cell is set up as  follows:
 
   Pt | H2(g) (1 atm) | 0.001 M HCl || 0.1 M HA | 
H2(g) (1 atm) | Pt
 
   EMF of this cell is zero because (pKa of HA = 5)
 
 (a)  molar concentrations of acids are di  erent
 
 (b) Temperature is constant
 
 (c) pH of two solutions are same
 
 (d)  both are standard hydrogen  electrodes
\end{problembox}

\begin{problembox}
\textbf{Problem 59} \hfill \textbf{[Neeraj Kumar [NRJ] EXERCISE I Q59]}
\label{q:ext-nrj_all_questions-59}

The reaction: H2(g) + 2AgCl(s) $/to$ 2H+(aq)  + 
2Cl-(aq) + 2Ag(s) occurs in the galvanic cell
 
 (a)  Ag | AgCl (s) | KCl(aq) | AgNO3(aq) | Ag
 
 (b) Pt | H2(g) | HCl(aq) | AgNO3(aq) | Ag
 
 (c) Pt | H2(g) | HCl(aq) | AgCl(s) | Ag
 
 (d) Pt | H2(g) | KCl(aq) | AgCl(s) | Ag
\end{problembox}

\begin{problembox}
\textbf{Problem 60} \hfill \textbf{[Neeraj Kumar [NRJ] EXERCISE I Q60]}
\label{q:ext-nrj_all_questions-60}

The standard reduction potentials in acidic 
conditions are 0.77 V and 0.53 V, respectively, for 
Fe3+ | Fe2+ and I3
- | I- couples. The equilibrium 
constant for the reaction: 2Fe3+ + 3I-   2Fe2+ + 
I3
-, is (2.303 RT/F = 0.06)
 
 (a) 2 $/times$ 108 
(b) 108
 
 (c) 104 
(d) 10-8
\end{problembox}

\begin{problembox}
\textbf{Problem 61} \hfill \textbf{[Neeraj Kumar [NRJ] EXERCISE I Q61]}
\label{q:ext-nrj_all_questions-61}

The following electrochemical cell has been set up: 
Pt(s)|Fe3+, Fe2+ (a = 1) || Ce4+, Ce3+ (a = 1) | Pt(s); 
E $/circ$(Fe3+|Fe2+) = 0.77 V; E $/circ$ (Ce4+|Ce3+) = 1.61 V. 
If an ammeter is connected between the two 
platinum electrodes, predict the direction of   ow 
of current. Will the current increase or decrease 
with time?
 
 (a) Ce electrode to Fe electrode, decrease
 
 (b) Ce electrode to Fe electrode, increase
 
 (c) Fe electrode to Ce electrode, decrease
 
 (d) Fe electrode to Ce electrode, increase
\end{problembox}

\begin{problembox}
\textbf{Problem 62} \hfill \textbf{[Neeraj Kumar [NRJ] EXERCISE I Q62]}
\label{q:ext-nrj_all_questions-62}

For the reaction: H2 (1 bar) + 2AgCl(s)   2Ag(s) 
+ 2H+ (0.1 M) + 2Cl- (0.1 M); $/Delta$G  0 = - 48,250 J at 
25$/circ$C. The EMF of cell in which the given reaction 
takes place is
 
 (a) 0.25 V 
(b) 0.37 V
 
 (c) 0.13 V 
(d) 0.0.49 V
 
 
Visit www.jeebooks.in for more
\end{problembox}

\begin{problembox}
\textbf{Problem 63} \hfill \textbf{[Neeraj Kumar [NRJ] EXERCISE I Q63]}
\label{q:ext-nrj_all_questions-63}

A cell contains two hydrogen  electrodes. The 
negative electrode is in contact with a solution 
of 10-6 M hydrogen ions. The EMF of the cell is 
0.118 V at 25$/circ$C. The concentration of hydrogen 
ions at the positive electrode is
 
 (a) 10-6 M 
(b) 10-3 M
 
 (c) 10-4 M 
(d) 10-5 M
\end{problembox}

\begin{problembox}
\textbf{Problem 64} \hfill \textbf{[Neeraj Kumar [NRJ] EXERCISE I Q64]}
\label{q:ext-nrj_all_questions-64}

The standard potentials of OCl-/Cl- and Cl-/Cl2 
are 0.94 V and -1.36 V, respectively. The E  value 
of OCl-/Cl2 will be
 
 (a) 3.24 V 
(b) -0.42 V
 
 (c) -2.30 V 
(d) 0.52 V
\end{problembox}

\begin{problembox}
\textbf{Problem 65} \hfill \textbf{[Neeraj Kumar [NRJ] EXERCISE I Q65]}
\label{q:ext-nrj_all_questions-65}

From the following E $/circ$ values for the half-cells:
 
   (i) D $/to$ D2+ + 2e-; E $/circ$ = -1.5 V
 
  (ii) B+ + e- $/to$ B; E $/circ$ = - 0.5 V
 
 (iii) A3- $/to$ A2- + e-; E $/circ$ = 1.5 V
 
 (iv) C2+ + e- $/to$ C+; E $/circ$ = +0.5 V
 
 Which combination of two half-cells would result 
in a cell with largest  potential?
 
 (a) i and iii 
(b) i and iv
 
 (c) iii and iv 
(d) ii and iv
Electrolysis
\end{problembox}

\begin{problembox}
\textbf{Problem 66} \hfill \textbf{[Neeraj Kumar [NRJ] EXERCISE I Q66]}
\label{q:ext-nrj_all_questions-66}

Which of the following statements does not 
di  erentiate between electrochemical cell and 
electrolytic cell?
 
 (a)  Spontaneous or non-spontaneous nature of 
the chemical process.
 
 (b)  Chemical reactions occurring at the electrodes
 
 (c) Positive and negative nature of anode.
 
 (d) EMF measurement.
\end{problembox}

\begin{problembox}
\textbf{Problem 67} \hfill \textbf{[Neeraj Kumar [NRJ] EXERCISE I Q67]}
\label{q:ext-nrj_all_questions-67}

The electrode through which electrons enter the 
electrolytic solution is
 
 (a) cathode
 
 (b) anode
 
 (c) may be anode or cathode
 
 (d) both, anode and cathode
\end{problembox}

\begin{problembox}
\textbf{Problem 68} \hfill \textbf{[Neeraj Kumar [NRJ] EXERCISE I Q68]}
\label{q:ext-nrj_all_questions-68}

Which process occurs in the electrolysis of an 
aqueous solution of nickel chloride at nickel anode?
 
 (a) Ni $/to$ Ni2+ + 2e- 
(b) Ni2+ + 2e- $/to$ Ni
 
 (c) 2Cl- $/to$ Cl2 + 2e- 
(d) 2H+ + 2e- $/to$ H2
\end{problembox}

\begin{problembox}
\textbf{Problem 69} \hfill \textbf{[Neeraj Kumar [NRJ] EXERCISE I Q69]}
\label{q:ext-nrj_all_questions-69}

If mercury is used as a cathode during the 
electrolysis of an aqueous NaCl solution, the ions 
discharged at cathode are
 
 (a) H+ 
(b) Na+
 
 (c) OH- 
(d) Cl-
\end{problembox}

\begin{problembox}
\textbf{Problem 70} \hfill \textbf{[Neeraj Kumar [NRJ] EXERCISE I Q70]}
\label{q:ext-nrj_all_questions-70}

Electrochemical equivalent is more for
 
 (a) Hydrogen 
(b) Silver
 
 (c) Copper 
(d) Zinc
\end{problembox}

\begin{problembox}
\textbf{Problem 71} \hfill \textbf{[Neeraj Kumar [NRJ] EXERCISE I Q71]}
\label{q:ext-nrj_all_questions-71}

The electrolytic bath used in gold plating of copper 
articles contains
 
 (a) Molten gold
 
 (b) Copper sulphate(aq)
 
 (c) AuCl3(aq)
 
 (d) AuCl3 + NaCN(aq)
\end{problembox}

\begin{problembox}
\textbf{Problem 72} \hfill \textbf{[Neeraj Kumar [NRJ] EXERCISE I Q72]}
\label{q:ext-nrj_all_questions-72}

Copper can be deposited from acidi  ed copper 
sulphate and alkaline cuprous cyanide. If the same 
current is passed for a de  nite time
 
 (a)  the amount of copper deposited from acidic 
copper sulphate will be higher
 
 (b)  the amount of copper deposited from alkaline 
cuprous cyanide will be higher
 
 (c)  the same amount of copper will be deposited
 
 (d) copper will not deposit in either case
\end{problembox}

\begin{problembox}
\textbf{Problem 73} \hfill \textbf{[Neeraj Kumar [NRJ] EXERCISE I Q73]}
\label{q:ext-nrj_all_questions-73}

In the electrolytic cell,   ow of electrons is from
 
 (a) Cathode to anode in solution
 
 (b)  Cathode to anode through external supply
 
 (c)  Cathode to anode through internal supply
 
 (d)  Anode to cathode through internal supply
\end{problembox}

\begin{problembox}
\textbf{Problem 74} \hfill \textbf{[Neeraj Kumar [NRJ] EXERCISE I Q74]}
\label{q:ext-nrj_all_questions-74}

Electrolytic cell is used to convert
 
 (a) Chemical energy to electrical energy
 
 (b) Electrical energy to chemical energy
 
 (c)  Chemical energy to mechanical energy
 
 (d)  Electrical energy to mechanical energy
\end{problembox}

\begin{problembox}
\textbf{Problem 75} \hfill \textbf{[Neeraj Kumar [NRJ] EXERCISE I Q75]}
\label{q:ext-nrj_all_questions-75}

Faraday's law of electrolysis fails when
 
 (a) temperature is increased
 
 (b) inert electrodes are used
 
 (c) a mixture of electrolytes is used
 
 (d) in none of these cases
\end{problembox}

\begin{problembox}
\textbf{Problem 76} \hfill \textbf{[Neeraj Kumar [NRJ] EXERCISE I Q76]}
\label{q:ext-nrj_all_questions-76}

Using same quantity of current, which among Na, 
Mg and Al is deposited more (by mass) during 
electrolysis of their molten salts?
 
 (a) Na 
(b) Mg
 
 (c) Al 
(d) All in same
 
 
Visit www.jeebooks.in for more
\end{problembox}

\begin{problembox}
\textbf{Problem 77} \hfill \textbf{[Neeraj Kumar [NRJ] EXERCISE I Q77]}
\label{q:ext-nrj_all_questions-77}

A certain current liberated 0.50 g of hydrogen in 
2 h. How many grams of copper can be liberated 
by the same current   owing for the same time in a 
copper sulphate solution? (Cu = 63.5)
 
 (a) 12.7 g 
(b) 15.88 g
 
 (c) 31.75 g 
(d) 63.5 g
\end{problembox}

\begin{problembox}
\textbf{Problem 78} \hfill \textbf{[Neeraj Kumar [NRJ] EXERCISE I Q78]}
\label{q:ext-nrj_all_questions-78}

A current of 3.7 A is passed for 6 h between nickel 
electrodes in 0.50 L of 2 M solution of Ni(NO3)2. 
The molarity of Ni2+ at the end of electrolysis is
 
 (a) 1.172 M 
(b) 0.172 M
 
 (c) 0.586 M 
(d) 2 M
\end{problembox}

\begin{problembox}
\textbf{Problem 79} \hfill \textbf{[Neeraj Kumar [NRJ] EXERCISE I Q79]}
\label{q:ext-nrj_all_questions-79}

The current e   ciency of an electrodeposition of 
copper metal in which 9.8 g of copper is deposited 
by a current of 3  A for 10000 s, from aqueous 
copper sulphate solution, is about
 
 (a) 60/% 
(b) 99/%
 
 (c) 92/% 
(d) 75/%
\end{problembox}

\begin{problembox}
\textbf{Problem 80} \hfill \textbf{[Neeraj Kumar [NRJ] EXERCISE I Q80]}
\label{q:ext-nrj_all_questions-80}

On passing electricity through dilute H2SO4 
solution, the mass of substances liberated at the 
cathode and anode are in the ratio of
 
 (a) 1:8 
(b) 8:1
 
 (c) 1:32 
(d) 1:16
\end{problembox}

\begin{problembox}
\textbf{Problem 81} \hfill \textbf{[Neeraj Kumar [NRJ] EXERCISE I Q81]}
\label{q:ext-nrj_all_questions-81}

The electrochemical equivalents of two substances 
are E1 and E2. The current that must pass to 
deposit the same amount at the cathodes in the 
same time must be in the ratio of
 
 (a) E1:E2 
(b) E2:E1
 
 (c) (E1 - E2):E2 
(d) E1:(E2 - E1)
\end{problembox}

\begin{problembox}
\textbf{Problem 82} \hfill \textbf{[Neeraj Kumar [NRJ] EXERCISE I Q82]}
\label{q:ext-nrj_all_questions-82}

The same quantity of electricity is passed through 
one molar solution of H2SO4 and one molar 
solution of HCl. The amount of hydrogen evolved 
from H2SO4 as compared to that from HCl is
 
 (a) the same
 
 (b) twice as such
 
 (c) one half as such
 
 (d) dependent on size of electrode
\end{problembox}

\begin{problembox}
\textbf{Problem 83} \hfill \textbf{[Neeraj Kumar [NRJ] EXERCISE I Q83]}
\label{q:ext-nrj_all_questions-83}

In the electrolysis of acidi  ed AgNO3 solution 
using Pt-electrodes, the anode reaction is
 
 (a) 2NO3
- $/to$ 2NO2 + O2 + 2e-
 
 (b) NO3
- $/to$ NO + 1
2 O2 + e-
 
 (c) 2H2O $/to$ 4H+ + O2 + 4e-
 
 (d) Pt $/to$ Pt3+ + 3e-
\end{problembox}

\begin{problembox}
\textbf{Problem 84} \hfill \textbf{[Neeraj Kumar [NRJ] EXERCISE I Q84]}
\label{q:ext-nrj_all_questions-84}

Two platinum electrodes were immersed in a 
solution of CuSO4 and electric current was passed 
through the solution. After some time, it was 
found that colour of CuSO4 disappeared with the 
evolution of gas at the electrode. The colourless 
solution contains
 
 (a) platinum sulphate
 
 (b) copper sulphate
 
 (c) copper hydroxide
 
 (d) sulphuric acid
\end{problembox}

\begin{problembox}
\textbf{Problem 85} \hfill \textbf{[Neeraj Kumar [NRJ] EXERCISE I Q85]}
\label{q:ext-nrj_all_questions-85}

A solution containing 1.0 M each of Cu(NO3)2, 
Mg(NO3)2, 
AgNO3, 
Hg(NO3)2 
is 
being 
electrolysed using inert electrodes. The values of 
standard electrode potential are: Ag+ | Ag = 0.80 V, 
Hg2+ | Hg  = 0.79 V, Cu2+ | Cu = 0.34 V, Mg2+ | Mg 
= - 2.37 V. With increasing voltage, the sequence 
of deposit of metals on the cathode will be
 
 (a) Ag, Hg, Cu, Mg 
 
 (b) Mg, Cu, Hg, Ag
 
 (c) Ag, Mg, Cu 
 
 (d) Cu, Hg, Ag
\end{problembox}

\begin{problembox}
\textbf{Problem 86} \hfill \textbf{[Neeraj Kumar [NRJ] EXERCISE I Q86]}
\label{q:ext-nrj_all_questions-86}

During the electrolysis of an aqueous salt 
solution, the pH in the space near one of the 
electrode was increased and the other one was 
decreased. The salt solution was
 
 (a) NaCl (very dilute) 
(b) ZnCl2
 
 (c) NaCl (Conc.) 
(d) Cu(NO3)2
\end{problembox}

\begin{problembox}
\textbf{Problem 87} \hfill \textbf{[Neeraj Kumar [NRJ] EXERCISE I Q87]}
\label{q:ext-nrj_all_questions-87}

Two electrolytic cells, one containing acidi  ed 
ferrous chloride and another acidi  ed ferric 
chloride, are connected in series. The mass ratio 
of iron deposited at cathodes in the two cells 
will be
 
 (a) 3:1 
(b) 2:3
 
 (c) 1:1 
(d) 3:2
\end{problembox}

\begin{problembox}
\textbf{Problem 88} \hfill \textbf{[Neeraj Kumar [NRJ] EXERCISE I Q88]}
\label{q:ext-nrj_all_questions-88}

A galvanic cell is set up from a zinc bar weighing 
100 g and 1.0 L of 1.0 M copper sulphate solution. 
How long would the cell run if it is assumed to 
deliver a steady current of 1.0 A? (Zn = 65.4)
 
 (a) 53.6 h 
(b) 26.8 h
 
 (c) 81.97 h 
(d) 40.99 h
\end{problembox}

\begin{problembox}
\textbf{Problem 89} \hfill \textbf{[Neeraj Kumar [NRJ] EXERCISE I Q89]}
\label{q:ext-nrj_all_questions-89}

An ion is a reduced to the element when it absorbs 
6 $/times$ 1020 electrons. The number of equivalents of 
the ion is
 
 (a) 0.10 
(b) 0.01
 
 (c) 0.001 
(d) 0.0001
 
 
Visit www.jeebooks.in for more
\end{problembox}

\begin{problembox}
\textbf{Problem 90} \hfill \textbf{[Neeraj Kumar [NRJ] EXERCISE I Q90]}
\label{q:ext-nrj_all_questions-90}

In the lead storage battery, the anode  reaction is 
Pb(s) + HSO4
- + H2O $/to$ PbSO4(s) + H3O+ + 2e-. 
How many grams of Pb will be used up to deliver 
1 A for 100 h? (Pb = 208)
 
 (a) 776 g 
(b) 388 g
 
 (c) 194 g 
(d) 0.1 g
\end{problembox}

\begin{problembox}
\textbf{Problem 91} \hfill \textbf{[Neeraj Kumar [NRJ] EXERCISE I Q91]}
\label{q:ext-nrj_all_questions-91}

The copper anode of a cell containing silver nitrate 
solution weighs 60.0 g. After passing current 
for some time, it is found that 3.24 g of silver if 
deposited on the platinum cathode. What is the 
  nal weight of the anode? (Ag = 108, Cu = 64)
 
 (a) 0.96 g 
(b) 60 g
 
 (c) 59.04 g 
(d) 60.96 g
\end{problembox}

\begin{problembox}
\textbf{Problem 92} \hfill \textbf{[Neeraj Kumar [NRJ] EXERCISE I Q92]}
\label{q:ext-nrj_all_questions-92}

The quantity of electricity required for the 
reduction of 1 mole of Fe2O3 to Fe is
 
 (a) 1F 
(b) 0.33F
 
 (c) 3F 
(d) 6F
\end{problembox}

\begin{problembox}
\textbf{Problem 93} \hfill \textbf{[Neeraj Kumar [NRJ] EXERCISE I Q93]}
\label{q:ext-nrj_all_questions-93}

Three faradays of electricity is passed through molten 
Al2O3, aqueous solutions of CuSO4 and molten 
NaCl. The amounts of Al, Cu and Na deposited at 
the cathodes will be in the molar ratio of
 
 (a) 1:2:3 
(b) 3:2:1
 
 (c) 1:1.5:3 
(d) 6:3:2
\end{problembox}

\begin{problembox}
\textbf{Problem 94} \hfill \textbf{[Neeraj Kumar [NRJ] EXERCISE I Q94]}
\label{q:ext-nrj_all_questions-94}

Electrolysis of a solution of HSO4
- ions produces 
S2O8
2-. 
Assuming 
75/% 
current 
e   ciency, 
what current should be employed to achieve a 
production rate of 1 mole of S2O8
2- per hour?
 
 (a) 71.5 A 
(b) 35.7 A
 
 (c) 53.0 A 
(d) 143 A
\end{problembox}

\begin{problembox}
\textbf{Problem 95} \hfill \textbf{[Neeraj Kumar [NRJ] EXERCISE I Q95]}
\label{q:ext-nrj_all_questions-95}

The number of Faradays required to produce 
1 g-atom of Mg from MgCl2 is
 
 (a) 1 
(b) 2
 
 (c) 0.5 
(d) 4
\end{problembox}

\begin{problembox}
\textbf{Problem 96} \hfill \textbf{[Neeraj Kumar [NRJ] EXERCISE I Q96]}
\label{q:ext-nrj_all_questions-96}

Passage of 96,500 coulomb of electricity liberates 
\underline{\hspace{1cm}} L of O2 at 273$/circ$C and 2  atm during 
electrolysis.
 
 (a) 5.6 
(b) 16.8
 
 (c) 22.4 
(d) 11.2
\end{problembox}

\begin{problembox}
\textbf{Problem 97} \hfill \textbf{[Neeraj Kumar [NRJ] EXERCISE I Q97]}
\label{q:ext-nrj_all_questions-97}

An electrolytic cell contains a solution of Ag2SO4 
and platinum electrodes. A current is passed 
until 1.6 g of O2 has been liberated at anode. The 
amount of silver deposited at cathode would be
 
 (a) 108.0 g 
(b) 1.6 g
 
 (c) 10.8 g 
(d) 21.6 g
\end{problembox}

\begin{problembox}
\textbf{Problem 98} \hfill \textbf{[Neeraj Kumar [NRJ] EXERCISE I Q98]}
\label{q:ext-nrj_all_questions-98}

When 12,000 coulombs of electricity is passed 
through the electrolyte, 3.0 g of a metal of atomic 
mass 96.5 g/mol is deposited. The electro-valency 
of the metal cation in the electrolyte is
 
 (a) +4 
(b) +3
 
 (c) +2 
(d) -4
\end{problembox}

\begin{problembox}
\textbf{Problem 99} \hfill \textbf{[Neeraj Kumar [NRJ] EXERCISE I Q99]}
\label{q:ext-nrj_all_questions-99}

How many electrons   ow when a current of 5 A is 
passed through a solution for 200 s?
 
 (a) 6.022 $/times$ 1023 
(b) 6.24 $/times$ 1021
 
 (c) 6.024 $/times$ 1021 
(d) 6.022 $/times$ 1020
\end{problembox}

\begin{problembox}
\textbf{Problem 100} \hfill \textbf{[Neeraj Kumar [NRJ] EXERCISE I Q100]}
\label{q:ext-nrj_all_questions-100}

The current of 9.65 A   owing for 10 min deposits 
3.0 g of a metal. The equivalent weight of the 
metal is
 
 (a) 10 
(b) 30
 
 (c) 50 
(d) 96.5
\end{problembox}

\begin{problembox}
\textbf{Problem 101} \hfill \textbf{[Neeraj Kumar [NRJ] EXERCISE I Q101]}
\label{q:ext-nrj_all_questions-101}

If a current of 1.0 A is drawn from the Daniel cell 
for 96.5 min, the cathode will gain in weight by 
(Cu = 63.5, Zn = 65.4)
 
 (a) 1.905 g 
(b) 1.962 g
 
 (c) 3.81 g 
(d) 3.924 g
\end{problembox}

\begin{problembox}
\textbf{Problem 102} \hfill \textbf{[Neeraj Kumar [NRJ] EXERCISE I Q102]}
\label{q:ext-nrj_all_questions-102}

The current required to produce oxygen at the 
rate of 2.8 ml  (0$/circ$C, 1 atm) per second during 
electrolysis of acidulated water is
 
 (a) 48.25 A/s 
(b) 24.12 A/s
 
 (c) 96.5 A/s 
(d) 0.0048 A/s
\end{problembox}

\begin{problembox}
\textbf{Problem 103} \hfill \textbf{[Neeraj Kumar [NRJ] EXERCISE I Q103]}
\label{q:ext-nrj_all_questions-103}

Sodium amalgam is prepared by  electrolysis of 
aqueous NaCl using 10 g mercury as cathode. 
How many Faraday of electricity is required to 
prepare 18.7/% Na-amalgam, by weight, with a 
current e   ciency of 50/%?
 
 (a) 0.1 F 
(b) 0.2 F
 
 (c) 0.05 F 
(d) 0.16 F
\end{problembox}

\begin{problembox}
\textbf{Problem 104} \hfill \textbf{[Neeraj Kumar [NRJ] EXERCISE I Q104]}
\label{q:ext-nrj_all_questions-104}

Element A (atomic mass = 112) and element B 
(atomic mass = 27) form chlorides. Solutions of 
these chlorides are electrolysed separately and it is 
found that when the same quantity of electricity 
is passed, 5.6 g of A was deposited while only 0.9 
g of B was deposited. The valency of B is 3. The 
valency of A is
 
 (a) 1 
(b) 2
 
 (c) 3 
(d) 4
 
 
Visit www.jeebooks.in for more
\end{problembox}

\begin{problembox}
\textbf{Problem 105} \hfill \textbf{[Neeraj Kumar [NRJ] EXERCISE I Q105]}
\label{q:ext-nrj_all_questions-105}

The same current was passed successively through 
solution 
of 
zinc-ammonium 
sulphate 
and 
nickel-ammonium sulphate rendered alkaline 
with ammonia. The weights of zinc and nickel 
deposited in a certain time were found to be 
22.89 g and 20.55 g, respectively. Given that the 
chemical equivalent weight of zinc is 32.7, what is 
the chemical equivalent weight of nickel?
 
 (a) 58.71 
(b) 29.36
 
 (c) 14.39 
(d) 36.42
Conductance
\end{problembox}

\begin{problembox}
\textbf{Problem 106} \hfill \textbf{[Neeraj Kumar [NRJ] EXERCISE I Q106]}
\label{q:ext-nrj_all_questions-106}

Which of the following solutions have highest 
resistance?
 
 (a) 1N - NaCl 
(b) 0.05N - NaCl
 
 (c) 2N - NaCl 
(d) 0.1N -- NaCl
\end{problembox}

\begin{problembox}
\textbf{Problem 107} \hfill \textbf{[Neeraj Kumar [NRJ] EXERCISE I Q107]}
\label{q:ext-nrj_all_questions-107}

Variation of molar conductance of an electrolytic 
solution with temperature is that it
 
 (a)  increases with increase of  temperature
 
 (b)  decreases with increase of  temperature
 
 (c)   rst increases then decreases
 
 (d) is not a  ected by temperature
\end{problembox}

\begin{problembox}
\textbf{Problem 108} \hfill \textbf{[Neeraj Kumar [NRJ] EXERCISE I Q108]}
\label{q:ext-nrj_all_questions-108}

Which pure substance will not conduct electricity?
 
 (a) Molten NaCl 
(b) Molten KOH
 
 (c) Lique  ed HCl 
(d) Liquid Hg
\end{problembox}

\begin{problembox}
\textbf{Problem 109} \hfill \textbf{[Neeraj Kumar [NRJ] EXERCISE I Q109]}
\label{q:ext-nrj_all_questions-109}

The correct order of molar conductance at in  nite 
dilution of LiCl, NaCl and KCl is
 
 (a) LiCl > NaCl > KCl 
(b) KCl > NaCl > LiCl
 
 (c) NaCl > KCl > LiCl 
(d) LiCl > KCl > NaCl
\end{problembox}

\begin{problembox}
\textbf{Problem 110} \hfill \textbf{[Neeraj Kumar [NRJ] EXERCISE I Q110]}
\label{q:ext-nrj_all_questions-110}

The molar conductance of a strong electrolyte at 
in  nite dilution
 
 (a)  tends to a   nite value, which is above that at 
higher concentration
 
 (b)  tends to a   nite value, which is below that at 
higher concentration
 
 (c) tends to zero
 
 (d)  tends to a   nite value, which is equal to that at 
high concentration
\end{problembox}

\begin{problembox}
\textbf{Problem 111} \hfill \textbf{[Neeraj Kumar [NRJ] EXERCISE I Q111]}
\label{q:ext-nrj_all_questions-111}

The best conductor of electricity is a 0.1  M 
solution of
 
 (a) Boric acid 
(b) Sulphuric acid
 
 (c) Acetic acid 
(d) Propanoic acid
\end{problembox}

\begin{problembox}
\textbf{Problem 112} \hfill \textbf{[Neeraj Kumar [NRJ] EXERCISE I Q112]}
\label{q:ext-nrj_all_questions-112}

The speci  c conductance of AgCl solution in 
water was determined to be 1.8 $/times$ 10-6 $/Omega$-1 cm-1 
at 298 K. The molar conductances at in  nite 
dilution, of Ag+ and Cl- are 67.9 and 82.1 $/Omega$-1 cm2 
mol-1, respectively. What is the solubility of AgCl 
in water?
 
 (a) 1.2 $/times$ 10-8 M 
(b) 1.44 $/times$ 10-10 M
 
 (c) 1.2 $/times$ 10-5 M 
(d) 1.44 $/times$ 10-16 M
\end{problembox}

\begin{problembox}
\textbf{Problem 113} \hfill \textbf{[Neeraj Kumar [NRJ] EXERCISE I Q113]}
\label{q:ext-nrj_all_questions-113}

Equivalence conductance at in  nite dilution 
of NH4Cl, NaOH and NaCl are 129.8, 217.4 
and 108.9 $/Omega$-1 cm2 mol-1, respectively. If the 
equivalent conductance of 0.01 N solution of 
NH4OH is 9.532 $/Omega$-1 cm2 mol-1, then the degree 
of dissociation of NH4OH at this temperature is
 
 (a) 0.04/% 
(b) 2.1/%
 
 (c) 4.0/% 
(d) 44.7/%
\end{problembox}

\begin{problembox}
\textbf{Problem 114} \hfill \textbf{[Neeraj Kumar [NRJ] EXERCISE I Q114]}
\label{q:ext-nrj_all_questions-114}

The resistance of 1 M -- CH3COOH solution is 250 
$/Omega$, when measured in a cell of cell constant 125 
m-1. The molar conductivity, in $/Omega$-1 m2 mol-1 is
 
 (a) 5.0 $/times$ 10-4 
(b) 500
 
 (c) 2 $/times$ 10-3 
(d) 200
\end{problembox}

\begin{problembox}
\textbf{Problem 115} \hfill \textbf{[Neeraj Kumar [NRJ] EXERCISE I Q115]}
\label{q:ext-nrj_all_questions-115}

How does the electrical conductivity of 20 ml of 
0.2 M -- MgSO4 change when 0.5 M -- Ba(OH)2 
solution is gradually added in it, to excess?
 
 (a) decreases continuously
 
 (b) increases continuously
 
 (c) increases and then decreases
 
 (d) decreases and then increases
\end{problembox}

\begin{problembox}
\textbf{Problem 116} \hfill \textbf{[Neeraj Kumar [NRJ] EXERCISE I Q116]}
\label{q:ext-nrj_all_questions-116}

The equivalent conductivity (in $/Omega$-1 cm2 eq-1) of 
1.0 M -- H2SO4 solution of speci  c conductance 
2.6 $/times$ 10-1 cm-1, is
 
 (a) 1.3 $/times$ 102 
(b) 6.5 $/times$ 101
 
 (c) 1.3 $/times$ 10-1 
(d) 2.6 $/times$ 102
\end{problembox}

\begin{problembox}
\textbf{Problem 117} \hfill \textbf{[Neeraj Kumar [NRJ] EXERCISE I Q117]}
\label{q:ext-nrj_all_questions-117}

The molar conductance of a 0.01 M solution of 
acetic acid was found to be 16.30 $/Omega$-1 cm-1 mol-1 
at 25$/circ$C. The ionic conductances of hydrogen 
and acetate ions at in  nite dilution are 349.8 and 
40.9  $/Omega$-1 cm-1 mol-1, respectively, at the same 
temperature. What percentage of acetic acid is 
dissociated at this concentration?
 
 (a) 0.04172/% 
(b) 4.172/%
 
 (c) 41.72/% 
(d) 0.4172/%
 
 
Visit www.jeebooks.in for more
\end{problembox}

\begin{problembox}
\textbf{Problem 118} \hfill \textbf{[Neeraj Kumar [NRJ] EXERCISE I Q118]}
\label{q:ext-nrj_all_questions-118}

The distance between two electrodes of a cell is 2.5 
cm and area of each electrode is 5 cm2. The cell 
constant is
 
 (a) 0.5 m-1 
(b) 12.5 cm3
 
 (c) 2.0 cm 
(d) 50 m-1
\end{problembox}

\begin{problembox}
\textbf{Problem 119} \hfill \textbf{[Neeraj Kumar [NRJ] EXERCISE I Q119]}
\label{q:ext-nrj_all_questions-119}

The molar conductivity of NH4Cl, OH- and Cl- 
at in  nite dilution is 150, 200 and 75 $/Omega$-1 cm2 
mol-1, respectively. If the molar conductivity of 
a 0.01 M-NH4OH solution is 22 $/Omega$-1 cm2 mol-1, 
then its degree of dissociation is
 
 (a) 0.146 
(b) 0.063
 
 (c) 0.080 
(d) 0.293
\end{problembox}

\begin{problembox}
\textbf{Problem 120} \hfill \textbf{[Neeraj Kumar [NRJ] EXERCISE I Q120]}
\label{q:ext-nrj_all_questions-120}

Calculate the ionic product of water at 25$/circ$C from 
the following data:
 
   Conductivity of water = 5.5 $/times$ 10-6 mho m-1
 
  $/lambda$$/circ$H
+ = 0.035 mho m2 mol-1
 
  $/lambda$$/circ$OH
- = 0.020 mho m2 mol-1
 
 (a) 2 $/times$ 10-14 M2 
(b) 1 $/times$ 10-7 M
 
 (c) 1 $/times$ 10-8 M2 
(d) 1 $/times$ 10-14 M2
\end{problembox}

\begin{problembox}
\textbf{Problem 121} \hfill \textbf{[Neeraj Kumar [NRJ] EXERCISE I Q121]}
\label{q:ext-nrj_all_questions-121}

Calculate Ka of acetic acid if its 0.05 M solution has 
molar conductivity of 7.814 $/times$ 10-4 $/Omega$-1 m2 mol-1 
at 25$/circ$C. Given: $/Lambda$m$/circ$ for CH3COOH = 3.907 $/times$ 10-2 
$/Omega$-1 m2 mol-1.
 
 (a) 2 $/times$ 10-5 
(b) 1.8 $/times$ 10-5
 
 (c) 4 $/times$ 10-4 
(d) 0.02
\end{problembox}

\begin{problembox}
\textbf{Problem 122} \hfill \textbf{[Neeraj Kumar [NRJ] EXERCISE I Q122]}
\label{q:ext-nrj_all_questions-122}

Equal volumes of 0.015 M -- CH3COOH and 
0.015 M -- NaOH solutions are mixed together. 
What would be the molar conductivity of mixture 
if conductivity of CH3COONa is 6.3 $/times$ 10-4 
S cm-1?
 
 (a) 0.84 S cm2 mol-1 
(b) 8.4 S cm2 mol-1
 
 (c) 84 S cm2 mol-1 
(d) 42 S cm2 mol-1
\end{problembox}

\begin{problembox}
\textbf{Problem 123} \hfill \textbf{[Neeraj Kumar [NRJ] EXERCISE I Q123]}
\label{q:ext-nrj_all_questions-123}

Calculate $/Lambda$m
$/infty$(in $/Omega$-1 cm2 mol-1) for SrCl2 at 25$/circ$C, 
from the following data:
Conc.
0.25 M
1.0 M
$/Lambda$m (in $/Omega$-1 cm2 mol-1)
260
250
 
 (a) 270 
(b) 265
 
 (c) 240 
(d) 275
\end{problembox}

\begin{problembox}
\textbf{Problem 124} \hfill \textbf{[Neeraj Kumar [NRJ] EXERCISE I Q124]}
\label{q:ext-nrj_all_questions-124}

The resistance of a solution A is 50 $/Omega$ and that 
of solution B is 100 $/Omega$, both solutions being taken 
in the same conductivity cell. If equal volumes 
of solution A and B are mixed, what will be the 
resistance of the mixture using the same cell? 
Assume that there is no increase in the degree of 
dissociation of A and B on mixing.
 
 (a) 150 $/Omega$ 
(b) 75 $/Omega$
 
 (c) 33.33 $/Omega$ 
(d) 66.67 $/Omega$
\end{problembox}

\begin{problembox}
\textbf{Problem 125} \hfill \textbf{[Neeraj Kumar [NRJ] EXERCISE I Q125]}
\label{q:ext-nrj_all_questions-125}

In a conductivity cell, the two platinum electrodes, 
each of area 10 cm2 are   xed 1.5 cm apart. The 
cell contained 0.05 N solution of a salt. If the two 
electrodes are just half dipped into the solution 
which has a resistance of 50 $/Omega$, the equivalent 
conductance of the salt solution, in $/Omega$-1 cm2 eq-1, is
 
 (a) 120 
(b) 60
 
 (c) 240 
(d) 3000
 
 
Visit www.jeebooks.in for more
\end{problembox}

\begin{problembox}
\textbf{Problem 126} \hfill \textbf{[Neeraj Kumar [NRJ] EXERCISE II (JEE ADVANCED) Q1]}
\label{q:ext-nrj_all_questions-126}

A student made the following observations in the 
laboratory:
 
    (I)  Clean copper metal did not react with 1 M -- 
Pb(NO3)2 solution.
 
   (II)  Clean lead metal dissolves in 1 M --  AgNO3 
solution and crystals of Ag metal appeared.
 
 (III)  Clean silver metal did not react with 1 M -- 
Cu(NO3)2 solution.
 
 The order of decreasing reducing character of the 
three metals is
 
 (a) Cu > Pb > Ag 
(b) Cu > Ag > Pb
 
 (c) Pb > Cu > Ag 
(d) Pb > Ag > Cu
\end{problembox}

\begin{problembox}
\textbf{Problem 127} \hfill \textbf{[Neeraj Kumar [NRJ] EXERCISE II (JEE ADVANCED) Q2]}
\label{q:ext-nrj_all_questions-127}

We 
have 
an 
oxidation-reduction 
system: 
[Fe(CN)6]3- + e-   [Fe(CN)6]4-; E  = +0.36 V. The 
ratio of concentrations of oxidized and reduced 
from at which the potential of the system becomes 
0.24 V, is [Given: 2.303 RT/F = 0.06)
 
 (a) 2:1 
(b) 1:2
 
 (c) 1:20 
(d) 1:100
\end{problembox}

\begin{problembox}
\textbf{Problem 128} \hfill \textbf{[Neeraj Kumar [NRJ] EXERCISE II (JEE ADVANCED) Q3]}
\label{q:ext-nrj_all_questions-128}

The standard reduction potential for the process: 
[Co(H2O)6]3+ + e- $/to$ [Co(H2O)6]2+ is 1.8 V. The 
standard reduction potential for the process: 
[Co(NH3)6]3+ + e- $/to$ [Co(NH3)6]2+ is 0.1 V. Which 
of the complex ion, [Co(H2O)6]2+ or [Co(NH3)6]2+ 
can be oxidized to the corresponding cobalt (III) 
complex, by oxygen, in basic medium, under 
standard condition? [Given: EO OH
2 |
-
$/circ$
 
= 0.4 V]
 
 (a) [Co(H2O)6]2+ 
(b) [Co(NH3)6]2+
 
 (c) both 
(d) none of these
\end{problembox}

\begin{problembox}
\textbf{Problem 129} \hfill \textbf{[Neeraj Kumar [NRJ] EXERCISE II (JEE ADVANCED) Q4]}
\label{q:ext-nrj_all_questions-129}

The 
standard 
reduction 
potential 
for 
the 
reactions: Ag+ + e- $/to$ Ag and Ag(NH3)2
+  + e- $/to$ 
Ag + 2NH3 are +0.79 V and +0.37 V, respectively. 
From these values and the Nernst equation, what 
should be Kf for the Ag(NH3)2
+ ion? [Given: 2.303
RT/F = 0.06]
 
 (a) 1.0 $/times$ 10-7 
(b) 1.0 $/times$ 107
 
 (c) 2.15 $/times$ 1019 
(d) 4.64 $/times$ 10-20
\end{problembox}

\begin{problembox}
\textbf{Problem 130} \hfill \textbf{[Neeraj Kumar [NRJ] EXERCISE II (JEE ADVANCED) Q5]}
\label{q:ext-nrj_all_questions-130}

The overall formation constant for the reaction of 
6 mole of CN- with cobalt (II) is 1 $/times$ 1019. What is 
the formation constant for the reaction of 6 moles 
of CN- with cobalt (III)? Given that
 
   Co(CN)6
3- + e- $/to$ Co (CN)6
4- ; E $/circ$ = -0.83 V
 
   Co3+ + e- $/to$ Co2+; E $/circ$ = +1.81 V
 
  and 2.303 RT/F = 0.06.
 
 (a) 1.0 $/times$ 1063 
(b) 1.0 $/times$ 1025
 
 (c) 1.0 $/times$ 10-25 
(d) 1.0 $/times$ 10-63
\end{problembox}

\begin{problembox}
\textbf{Problem 131} \hfill \textbf{[Neeraj Kumar [NRJ] EXERCISE II (JEE ADVANCED) Q6]}
\label{q:ext-nrj_all_questions-131}

For the process: Cu2+ + 2e- $/to$ Cu; log[Cu2+] vs. 
Ered graph is shown in the   gure, where OA = 
0.34 V. The electrode potential of the half-cell of 
Cu|Cu2+ (0.1 M) will be [2.303 RT/F = 0.06]
A
Ered
O
log[Cu2+]
 
 (a) -0.31 V 
(b) +0.31 V
 
 (c) -0.37 V 
(d) +0.37 V
\end{problembox}

\begin{problembox}
\textbf{Problem 132} \hfill \textbf{[Neeraj Kumar [NRJ] EXERCISE II (JEE ADVANCED) Q7]}
\label{q:ext-nrj_all_questions-132}

Tell which of the following statements accurately 
describes the e  ect of adding CN- to the cathode 
of a cell with a cell reaction: Cd + 2Ag+ $/to$ 2Ag + 
Cd2+, E $/circ$ = 1.2 V
 
 (a)  E $/circ$ increases because Cd(CN)4
2- forms
 
 (b)  E $/circ$ decreases because Cd(CN)4
2- forms
 
 (c) E $/circ$ increases because Ag(CN)2
- forms
 
 (d)  E $/circ$ decreases because Ag(CN)2
- forms
\end{problembox}

\begin{problembox}
\textbf{Problem 133} \hfill \textbf{[Neeraj Kumar [NRJ] EXERCISE II (JEE ADVANCED) Q8]}
\label{q:ext-nrj_all_questions-133}

For a electrochemical cell Zn | Zn2+ (C1 M) || Cu2+ 
(C2 M) | Cu, the decrease in free energy at a given 
temperature is a function of
 
 (a) ln C1 
(b) ln C2
 
 (c) ln C2$/cdot$C1 
(d) ln C1/C2
\end{problembox}

\begin{problembox}
\textbf{Problem 134} \hfill \textbf{[Neeraj Kumar [NRJ] EXERCISE II (JEE ADVANCED) Q9]}
\label{q:ext-nrj_all_questions-134}

Select the correct option if it is known that Ksp 
(AgCl) > Ksp (AgBr) > Ksp (AgI)
 
 (a) EI
AgI Ag
-$/circ$
|
|
 > EBr
AgBr Ag
-
$/circ$
|
|
 > ECl
AgCl Ag
-
$/circ$
|
|
 
 (b) EI
AgI Ag
-$/circ$
|
|
 < EBr
AgBr Ag
-
$/circ$
|
|
 < ECl
AgCl Ag
-
$/circ$
|
|
 
 (c) EI
AgI Ag
-$/circ$
|
|
 < ECl
AgCl Ag
-
$/circ$
|
|
 < EBr
AgBr Ag
-
$/circ$
|
|
 
 (d) EI
AgI Ag
-$/circ$
|
|
 = EBr
AgBr Ag
-
$/circ$
|
|
 = ECl
AgCl Ag
-
$/circ$
|
|
EXERCISE II (JEE ADVANCED)
 
 
Visit www.jeebooks.in for more
\end{problembox}

\begin{problembox}
\textbf{Problem 135} \hfill \textbf{[Neeraj Kumar [NRJ] EXERCISE II Q10]}
\label{q:ext-nrj_all_questions-135}

The EMF of a galvanic cell composed of two 
hydrogen electrodes is 272 mV. What is the pH of 
the solution in which the anode is immersed if the 
cathode is in contact with a solution of pH = 3?
 
 (a) 3 
(b) 6.7
 
 (c) 7.6 
(d) 1.6
\end{problembox}

\begin{problembox}
\textbf{Problem 136} \hfill \textbf{[Neeraj Kumar [NRJ] EXERCISE II Q11]}
\label{q:ext-nrj_all_questions-136}

At what pH does the potential (EMF) for the 
disproportionation of chlorine change from a 
negative value to a positive value, assuming 1.0 
bar pressure and 1.0  M concentration for all 
species except hydrogen ion? Given:
 
  Cl2 + 2e- $/to$ 2Cl- ; E $/circ$ = 1.36 V
 
   2 OCl- + 4H+ + 2e- $/to$ Cl2 + 2H2O; E $/circ$ = 1.63 V.
 
  [2.303 RT/F = 0.06]
 
 (a) 4.5 
(b) 1.5
 
 (c) 2.25 
(d) 9.0
\end{problembox}

\begin{problembox}
\textbf{Problem 137} \hfill \textbf{[Neeraj Kumar [NRJ] EXERCISE II Q12]}
\label{q:ext-nrj_all_questions-137}

The EMF of the cell: Hg(l) | Hg2Cl2(s), KCl sol. 
(1.0N) | Quinohydrone | Pt, is 0.210 V at 298 K. 
What is the pH of the quinohydrone solution, 
the potential of the normal calomel electrode is 
0.279 V and E $/circ$ for the quinohydrone electrode 
is 0.699 V, both at the same temperature. [2.303 
RT/F = 0.06]
 
 (a) 3.5 
(b) 7.0
 
 (c) 1.85 
(d) -3.5
\end{problembox}

\begin{problembox}
\textbf{Problem 138} \hfill \textbf{[Neeraj Kumar [NRJ] EXERCISE II Q13]}
\label{q:ext-nrj_all_questions-138}

The potential (EMF) of a cell consisting of an anode 
of silver in 0.10 M -- AgNO3 solution and a cathode 
of Pt immersed in a solution of 1.6 M -- Cr2O7
2-, 
0.4 M -- Cr3+ and 0.1 M -- H+ is (E $/circ$Ag+|Ag = 0.80 V 
and ECr O
Cr
2
7
2
3+
-
$/circ$
|
 = 1.33 V) [2.303 RT/F = 0.06]
 
 (a) 0.46 V 
(b) 0.60 V
 
 (c) 0.53 V 
(d) -0.17 V
\end{problembox}

\begin{problembox}
\textbf{Problem 139} \hfill \textbf{[Neeraj Kumar [NRJ] EXERCISE II Q14]}
\label{q:ext-nrj_all_questions-139}

The EMF of the cell: Zn-Hg (C1M) | Zn2+(aq) | 
Zn-Hg (C2M) at 25$/circ$C, if the concentrations of 
the zinc amalgams are: C1 = 10 g per 100 g of 
mercury, C2 = 1 g per 100 g mercury, is
 
 (a) 0.059 V 
(b) 0.0295 V
 
 (c) 0.59 V 
(d) 0.295 V
\end{problembox}

\begin{problembox}
\textbf{Problem 140} \hfill \textbf{[Neeraj Kumar [NRJ] EXERCISE II Q15]}
\label{q:ext-nrj_all_questions-140}

What is the equilibrium constant of the reaction: 
2Fe3+ + Au+ $/to$ 2Fe2+ + Au3+? Given EAu
Au
+ |
$/circ$
 = 1.68 
V, EAu
Au
3+ |
$/circ$
 = 1.50 V, EFe
Fe
3+
2+
|
$/circ$
 = 0.75 V and 2.303 
RT/F = 0.06.
 
 (a) 1 $/times$ 1022 
(b) 1 $/times$ 10-22
 
 (c) 1 $/times$ 10-11 
(d) 1 $/times$ 10-72
\end{problembox}

\begin{problembox}
\textbf{Problem 141} \hfill \textbf{[Neeraj Kumar [NRJ] EXERCISE II Q16]}
\label{q:ext-nrj_all_questions-141}

What is the EMF of the cell: Pt, H2 
(1 atm) | CH3COOH (0.1 M) || (0.01 M) NH4OH 
| H2 (1 atm), Pt? Given: Ka for CH3COOH = 1.8 $/times$ 
10-5, Kb for NH4OH = 1.8 $/times$ 10-5, 2.303 RT/F = 
0.06, log 1.8 = 0.25)
 
 (a) 0.465 V 
(b) -0.465 V
 
 (c) -0.2325 V 
(d) -0.93 V
\end{problembox}

\begin{problembox}
\textbf{Problem 142} \hfill \textbf{[Neeraj Kumar [NRJ] EXERCISE II Q17]}
\label{q:ext-nrj_all_questions-142}

Two electrochemical cells are assembled in which 
the following reactions occur:
 
   V2+ + VO2+ + 2H+ $/to$ 2V3+ + H2O; E $/circ$Cell = 
0.616 V
 
   V3+ + Ag+ + H2O $/to$ VO2+ + Ag(s) + 2H+; E $/circ$Cell 
= 0.439 V
 
 If EAg
Ag
+ |
$/circ$
 = 0.799 V, what is EV
V
3+
2+
|
$/circ$
?
 
 (a) -0.256 V 
(b) +0.256 V
 
 (c) +1.854 V 
(d) -1.854 V
\end{problembox}

\begin{problembox}
\textbf{Problem 143} \hfill \textbf{[Neeraj Kumar [NRJ] EXERCISE II Q18]}
\label{q:ext-nrj_all_questions-143}

The EMF of cell: H2(g) | Bu  er || Normal calomel 
electrode, is 0.70 V at 25$/circ$C, when barometric 
pressure is 760 mm. What is the pH of the bu  er 
solution? E $/circ$Calomel = 0.28 V. [2.303 RT/F = 0.06]
 
 (a) 3.5
 
 (b) 7.0
 
 (c) tending to zero
 
 (d) tending to 14.0
\end{problembox}

\begin{problembox}
\textbf{Problem 144} \hfill \textbf{[Neeraj Kumar [NRJ] EXERCISE II Q19]}
\label{q:ext-nrj_all_questions-144}

What is the solubility product of a saturated 
solution of Ag2CrO4 in water at 298 K if the EMF 
of the cell: Ag | Ag+ (satd. Ag2CrO4) || Ag+ (0.1 
M) | Ag is 0.162 V at 298 K? [2.303 RT/F = 0.06, 
log 2 = 0.3]
 
 (a) 2.0 $/times$ 10-4 
(b) 3.2 $/times$ 10-11
 
 (c) 8.0 $/times$ 10-12 
(d) 4.0 $/times$ 10-12
\end{problembox}

\begin{problembox}
\textbf{Problem 145} \hfill \textbf{[Neeraj Kumar [NRJ] EXERCISE II Q20]}
\label{q:ext-nrj_all_questions-145}

A Tl+|Tl couple was prepared by saturating 
0.10 M -- KBr with TlBr and  allowing Tl+ ions 
form the insoluble bromide to equilibrate. This 
couple was observed to have a potential -0.444 
V with respect to Pb2+ | Pb couple in which 
Pb2+ was 0.10 M. What is the Ksp of TlBr. 
[Given: EPb
Pb
2+
$/circ$
|
  = - 0.126 V, ETl
Tl
+
$/circ$
|
 = -0.336 V, 
log 2.5 = 0.4, 2.303 RT/F = 0.06]
 
 (a) 4.0 $/times$ 10-6 
(b) 2.5 $/times$ 10-4
 
 (c) 4.0 $/times$ 10-5 
(d) 6.3 $/times$ 10-3
 
 
Visit www.jeebooks.in for more
\end{problembox}

\begin{problembox}
\textbf{Problem 146} \hfill \textbf{[Neeraj Kumar [NRJ] EXERCISE II Q21]}
\label{q:ext-nrj_all_questions-146}

The EMFs of the cell obtained by combining 
separately Zn and Cu electrodes of a Daniel cell 
with normal calomel  electrodes are 1.083 V and 
- 0.018 V, respectively, at 25  C. If the potential of 
normal calomel electrode is - 0.28 V, the EMF of 
the Daniel cell is
 
 (a) 1.065 V 
(b) 1.101 V
 
 (c) 0.803 V 
(d) 0.262 V
\end{problembox}

\begin{problembox}
\textbf{Problem 147} \hfill \textbf{[Neeraj Kumar [NRJ] EXERCISE II Q22]}
\label{q:ext-nrj_all_questions-147}

A hydrogen electrode placed in a bu  er solution 
of CH3COONa and CH3COOH in the ratio's x:y 
and y:x has electrode potential values E1 and E2 
volts,  respectively, at 25$/circ$C. The pKa value of acetic 
acid is
 
 (a) (E1 + E2)/ 0.118 
(b) (E2 - E1)/ 0.118
 
 (c) -(E1 + E2)/ 0.118 
(d) (E1 - E2)/ 0.118
\end{problembox}

\begin{problembox}
\textbf{Problem 148} \hfill \textbf{[Neeraj Kumar [NRJ] EXERCISE II Q23]}
\label{q:ext-nrj_all_questions-148}

The standard reduction potential at 25$/circ$C of the 
reaction 2H2O + 2e-   H2 + 2OH- is -0.84 V. 
Calculate equilibrium constant for the reaction: 
2H2O   H3O+ + OH- at 25$/circ$C. (2.303 RT/F = 0.06)
 
 (a) 10-14 
(b) 1014
 
 (c) 107 
(d) 10-7
\end{problembox}

\begin{problembox}
\textbf{Problem 149} \hfill \textbf{[Neeraj Kumar [NRJ] EXERCISE II Q24]}
\label{q:ext-nrj_all_questions-149}

The voltage of the cell given below is -0.61 V.
 
   Pt | H2 (1 bar) | NaHSO3 (0.4 M), Na2SO3 
(6.4 $/times$ 10-2 M) || Zn2+ (0.4 M) | Zn
 
  If E $/circ$Zn2+|Zn = - 0.76 V, Calculate Ka2 of H2SO4. 
(2.303 RT/F = 0.06)
 
 (a) 3.2 $/times$ 10-4 
(b) 3.2 $/times$ 10-2
 
 (c) 3.2 $/times$ 10-3 
(d) 6.4 $/times$ 10-7
\end{problembox}

\begin{problembox}
\textbf{Problem 150} \hfill \textbf{[Neeraj Kumar [NRJ] EXERCISE II Q25]}
\label{q:ext-nrj_all_questions-150}

Estimate the cell potential of a Daniel cell having 
1.0 M -- Zn2+ and originally having 1.0 M -- 
Cu2+ after su   cient ammonia has been added 
to the cathode compartment to make the NH3 
concentration 2.0 M. Given: E $/circ$Zn2+|Zn = - 0.76 V, 
E $/circ$Cu2+|Cu = +0.34 V, Kf for Cu (NH3)4
2+ = 1 $/times$ 1012 
[2.303 RT/F = 0.06].
 
 (a) 1.10 V 
(b) 0.704 V
 
 (c) 0.396 V 
(d) 1.496 V
\end{problembox}

\begin{problembox}
\textbf{Problem 151} \hfill \textbf{[Neeraj Kumar [NRJ] EXERCISE II Q26]}
\label{q:ext-nrj_all_questions-151}

The minimum mass of NaOH required to be 
added in RHS to consume all the H+ present in 
RHS of the cell of EMF, +0.70 V at 25$/circ$C, before 
its use.
 
   Zn | Zn2+ (0.01 M) || HCl (V = 500 ml) | H2 
(1 bar) | Pt
 
  Given: E $/circ$Zn2+|Zn = - 0.76 V; 2.303 RT/F = 0.06
 
 (a) 2.0 g 
(b) 4.0 g
 
 (c) 0.2 g 
(d) 0.4 g
\end{problembox}

\begin{problembox}
\textbf{Problem 152} \hfill \textbf{[Neeraj Kumar [NRJ] EXERCISE II Q27]}
\label{q:ext-nrj_all_questions-152}

The EMF of the cell: Ag, AgCl in 0.1 M -- KCl || 
satd. NH4NO3 || 0.1 M -- AgNO3, Ag is 0.42 V at 
25$/circ$C. 0.1 M -- KCl is 50/% dissociated and 0.1 M -- 
AgNO3 is 40/% dissociated. The solubility product 
of AgCl is (2.303 RT/F = 0.06)
 
 (a) 1.0 $/times$ 10-10 
(b) 2.0 $/times$ 10-9
 
 (c) 1.0 $/times$ 10-9 
(d) 2.0 $/times$ 10-10
\end{problembox}

\begin{problembox}
\textbf{Problem 153} \hfill \textbf{[Neeraj Kumar [NRJ] EXERCISE II Q28]}
\label{q:ext-nrj_all_questions-153}

On the basis of information available from the 
reaction
 
   4Al + 3O2 $/to$ 2Al2O3; $/Delta$G = -965  kJ/ mol of O2
 
 The minimum EMF required to carry out 
electrolysis of Al2O3 is
 
 (a) 0.833 V 
(b) 2.5 V
 
 (c) 5.0 V 
(d) 1.67 V
\end{problembox}

\begin{problembox}
\textbf{Problem 154} \hfill \textbf{[Neeraj Kumar [NRJ] EXERCISE II Q29]}
\label{q:ext-nrj_all_questions-154}

The theoretical e   ciency of a hypothetical cell is 
about 84/% which involves the following reaction
 
   A(s) + B2+(aq) $/to$ A2+(aq) + B(s); $/Delta$H = -285 kJ,
 
 then, the standard EMF of the cell is
 
 (a) 1.10 V 
(b) 1.24 V
 
 (c) 2.48 V 
(d) 2.20 V
\end{problembox}

\begin{problembox}
\textbf{Problem 155} \hfill \textbf{[Neeraj Kumar [NRJ] EXERCISE II Q30]}
\label{q:ext-nrj_all_questions-155}

To a beaker containing 0.1 M -- HCl, a little 
pure solid AgCl is added. Both a silver and a 
hydrogen electrode (PH2 = 1.0 bar) are then placed 
in the solution. What is the approximate value of 
EMF measured between the electrodes at 25$/circ$C? 
(Given: EAg
Ag
+
$/circ$
|
 = 0.80 V and Ksp (AgCl) = 10-10) 
[2.303 RT/F = 0.06]
 
 (a) 0.92 V 
(b) 0.32 V
 
 (c) 1.28 V 
(d) 0.56 V
\end{problembox}

\begin{problembox}
\textbf{Problem 156} \hfill \textbf{[Neeraj Kumar [NRJ] EXERCISE II Q31]}
\label{q:ext-nrj_all_questions-156}

Although aluminium is above hydrogen in the 
electrochemical series, it is stable in air and water. 
Why?
 
 (a) Aluminium is non-reactive metal.
 
 (b)  A layer of oxide is formed at metal surface, 
which prevent further reaction.
 
 (c)  The 
reaction 
is 
thermodynamically 
unfavourable, i.e., $/Delta$G is not negative.
 
 (d)  Kinetically, the reaction has very low 
activation energy and hence do not occur.
 
 
Visit www.jeebooks.in for more
\end{problembox}

\begin{problembox}
\textbf{Problem 157} \hfill \textbf{[Neeraj Kumar [NRJ] EXERCISE II Q32]}
\label{q:ext-nrj_all_questions-157}

When silver chloride is dissolved in a large excess 
of ammonia, practically all silver ion can be 
assumed to exist in form of a single ionic species 
[Agx(NH3)y]x+. Compute the values of x and y 
using the two following cells:
 
 Cell I:  Ag | 4.0 $/times$ 10-4 M -- AgCl, 1 M 
--  NH3  || 4 $/times$ 10-2 M -- AgCl, 1 M 
-- NH3 | Ag; Ecell = 0.118 V at 298 K.
 
 Cell II:  Ag | 3 $/times$ 10-3 M -- AgCl, 1 M 
-- NH3 || 3 $/times$ 10-3 M -- AgCl, 0.1  M 
-- NH3 | Ag; Ecell = 0.118 V at 298 K.
 
 (a) x = 1, y = 2 
(b) x = 1, y = 4
 
 (c) x = 2, y = 4 
(d) x = 1, y = 6
\end{problembox}

\begin{problembox}
\textbf{Problem 158} \hfill \textbf{[Neeraj Kumar [NRJ] EXERCISE II Q33]}
\label{q:ext-nrj_all_questions-158}

When metallic copper is shaken with a solution 
of a copper salt, the reaction Cu + Cu2+   2Cu+ 
proceeds. When equilibrium is established at 298 
K, [Cu2+]/ [Cu+]2 = 1.667 $/times$ 106 M-1. If the standard 
potential of the Cu2+ | Cu half-cell is +0.3376 V, 
what is the standard potential of Cu+| Cu half-cell? 
(Given: 2.303 RT/F = 0.06, log 2 = 0.3, log 3 = 0.48)
 
 (a) -0.3732 V 
(b) 0.6752 V
 
 (c) 0.5242 V 
(d) 0.151 V
\end{problembox}

\begin{problembox}
\textbf{Problem 159} \hfill \textbf{[Neeraj Kumar [NRJ] EXERCISE II Q34]}
\label{q:ext-nrj_all_questions-159}

Zinc granules are added in excess to a 500 ml of 
1.0 M nickel nitrate solution at 25$/circ$C until the 
equilibrium is reached. If the standard reduction 
potential of Zn2+|Zn and Ni2+|Ni are - 0.75 V and 
-0.24 V, respectively, the concentration of Ni2+ in 
solution at equilibrium is (2.303 RT/F = 0.06)
 
 (a) 1.0 $/times$ 10-17 M 
(b) 1.0 $/times$ 1017 M
 
 (c) 5 $/times$ 10-17 M 
(d) 2 $/times$ 10-17 M
\end{problembox}

\begin{problembox}
\textbf{Problem 160} \hfill \textbf{[Neeraj Kumar [NRJ] EXERCISE II Q35]}
\label{q:ext-nrj_all_questions-160}

The EMF for the cell: Ag(s) | AgCl(s) | KCl(0.2 
M) || KBr (0.001 M) | AgBr(s) | Ag(s) at 25$/circ$C is 
(Ksp(AgCl) = 2.0 $/times$ 10-10; Ksp(AgBr) = 4.0 $/times$ 10-13, 
2.303 RT/F = 0.06, log 2 = 0.3)
 
 (a) 0.024 V 
(b) -0.024 V
 
 (c) -0.24 V 
(d) - 0.012 V
\end{problembox}

\begin{problembox}
\textbf{Problem 161} \hfill \textbf{[Neeraj Kumar [NRJ] EXERCISE II Q36]}
\label{q:ext-nrj_all_questions-161}

Two students use same stock solution of ZnSO4 
but di  erent solutions of CuSO4. The EMF 
of one cell is 0.03 V higher than the other. The 
concentration of CuSO4 in the cell with higher 
EMF value is 0.5 M. The concentration of CuSO4 
in the other cell is (2.303 RT/F = 0.06)
 
 (a) 0.05 M 
(b) 5.0 M
 
 (c) 0.5 M 
(d) 0.005 M
\end{problembox}

\begin{problembox}
\textbf{Problem 162} \hfill \textbf{[Neeraj Kumar [NRJ] EXERCISE II Q37]}
\label{q:ext-nrj_all_questions-162}

Two weak acid solutions HA1 and HA2 each with 
the same concentration and having pKa values 
3 and 5 are placed in contact with hydrogen 
electrodes (1 atm, 25$/circ$C) and are interconnected 
through a salt bridge. EMF of the cell is
 
 (a) 0.0295 V 
(b) 0.118 V
 
 (c) 0.0885 V 
(d) 0.059 V
\end{problembox}

\begin{problembox}
\textbf{Problem 163} \hfill \textbf{[Neeraj Kumar [NRJ] EXERCISE II Q38]}
\label{q:ext-nrj_all_questions-163}

Consider the reaction of extraction of gold from 
its ore:
 
   Au(s) + 2CN- (aq) + 1
4 O2(g) + 1
2 H2O(l) $/to$ 
Au(CN)2
-(aq) + OH-(aq)
 
 Use the following data to calculate $/Delta$Go for the 
above reaction.
 
  Kf [Au(CN)2
-] 
= X
 
  O2 + 2H2O + 4e- $/to$ 4OH-;  E $/circ$ = +0.41 V
 
  Au3+ + 3e- $/to$ Au;                E $/circ$ = +1.50 V
 
  Au3+ + 2e- $/to$ Au+;              E $/circ$ = +1.40 V
 
 (a) -RT lnX + 1.29 F 
(b) -RT lnX - 1.29 F
 
 (c) +RT lnX + 2.11 F 
(d) -RT lnX - 2.11 F
\end{problembox}

\begin{problembox}
\textbf{Problem 164} \hfill \textbf{[Neeraj Kumar [NRJ] EXERCISE II Q39]}
\label{q:ext-nrj_all_questions-164}

F2 gas cannot be obtained by the electrolysis of 
any aqueous   uoride salt because
 
 (a) F2 is the strongest oxidizing agent
 
 (b) F2 easily combines with water
 
 (c)  F2 readily combines with the  electrodes
 
 (d) F- can never be oxidized
\end{problembox}

\begin{problembox}
\textbf{Problem 165} \hfill \textbf{[Neeraj Kumar [NRJ] EXERCISE II Q40]}
\label{q:ext-nrj_all_questions-165}

A volume of 100 ml of a bu  er of 1 M -- NH3 and 1 
M -- NH4
+ is placed in two  half-cells connected by 
a salt bridge. A current of 1.5 A is passed through 
the cell for 20 min. If electrolysis of water takes 
place only and the electrode reactions are:
 
      Right: 2H2O + O2 + 4e $/to$ 4 OH- 
and Left: 2H2O $/to$ 4H+ + O2 + 4e,
 
 then, the pH of the
 
 (a) right electrode will increase
 
 (b) left electrode will increase
 
 (c) both electrode will increase
 
 (d) both electrode will decrease
\end{problembox}

\begin{problembox}
\textbf{Problem 166} \hfill \textbf{[Neeraj Kumar [NRJ] EXERCISE II Q41]}
\label{q:ext-nrj_all_questions-166}

When molten ICI3 is electrolysed using platinum 
electrodes
 
 (a)  I2 is evolved at cathode and Cl2 at anode
 
 (b)  Cl2 is evolved at cathode and I2 at anode
 
 (c)  I2 is evolved at cathode and both I2 and Cl2 at 
anode
 
 (d) electrolysis does not take place
  
 
 
Visit www.jeebooks.in for more
\end{problembox}

\begin{problembox}
\textbf{Problem 167} \hfill \textbf{[Neeraj Kumar [NRJ] EXERCISE II Q42]}
\label{q:ext-nrj_all_questions-167}

When electric current is passed through a cell 
having an electrolyte, positive ions move towards 
the cathode and negative ions towards the anode. 
If the cathode is pulled out of the solution, 
then the
 
 (a)  positive and negative ions will move towards 
the anode
 
 (b)  positive ions will start moving towards the 
anode, the negative ions will stop moving
 
 (c)  negative ions will continue to move towards 
the anode, the positive ions will start moving 
randomly
 
 (d)  positive and negative ions will start moving 
randomly
\end{problembox}

\begin{problembox}
\textbf{Problem 168} \hfill \textbf{[Neeraj Kumar [NRJ] EXERCISE II Q43]}
\label{q:ext-nrj_all_questions-168}

When 10-6 M -- HCl is electrolysed
 
 (a) O2 is produced at the anode
 
 (b) H2 is produced at the anode
 
 (c) Cl2 is produced at the anode
 
 (d) Cl2 and O2 are produced at the anode
\end{problembox}

\begin{problembox}
\textbf{Problem 169} \hfill \textbf{[Neeraj Kumar [NRJ] EXERCISE II Q44]}
\label{q:ext-nrj_all_questions-169}

A dilute aqueous solution of Na2SO4 is electrolysed 
using platinum electrodes. The products at the 
anode and cathode are, respectively,
 
 (a) O2, H2 
(b) S2O8
2-, Na
 
 (c) O2, Na 
(d) S2O8
2-, H2
\end{problembox}

\begin{problembox}
\textbf{Problem 170} \hfill \textbf{[Neeraj Kumar [NRJ] EXERCISE II Q45]}
\label{q:ext-nrj_all_questions-170}

In the electrolysis of a fused salt, the mass of 
substance deposited on an electrode will not 
depend on
 
 (a) temperature of bath
 
 (b) current intensity
 
 (c) time of electrolysis
 
 (d) electrochemical equivalent of ions
\end{problembox}

\begin{problembox}
\textbf{Problem 171} \hfill \textbf{[Neeraj Kumar [NRJ] EXERCISE II Q46]}
\label{q:ext-nrj_all_questions-171}

When an electric current is passed through 
an aqueous solution of the following, the 
concentrations of cation as well as anion will 
not change for (assume constant volume of the 
electrolytic solution)
 
 (a) CsCl 
(b) KNO3
 
 (c) AgNO3 
(d) HCl
\end{problembox}

\begin{problembox}
\textbf{Problem 172} \hfill \textbf{[Neeraj Kumar [NRJ] EXERCISE II Q47]}
\label{q:ext-nrj_all_questions-172}

Which of the following substances: Na, Hg, 
S, Pt and graphite can be used as electrodes in 
electrolytic cells having aqueous solution?
 
 (a) Na, Pt and graphite 
(b) Na and Hg
 
 (c) Hg, Pt and graphite 
(d) Na and S
\end{problembox}

\begin{problembox}
\textbf{Problem 173} \hfill \textbf{[Neeraj Kumar [NRJ] EXERCISE II Q48]}
\label{q:ext-nrj_all_questions-173}

An ammeter and a copper voltameter are connected 
in series in an electric  circuit through which a 
constant direct  current   ows. The ammeter shows 
0.5 A. If 0.635 g of Cu is deposited in 0.965 h, what 
is the percentage error of ammeter? (Cu = 63.5)
 
 (a) 5 /% 
(b) 10 /%
 
 (c) 9 /% 
(d) 90 /%
\end{problembox}

\begin{problembox}
\textbf{Problem 174} \hfill \textbf{[Neeraj Kumar [NRJ] EXERCISE II Q49]}
\label{q:ext-nrj_all_questions-174}

The same current is passed through acidulated 
water and stannous chloride  solution. What 
volume of dry detonating gas at 0$/circ$C and 1 atm is 
evolved from water, when 1.20 g of tin is deposited 
from the other solution? (Sn = 120)
 
 (a) 112 ml 
(b) 336 ml
 
 (c) 224 ml 
(d) 672 ml
\end{problembox}

\begin{problembox}
\textbf{Problem 175} \hfill \textbf{[Neeraj Kumar [NRJ] EXERCISE II Q50]}
\label{q:ext-nrj_all_questions-175}

The preparation of LiOH by the electrolysis of a 
35/% solution of LiCl using a platinum anode led to 
a current e   ciency of 80/%. What weight of LiOH 
was formed by the passage of 2.5 A for 4825 s?
 
 (a) 1.92 g 
(b) 2.40 g
 
 (c) 0.96 g 
(d) 0.672 g
\end{problembox}

\begin{problembox}
\textbf{Problem 176} \hfill \textbf{[Neeraj Kumar [NRJ] EXERCISE II Q51]}
\label{q:ext-nrj_all_questions-176}

Assuming that copper contains only iron, silver 
and gold as impurities. After passage of 12.4 A for 
4825 s, the mass of anode decreased by 20.00 g and 
the cathode increased by 19.05 g. The percentages 
of iron and copper in the original sample are, 
respectively, (Cu = 63.5, Fe = 56)
 
 (a) 4.75/%, 95.25/% 
(b) 2.8/%, 95.25/%
 
 (c) 95.25/%, 4.75/% 
(d) 95.25/%, 1.91/%
\end{problembox}

\begin{problembox}
\textbf{Problem 177} \hfill \textbf{[Neeraj Kumar [NRJ] EXERCISE II Q52]}
\label{q:ext-nrj_all_questions-177}

After electrolysis of an aqueous sodium chloride 
solution, it was found that the solution is being 
neutralized by 60 ml N -- HCl solution. During the 
same period of electrolysis, 3.18 g of copper was 
deposited in a copper voltameter in series. What is 
the percentage of the theoretical yield of sodium 
hydroxide obtained? (Cu = 63.6)
 
 (a) 40/% 
(b) 60/%
 
 (c) 30/% 
(d) 80/%
\end{problembox}

\begin{problembox}
\textbf{Problem 178} \hfill \textbf{[Neeraj Kumar [NRJ] EXERCISE II Q53]}
\label{q:ext-nrj_all_questions-178}

A lead storage battery has initially 200  g of 
lead and 200 g of PbO2 plus excess H2SO4. 
Theoretically, how long could this cell deliver a 
current of 10.0 A, without recharging, if it were 
possible to operate it so that the reaction goes to 
completion? (Pb = 208)
 
 (a) 16083.33 s 
(b) 8041.67 s
 
 (c) 44.68 s 
(d) 18557.7 s
 
 
Visit www.jeebooks.in for more
\end{problembox}

\begin{problembox}
\textbf{Problem 179} \hfill \textbf{[Neeraj Kumar [NRJ] EXERCISE II Q54]}
\label{q:ext-nrj_all_questions-179}

A volume of 100 ml of 0.6 N-CuSO4 solution is 
electrolysed between two platinum electrodes till the 
concentration in the residual liquid is 0.1 N, when 
a steady current of 5.0 A is used. How long should 
the current be passed to get the above change?
 
 (a) 965 s 
(b) 96500 s
 
 (c) 1930 s 
(d) 482.5 s
\end{problembox}

\begin{problembox}
\textbf{Problem 180} \hfill \textbf{[Neeraj Kumar [NRJ] EXERCISE II Q55]}
\label{q:ext-nrj_all_questions-180}

A constant current   owed for 2 h through a 
potassium iodide solution, oxidizing the iodide 
ion to iodine. At the end of the experiment, the 
iodine was titrated with 72 ml of 1.0 M-Na2S2O3 
solution. What was the average rate of current 
  ow, in ampere?
 
 (a) 0.965 A 
(b) 1.93 A
 
 (c) 0.483 A 
(d) 0.0965 A
\end{problembox}

\begin{problembox}
\textbf{Problem 181} \hfill \textbf{[Neeraj Kumar [NRJ] EXERCISE II Q56]}
\label{q:ext-nrj_all_questions-181}

Anthracene, 
C14H10, 
can 
be 
oxidized 
to 
anthraquinone, 
C14H8O2. 
What 
weight 
of 
anthraquinone can be produced by the passage of 
a current of 1 A of 40 min if the current e   ciency 
is 96.5/%?
 
 (a) 0.862 g 
(b) 0.832 g
 
 (c) 0.624 g 
(d) 0.738 g
\end{problembox}

\begin{problembox}
\textbf{Problem 182} \hfill \textbf{[Neeraj Kumar [NRJ] EXERCISE II Q57]}
\label{q:ext-nrj_all_questions-182}

Most of the copper used to make wire has been 
electrically re  ned by depositing it from copper 
salts solution (divalent) on to a cathode. What 
is the cost of electrical energy required per kg of 
copper if the cost of electricity is Rs 4.00 per kWh 
and the cell operates at 0.33 V? The electrochemical 
equivalent of copper is 0.00033 g/coulomb.
 
 (a) Rs 11.11 
(b) Rs 5.55
 
 (c) Rs 2.22 
(d) Rs 1.11
\end{problembox}

\begin{problembox}
\textbf{Problem 183} \hfill \textbf{[Neeraj Kumar [NRJ] EXERCISE II Q58]}
\label{q:ext-nrj_all_questions-183}

Electrolysis of an acetate solution produces ethane 
according to the reaction:
 
  2CH3COO- $/to$ C2H6(g) + 2CO2(g) + 2e-
 
 What total volume of ethane and CO2 would be 
produced at 0$/circ$C and 1 atm, if a current of 0.5 
A is passed through the solution for 482.5 min? 
Assume current e   ciency 80/%.
 
 (a) 1.344 L 
(b) 2.688 L
 
 (c) 4.032 L 
(d) 1.792 L
\end{problembox}

\begin{problembox}
\textbf{Problem 184} \hfill \textbf{[Neeraj Kumar [NRJ] EXERCISE II Q59]}
\label{q:ext-nrj_all_questions-184}

To perform an analysis of a mixture of metal ions by 
electro-deposition, the second metal to be deposited 
must not being plating out until the concentration 
ratio of the second to the   rst is about 106. What 
must be the minimum di  erence in standard 
potential of the two metals which form dipositive 
ions in order for such an analysis to be feasible?
 
 (a) 0.177 V 
(b) 0.354 V
 
 (c) 0.708 V 
(d) 0.088 V
\end{problembox}

\begin{problembox}
\textbf{Problem 185} \hfill \textbf{[Neeraj Kumar [NRJ] EXERCISE II Q60]}
\label{q:ext-nrj_all_questions-185}

During the electrolysis of 0.1 M -- CuSO4 solution 
using copper electrodes, a depletion of Cu2+ 
occurs near the cathode with a corresponding 
excess near the anode, owing to ine   cient stirring 
of the solution. If the local concentration of Cu2+ 
near the anode and cathode are, respectively, 0.12 
M and 0.08 M, the back EMF developed at 298 K 
is (log 1.5 = 0.18, 2.303 RT/F = 0.06)
 
 (a) 0.33 V 
(b) 5.4 mV
 
 (c) 2.7 mV 
(d) 10.8 mV
\end{problembox}

\begin{problembox}
\textbf{Problem 186} \hfill \textbf{[Neeraj Kumar [NRJ] EXERCISE II Q61]}
\label{q:ext-nrj_all_questions-186}

The following galvanic cell:
 
   Zn | Zn(NO3)2 (100 ml, 1 M) || Cu(NO3)2 (100 
ml, 1 M) | Cu
 
 was operated as an electrolytic cell as Cu as 
the anode and Zn as the cathode. A current 
of 0.4825 A was passed for 10 h and then the 
cell was allowed to function as galvanic cell. 
What would be the   nal EMF of the cell at 
25$/circ$C? Assume that the only electrode reactions 
occurring were those involving Cu | Cu2+ and 
Zn | Zn2+. Given: E $/circ$Cu2+|Cu = +0.34 V, E $/circ$Zn2+|Zn 
= -0.76 V, 2.303 RT/F = 0.06, log 1.9 = 0.28)
 
 (a) 1.10 V 
(b) 1.0616 V
 
 (c) 1.1084 V 
(d) 1.1768 V
\end{problembox}

\begin{problembox}
\textbf{Problem 187} \hfill \textbf{[Neeraj Kumar [NRJ] EXERCISE II Q62]}
\label{q:ext-nrj_all_questions-187}

At the Nangal fertilizer plant in Punjab, hydrogen 
is produced by the electrolysis of water. The 
hydrogen is used for the production of ammonia 
and nitric acid (by the oxidation of ammonia). If 
the average production of ammonium nitrate is 
1152 kg/day, the daily consumption of electricity 
(in A /day) is
 
 (a) 96,500 
(b) 48,250
 
 (c) 32,166.67 
(d) 24,125
\end{problembox}

\begin{problembox}
\textbf{Problem 188} \hfill \textbf{[Neeraj Kumar [NRJ] EXERCISE II Q63]}
\label{q:ext-nrj_all_questions-188}

Lactic acid, HC3H5O3, produced in 1 g sample of 
muscle tissue was titrated using phenolphthalein 
as indicator against OH- ions which were obtained 
by the electrolysis of water. As soon as OH- ions 
are produced, they react with lactic acid and at 
complete neutralization, immediately a pink colour 
is noticed. If electrolysis was made for 1158 s using 
50.0 mA current to reach the end point, what was 
the percentage of lactic acid in muscle  tissue?
 
 (a) 5.4/% 
(b) 2.7/%
 
 (c) 10.8/% 
(d) 0.054/%
 
 
Visit www.jeebooks.in for more
\end{problembox}

\begin{problembox}
\textbf{Problem 189} \hfill \textbf{[Neeraj Kumar [NRJ] EXERCISE II Q64]}
\label{q:ext-nrj_all_questions-189}

H2O2 can be prepared by successive reactions:
 
  2NH4HSO4 $/to$ H2 + (NH4)2S2O8
 
   (NH4)2S2O8 + 2H2O $/to$ 2NH4HSO4 + H2O2
 
 The   rst reaction is an electrolytic reaction and 
second is steam distillation. What amount of 
current would have to be used in   rst reaction to 
produce enough intermediate to yield 102 g pure 
H2O2 per hour. Assume current e   ciency 50/%.
 
 (a) 643.33 A 
(b) 321.67 A
 
 (c) 160.83 A 
(d) 1286.67 A
\end{problembox}

\begin{problembox}
\textbf{Problem 190} \hfill \textbf{[Neeraj Kumar [NRJ] EXERCISE II Q65]}
\label{q:ext-nrj_all_questions-190}

In an analytical determination of arsenic, a 
solution containing arsenious acid, H3AsO3, KI 
and a small amount of starch is electrolysed. 
The electrolysis produces free I2 from I- ions and 
the I2 immediately oxidizes the arsenious acid to 
hydrogen arsenate ion, HAsO4
2-.
 
   I2 + H3AsO3 + H2O $/to$ 2I- + HAsO4
2- + 4H+
 
 When the oxidation of arsenic is complete, the 
free iodine combines with the starch to give a deep 
blue colour. If during a particular run, it takes 
96.5 s for a current of 1.68 mA to give an end 
point (indicated by the blue colour), how many g 
of arsenic are present in the solution. (As = 75)
 
 (a) 6.3 $/times$ 10-4 g 
(b) 1.26 $/times$ 10-4 g
 
 (c) 3.15 $/times$ 10-5 g 
(d) 6.3 $/times$ 10-5 g
\end{problembox}

\begin{problembox}
\textbf{Problem 191} \hfill \textbf{[Neeraj Kumar [NRJ] EXERCISE II Q66]}
\label{q:ext-nrj_all_questions-191}

Electrolysis of a solution of MnSO4 in aqueous 
sulphuric acid is a method for the preparation 
of MnO2 as per the reaction: Mn2+(aq) + 2H2O 
$/to$ MnO2(s) + 2H+(aq) + H2(g). Passing a current 
of 19.3 A for 2 h gives only 52.2 g of MnO2. The 
current e   ciency is (Mn = 55)
 
 (a) 8.33/% 
(b) 83.33/%
 
 (c) 41.67/% 
(d) 100/%
\end{problembox}

\begin{problembox}
\textbf{Problem 192} \hfill \textbf{[Neeraj Kumar [NRJ] EXERCISE II Q67]}
\label{q:ext-nrj_all_questions-192}

Copper sulphate solution (250 ml) was electrolysed 
using a platinum anode and a copper cathode. A 
constant current of 2 mA was passed for 19.3 min. 
It was found that after electrolysis the absorbance 
of the solution was reduced to 50/% of its original 
value. Calculate the concentration of copper 
sulphate in the solution to begin with.
 
 (a) 9.6 $/times$ 10-5 M 
(b) 4.8 $/times$ 10-5 M
 
 (c) 2.4 $/times$ 10-5 M 
(d) 1.2 $/times$ 10-5 M
\end{problembox}

\begin{problembox}
\textbf{Problem 193} \hfill \textbf{[Neeraj Kumar [NRJ] EXERCISE II Q68]}
\label{q:ext-nrj_all_questions-193}

Calculate the quantity of electricity that would 
be required to reduce 12.3 g of nitrobenzene to 
aniline if the current e   ciency for the process is 
50/%. If the potential drop across the cell is 3.0 V, 
how much energy will be consumed?
 
 (a) 0.3 F, 86.85 kJ 
(b) 0.6 F, 173.7 kJ
 
 (c) 1.2 F, 347.4 kJ 
(d) 0.6 F, 173.7 kJ
\end{problembox}

\begin{problembox}
\textbf{Problem 194} \hfill \textbf{[Neeraj Kumar [NRJ] EXERCISE II Q69]}
\label{q:ext-nrj_all_questions-194}

Perdisulphuric acid, H2S2O8 can be prepared 
by electrolytic oxidation of H2SO4, oxygen 
and hydrogen gases are by products. In such an 
electrolysis, 9.08 L of H2 and 2.27 L of O2 were 
generated at STP. What is the mass of H2S2O8 
formed?
 
 (a) 0 
(b) 77.6 g
 
 (c) 38.8 g 
(d) 19.4 g
\end{problembox}

\begin{problembox}
\textbf{Problem 195} \hfill \textbf{[Neeraj Kumar [NRJ] EXERCISE II Q70]}
\label{q:ext-nrj_all_questions-195}

During the discharge of a lead storage battery, the 
density of sulphuric acid fell from 1.5 to 1.1 g/ml. 
Sulphuric acid of density 1.5 g/ml is 40/% H2SO4, 
by weight, and that of density 1.1 g/ml is 10/% 
H2SO4, by weight. The battery holds 3.6 L of the 
acid and the volume remained practically constant 
during the discharge. Calculate the number of 
ampere-hours which the battery should have been 
used. The electrode reactions are:
 
  Pb + SO4
2- $/to$ PbSO4 + 2e
 
   PbO2 + 4H+ + SO4
2- + 2e $/to$ PbSO4 + 2H2O
 
 (a) 965 
(b) 482.5
 
 (c) 1930 
(d) 241.25
\end{problembox}

\begin{problembox}
\textbf{Problem 196} \hfill \textbf{[Neeraj Kumar [NRJ] EXERCISE II Q71]}
\label{q:ext-nrj_all_questions-196}

A cell whose resistance, when   lled with 0.1 M 
-- KCl is 200 $/Omega$, is measured to be 6400 $/Omega$, when 
  lled with 0.003 M -- NaCl solution. What is the 
molar conductance of NaCl solution, in $/Omega$-1 cm2 
mol-1 if the molar conductance of 0.1 M -- KCl is 
120 $/Omega$-1 cm2 mol-1?
 
 (a) 41.67 
(b) 250
 
 (c) 125 
(d) 375
\end{problembox}

\begin{problembox}
\textbf{Problem 197} \hfill \textbf{[Neeraj Kumar [NRJ] EXERCISE II Q72]}
\label{q:ext-nrj_all_questions-197}

At 18$/circ$C, the ionic mobilities of NH4
+ and CrO4
2- 
ions are 6.6 $/times$ 10-8 and 5.4 $/times$ 10-8 m2 volt-1 s-1 at 
in  nite dilution. What is the molar conductance 
of ammonium chromate solution in $/Omega$-1 m2 mol-1?
 
 (a) 0.01795 
(b) 0.01158
 
 (c) 1.2 $/times$ 10-8 
(d) 1.86 $/times$ 10-7
\end{problembox}

\begin{problembox}
\textbf{Problem 198} \hfill \textbf{[Neeraj Kumar [NRJ] EXERCISE II Q73]}
\label{q:ext-nrj_all_questions-198}

At 25$/circ$C, the molar conductance at in  nite dilution 
for HCl solution is 4.25 $/Omega$-1 m2 mol-1, while its 
speci  c conductance is 382.5 $/Omega$-1 m-1. If the degree 
of dissociation is 90/%, the molarity of solution is
 
 (a) 0.9 M 
(b) 1.0 M
 
 (c) 0.1 M 
(d) 1.1 M
 
 
Visit www.jeebooks.in for more
\end{problembox}

\begin{problembox}
\textbf{Problem 199} \hfill \textbf{[Neeraj Kumar [NRJ] EXERCISE II Q74]}
\label{q:ext-nrj_all_questions-199}

A big irregular shaped vessel contained water, 
conductivity of which was 2.56 $/times$ 10-3 S-1m-1. 
585 g of NaCl was then added to the water and 
conductivity after the addition of NaCl, was found 
to be 3.06 $/times$ 10-3 S-1m-1. The molar conductivity 
of NaCl at this concentration is 1.5 $/times$ 10-2 S-1m2 
mol-1. The capacity of vessel if it is ful  lled with 
water, is
 
 (a) 3 $/times$ 104 L 
(b) 30 L
 
 (c) 3 $/times$ 108 L 
(d) 3 $/times$ 105 L
\end{problembox}

\begin{problembox}
\textbf{Problem 200} \hfill \textbf{[Neeraj Kumar [NRJ] EXERCISE II Q75]}
\label{q:ext-nrj_all_questions-200}

The molar conductivity of 0.10 M solution of 
MgCl2 is 100 mho cm2 mol-1, at 25$/circ$C. A cell with 
electrodes that are 1.50 cm2 in surface area and 0.50 
cm apart is   lled with 0.10 M -- MgCl2 solution. 
How much current will   ow when the potential 
di  erence between the electrodes is 5 volts?
 
 (a) 0.03 A 
(b) 3.0 A
 
 (c) 0.15 A 
(d) 15 A
\end{problembox}

\begin{problembox}
\textbf{Problem 201} \hfill \textbf{[Neeraj Kumar [NRJ] EXERCISE II Q76]}
\label{q:ext-nrj_all_questions-201}

For Na+, the value of symbol $/lambda$m$/circ$ is 50.0 $/Omega$-1 cm2 
mol-1. The speed of Na+ ion in the solution, if in 
the cell, electrodes are 5 cm apart and to which a 
potential of 19.3 volt is applied is
 
 (a) 2 $/times$ 10-3 cm/s 
(b) 1 $/times$ 10-3 cm/s
 
 (c) 2 $/times$ 10-4 cm/s 
(d) 2 $/times$ 10-2 cm/s
\end{problembox}

\begin{problembox}
\textbf{Problem 202} \hfill \textbf{[Neeraj Kumar [NRJ] EXERCISE II Q77]}
\label{q:ext-nrj_all_questions-202}

The conductivity of a saturated solution of AgCl 
at 298 K was found to be 3.40 $/times$ 10-6 $/Omega$-1 cm-1; the 
conductivity of water used to make up the solution 
was 1.60 $/times$ 10-6 $/Omega$-1 cm-1. Determine the solubility 
of AgCl in water in mole per litre at 298 K. The 
equivalent conductivity of AgCl at in  nite dilution 
is 150.0 $/Omega$-1 cm-2 eq-1.
 
 (a) 1.44 $/times$ 10-10 
(b) 1.2 $/times$ 10-5
 
 (c) 3.33 $/times$ 10-5 
(d) 1.2 $/times$ 10-8
\end{problembox}

\begin{problembox}
\textbf{Problem 203} \hfill \textbf{[Neeraj Kumar [NRJ] EXERCISE II Q78]}
\label{q:ext-nrj_all_questions-203}

Resistance of 0.2 M solution of an electrolyte is 50 
$/Omega$. The speci  c conductance of the solution is 1.4 S 
m-1. The resistance of 0.5 M solution of the same 
electrolyte is 280 $/Omega$. The molar conductivity of 
0.5 M solution of the electrolyte (in S m2 mol-1) is
 
 (a) 5 $/times$ 10-3 
(b) 5 $/times$ 103
 
 (c) 5 $/times$ 102 
(d) 5 $/times$ 10-4
\end{problembox}

\begin{problembox}
\textbf{Problem 204} \hfill \textbf{[Neeraj Kumar [NRJ] EXERCISE II Q79]}
\label{q:ext-nrj_all_questions-204}

The 
conductivity 
of 
a 
saturated 
solution 
containing AgA (Ksp = 3 $/times$ 10-14) and AgB (Ksp 
= 1 $/times$ 10-14) is 3.75 $/times$ 10-8 $/Omega$-1cm-1. If the limiting 
molar conductivity of Ag+ and A- ion is 60 and 
80 $/Omega$-1
 cm2 mol-1, respectively, the limiting molar 
conductivity of B- (in $/Omega$-1 cm2 mol-1) is
 
 (a) 135 
(b) 67.5
 
 (c) 270 
(d) 190
\end{problembox}

\begin{problembox}
\textbf{Problem 205} \hfill \textbf{[Neeraj Kumar [NRJ] EXERCISE II Q80]}
\label{q:ext-nrj_all_questions-205}

The conductivity of saturated solution of 
Ba3(PO4)2 is 1.2 $/times$ 10-5 $/Omega$-1 cm-1. The limiting 
equivalent conductivities of BaCl2, K3PO4 and 
KCl are 160, 140 and 100 $/Omega$-1 cm2 eq-1, respectively. 
The solubility product of Ba3(PO4)2 is
 
 (a) 10-5 
(b) 1.08 $/times$ 10-23
 
 (c) 1.08 $/times$ 10-25 
(d) 1.08 $/times$ 10-27
Section B (One or More than one Correct)
\end{problembox}

\begin{problembox}
\textbf{Problem 206} \hfill \textbf{[Neeraj Kumar [NRJ] EXERCISE II Q1]}
\label{q:ext-nrj_all_questions-206}

Which of the following statement(s) di  erentiate 
between electrochemical cell and electrolytic cell?
 
 (a)  Spontaneous or non-spontaneous nature of 
the chemical process.
 
 (b)  Chemical reactions occurring at the electrodes.
 
 (c)  Positive and negative nature of anode.
 
 (d) Dependence on Faraday's law.
\end{problembox}

\begin{problembox}
\textbf{Problem 207} \hfill \textbf{[Neeraj Kumar [NRJ] EXERCISE II Q2]}
\label{q:ext-nrj_all_questions-207}

Pick up the false statement(s):
 
 (a)  The net chemical change in a galvanic cell 
reaction is always redox reactions.
 
 (b)  In a galvanic cell made of copper and 
cadmium electrodes, cadmium electrode may 
act as anode.
 
 (c)  Standard potential increases with increasing 
concentration of the electrolyte.
 
 (d)  Calomel electrode is a reference electrode 
having 0.00 volt potential.
\end{problembox}

\begin{problembox}
\textbf{Problem 208} \hfill \textbf{[Neeraj Kumar [NRJ] EXERCISE II Q3]}
\label{q:ext-nrj_all_questions-208}

Consider the cell: Ag(s), AgCl(s) | KCl (0.1 M) | 
Hg2Cl2(s), Hg(l). The cell  potential
 
 (a)  increases on increasing concentration of 
Cl- ions.
 
 (b)  decreases on decreasing concentration of 
Cl- ions.
 
 (c)  is independent of concentration of Cl- ions.
 
 (d)  is independent of amounts of AgCl and 
Hg2Cl2.
\end{problembox}

\begin{problembox}
\textbf{Problem 209} \hfill \textbf{[Neeraj Kumar [NRJ] EXERCISE II Q4]}
\label{q:ext-nrj_all_questions-209}

The passage of electricity in the Daniel cell when 
Zn and Cu electrodes are connected
 
 (a) from Cu to Zn inside the cell
 
 (b) from Cu to Zn outside the cell
 
 (c) from Zn to Cu inside the cell
 
 (d) from Zn to Cu outside the cell
 
 
Visit www.jeebooks.in for more
\end{problembox}

\begin{problembox}
\textbf{Problem 210} \hfill \textbf{[Neeraj Kumar [NRJ] EXERCISE II Q5]}
\label{q:ext-nrj_all_questions-210}

From an electrolyte, one mole of electron will 
deposit at cathode
 
 (a) 63.5 g of Cu 
(b) 24 g of Mg
 
 (c) 11.5 g of Na 
(d) 9.0 g of Al
\end{problembox}

\begin{problembox}
\textbf{Problem 211} \hfill \textbf{[Neeraj Kumar [NRJ] EXERCISE II Q6]}
\label{q:ext-nrj_all_questions-211}

Which of the following statements is/are correct?
 
 (a)  The conductivity of molten NaCl is due to 
movement of Na+ and Cl- ions.
 
 (b)  Solid NaCl is good conductor of electricity.
 
 (c)  Molten sodium is a good conductor because 
of mobile electrons.
 
 (d)  Resistivity is reciprocal of molar conductivity 
of electrolyte.
\end{problembox}

\begin{problembox}
\textbf{Problem 212} \hfill \textbf{[Neeraj Kumar [NRJ] EXERCISE II Q7]}
\label{q:ext-nrj_all_questions-212}

Ionic conductance at in  nite dilution of A13+ and 
SO4
2- ions are 60 and 80  $/Omega$-1 cm2 eq-1, respectively. 
The correct detail(s) regarding Al2(SO4)3 is/are
 
 (a)  the molar conductance is 140 $/Omega$-1 cm2 mol-1.
 
 (b)  the equivalent conductance is 140$/Omega$-1 cm2 eq-1.
 
 (c)  the molar conductance is 840 $/Omega$-1 cm2 mol-1.
 
 (d)  the molar conductance is 23.33$/Omega$-1 cm2 mol-1.
\end{problembox}

\begin{problembox}
\textbf{Problem 213} \hfill \textbf{[Neeraj Kumar [NRJ] EXERCISE II Q8]}
\label{q:ext-nrj_all_questions-213}

The function(s) of salt bridge in a cell is/ are
 
 (a)  It maintains standard electrode potential of 
cell constant which depends on several factors.
 
 (b)  It completes the electrical circuit.
 
 (c)  It departs both the solutions from each other.
 
 (d)  It maintains the electrical neutrality of both 
electrolytic solutions.
\end{problembox}

\begin{problembox}
\textbf{Problem 214} \hfill \textbf{[Neeraj Kumar [NRJ] EXERCISE II Q9]}
\label{q:ext-nrj_all_questions-214}

The reactions taking place in the dry cell are:
 
  Anode: Zn $/to$ Zn2+ + 2e-
 
  Cathode:  2MnO2 + 2NH4
+ + 2e- $/to$ Mn2O3 + 
2NH3 + H2O
 
 The minimum mass of reactants, if a dry cell is to 
generate 0.25A for 9.65 h, are (Mn = 55, Zn = 65.4) 
(neglect any other chemical reactions occurring in 
the cell)
 
 (a) 2.943 g Zn 
(b) 7.83 g MnO2
 
 (c) 1.62 g NH4
+ 
(d) 3.915 g MnO2
\end{problembox}

\begin{problembox}
\textbf{Problem 215} \hfill \textbf{[Neeraj Kumar [NRJ] EXERCISE II Q10]}
\label{q:ext-nrj_all_questions-215}

EMF of the cell: Cd(s) | CdCl2.5H2O (sat.) | AgCl(s) 
| Ag(s) is +0.70 V at 0$/circ$C and +0.60 V at 50$/circ$C. If 
$/Delta$Ho and $/Delta$So are temperature independent, then 
the correct information(s) regarding the cell 
reaction is/are
 
 (a) $/Delta$G $/circ$ = -115.8 kJ at 50$/circ$C
 
 (b) $/Delta$G $/circ$ = 135.1 kJ at 0$/circ$C
 
 (c) $/Delta$S $/circ$ = -386 J/K
 
 (d) $/Delta$H $/circ$ = -221.178 kJ
\end{problembox}

\begin{problembox}
\textbf{Problem 216} \hfill \textbf{[Neeraj Kumar [NRJ] EXERCISE II Q11]}
\label{q:ext-nrj_all_questions-216}

Two litre solution of a bu  er mixture containing 
1.0 M -- NaH2PO4 and 1.0 M -- Na2HPO4 is 
placed in two compartments (one litre in each) of 
an electrolytic cell. The platinum electrodes are 
inserted in each compartment and 1.25 A current 
is passed for 965 min. Assuming electrolysis of 
only water at each compartment, what will be 
pH in each compartment after passage of above 
charge? (pKa for H2PO4
- = 2.15, log 7 = 0.85)
 
 (a) Anode: 3.00 
(b) Cathode: 3.00
 
 (c) Anode: 1.30 
(d) Cathode: 1.30
\end{problembox}

\begin{problembox}
\textbf{Problem 217} \hfill \textbf{[Neeraj Kumar [NRJ] EXERCISE II Q12]}
\label{q:ext-nrj_all_questions-217}

For the reduction of NO3
- ion in an aqueous 
solution, E $/circ$ is +0.96 V. Values of E $/circ$ for some 
metal ions are given below:
 
  V2+(aq) + 2e- $/to$ V(s) 
E $/circ$ = -1.19 V
 
  Fe3+(aq) + 3e- $/to$ Fe(s) 
E $/circ$ = -0.04 V
 
  Au3+(aq) + 3e- $/to$ Au(s) E $/circ$ = +1.40 V
 
  Hg2+(aq) + 2e- $/to$ Hg(l) E $/circ$ = +0.86 V
 
 The pair(s) of metal that is(are) oxidized by NO3
- 
in aqueous solution is(are)
 
 (a) V and Hg 
(b) Hg and Fe
 
 (c) Fe and Au 
(d) Fe and V
\end{problembox}

\begin{problembox}
\textbf{Problem 218} \hfill \textbf{[Neeraj Kumar [NRJ] EXERCISE II Q13]}
\label{q:ext-nrj_all_questions-218}

Among the following, the intensive property is 
(properties are)
 
 (a) Molar conductivity 
(b) Electromotive force
 
 (c) Resistance 
(d) Heat capacity
\end{problembox}

\begin{problembox}
\textbf{Problem 219} \hfill \textbf{[Neeraj Kumar [NRJ] EXERCISE II Q14]}
\label{q:ext-nrj_all_questions-219}

In a galvanic cell, the salt bridge
 
 (a)  does not participate chemically in the cell 
reaction.
 
 (b)  stops the di  usion of ions from one electrode 
to another.
 
 (c)  is necessary for the occurrence of the cell 
reaction.
 
 (d)  ensures mixing of the two electrolytic 
solutions.
\end{problembox}

\begin{problembox}
\textbf{Problem 220} \hfill \textbf{[Neeraj Kumar [NRJ] EXERCISE II Q15]}
\label{q:ext-nrj_all_questions-220}

A lead storage cell is discharged which causes the 
H2SO4 electrolyte to change from a concentration 
of 40/% by weight (density = 1.260 g/ml) to 28/%, 
by weight. The original volume of solution was 
1 L. Identify the correct statement(s):
 
 (a)  The overall cell reaction is: Pb(s) + PbO2(s) + 
2H2SO4(aq) $/to$ 2PbSO4(s) + 2H2O(l).
 
 (b)  A total of 2.0 moles of H2SO4 is reacted.
 
 (c)  The total charge released from anode of the 
cell is 1.93 $/times$ 105 coulomb.
 
 (d)  The mass of electrolytic solution has 
decreased.
 
 
Visit www.jeebooks.in for more
\end{problembox}

\begin{problembox}
\textbf{Problem 221} \hfill \textbf{[Neeraj Kumar [NRJ] EXERCISE II Q1]}
\label{q:ext-nrj_all_questions-221}

EZn
Zn
2+
$/circ$
|
 
 (a) +0.24 V 
(b) -0.76 V
 
 (c) +1.76 V 
(d) -0.34 V
\end{problembox}

\begin{problembox}
\textbf{Problem 222} \hfill \textbf{[Neeraj Kumar [NRJ] EXERCISE II Q2]}
\label{q:ext-nrj_all_questions-222}

ECu
Cu
2+
$/circ$
|
 
 (a) -0.76 V 
(b) +1.34 V
 
 (c) +1.76 V 
(d) +0.34 V
\end{problembox}

\begin{problembox}
\textbf{Problem 223} \hfill \textbf{[Neeraj Kumar [NRJ] EXERCISE II Q3]}
\label{q:ext-nrj_all_questions-223}

Zn-Cu Cell
 
 (a) +1.10 V
 
 (b) -1.10 V
 
 (c) 0.0 V
 
 (d) Indeterminate
Comprehension II
The standard reduction potential of the Ag+|Ag electrode at 298 K is 0.80 V. The solubility product of AgI is 6.4 $/times$ 10-17 
at 298 K. (2.303 RT/F = 0.06, log 2 = 0.3)
\end{problembox}

\begin{problembox}
\textbf{Problem 224} \hfill \textbf{[Neeraj Kumar [NRJ] EXERCISE II Q4]}
\label{q:ext-nrj_all_questions-224}

The potential of Ag+ | Ag electrode in a saturated 
solution of AgI at 298 K is
 
 (a) -0.314 V 
(b) +0.314 V
 
 (c) -0.172 V 
(d) +0.172 V
\end{problembox}

\begin{problembox}
\textbf{Problem 225} \hfill \textbf{[Neeraj Kumar [NRJ] EXERCISE II Q5]}
\label{q:ext-nrj_all_questions-225}

The standard reduction potential of I-|AgI|Ag 
electrode at 298 K is
 
 (a) -0.314 V 
(b) +0.314 V
 
 (c) -0.172 V 
(d) +0.172 V
\end{problembox}

\begin{problembox}
\textbf{Problem 226} \hfill \textbf{[Neeraj Kumar [NRJ] EXERCISE II Q6]}
\label{q:ext-nrj_all_questions-226}

The potential of I- (0.04 M)|AgI|Ag electrode at 
198 K is
 
 (a) -0.088 V
 
 (b) +0.088 V
 
 (c) -0.172 V
 
 (d) +0.172 V
Comprehension III
The Edison storage cell is represented as Fe(s) | FeO(s) | KOH(aq) | Ni2O3(s) | NiO(s) | Ni(s). The half-cell reactions are
Ni2O3(s) + H2O(l) + 2e   2NiO(s) + 2OH-; E $/circ$ = +0.40 V
FeO(s) + H2O(l) + 2e   Fe(s) + 2OH-; E  = - 0.87 V
\end{problembox}

\begin{problembox}
\textbf{Problem 227} \hfill \textbf{[Neeraj Kumar [NRJ] EXERCISE II Q7]}
\label{q:ext-nrj_all_questions-227}

What is the cell reaction?
 
 (a) Ni2O3(s) + Fe(s) $/to$ 2 NiO(s) + FeO(s)
 
 (b) 2NiO(s) + FeO(s) $/to$ Ni2O3(s) + Fe(s)
 
 (c)  Ni2O3(s) + FeO(s) $/to$ 2NiO(s) + Fe(s) + O2
 
 (d) None of these
\end{problembox}

\begin{problembox}
\textbf{Problem 228} \hfill \textbf{[Neeraj Kumar [NRJ] EXERCISE II Q8]}
\label{q:ext-nrj_all_questions-228}

What is the standard cell EMF?
 
 (a) 1.27 V
 
 (b) 0.47 V
 
 (c) -1.27 V
 
 (d) -0.47 V
\end{problembox}

\begin{problembox}
\textbf{Problem 229} \hfill \textbf{[Neeraj Kumar [NRJ] EXERCISE II Q9]}
\label{q:ext-nrj_all_questions-229}

How does cell EMF depend on increasing the 
concentration of KOH?
 
 (a) increases
 
 (b) decreases
 
 (c) remains una  ected
 
 (d) none of these
\end{problembox}

\begin{problembox}
\textbf{Problem 230} \hfill \textbf{[Neeraj Kumar [NRJ] EXERCISE II Q10]}
\label{q:ext-nrj_all_questions-230}

What is the maximum amount of electrical energy 
that can be obtained from one mole of Ni2O3?
 
 (a) 2.54 J 
(b) 245.11 kJ
 
 (c) 122.56 kJ 
(d) 61.28 kJ
 
 
Visit www.jeebooks.in for more
\end{problembox}

\begin{problembox}
\textbf{Problem 231} \hfill \textbf{[Neeraj Kumar [NRJ] EXERCISE II Q11]}
\label{q:ext-nrj_all_questions-231}

In which dish, the solution is acidic?
 
 (a) Dish A 
(b) Dish B
 
 (c) Both 
(d) None
\end{problembox}

\begin{problembox}
\textbf{Problem 232} \hfill \textbf{[Neeraj Kumar [NRJ] EXERCISE II Q12]}
\label{q:ext-nrj_all_questions-232}

How many moles of electrons pass through the 
circuit when 0.60 mole of Hg2+ and 0.30 mole 
of HNO2 are produced in the cell that contains 
0.50 mole of Hg2
2+ and 0.40 mole of NO3
- at the 
beginning of the reaction?
 
 (a) 0.30 
(b) 0.60
 
 (c) 0.15 
(d) 1.20
\end{problembox}

\begin{problembox}
\textbf{Problem 233} \hfill \textbf{[Neeraj Kumar [NRJ] EXERCISE II Q13]}
\label{q:ext-nrj_all_questions-233}

How long will it take to produce 0.10  mole of 
HNO2 by this reaction if a current of 10 A passes 
through the cell?
 
 (a) 965 s 
(b) 193 s
 
 (c) 1930 s 
(d) 482.5 s
Comprehension V
A fuel cell is the device to convert the energy of a fuel into electrical energy without the use of heat  engine, where the 
fuel is burnt directly. Such conversions are possible because the combustion reactions are essentially redox reactions and 
highly exothermic as well as highly exergonic. Electrical  energy can be obtained inde  nitely from a fuel cell as along as 
the outside supply of fuel is maintained. In hydrogen-oxygen fuel cell, the following reactions take place:
 
 Anode reaction: 
2H2(g) + 4OH-(aq) $/to$ 4H2O(l) + 4e-
 
 Cathode reaction: 
O2(g) + 2H2O(l) + 4e- $/to$ 4OH-(aq)
 
 Overall reaction: 
2H2(g) + O2(g) $/to$ 2H2O(l)
 
 The overall reaction has a value of $/Delta$H$/circ$ = -285.8 kJ and $/Delta$G$/circ$ = -237.39 kJ at 25$/circ$C per mole of H2O(l).
\end{problembox}

\begin{problembox}
\textbf{Problem 234} \hfill \textbf{[Neeraj Kumar [NRJ] EXERCISE II Q14]}
\label{q:ext-nrj_all_questions-234}

What is the standard EMF of the cell?
 
 (a) 0.615 V 
(b) 1.23 V
 
 (c) 2.46 V 
(d) 0.74 V
\end{problembox}

\begin{problembox}
\textbf{Problem 235} \hfill \textbf{[Neeraj Kumar [NRJ] EXERCISE II Q15]}
\label{q:ext-nrj_all_questions-235}

What volume of gaseous H2 (at STP), when 
combined with excess O2 in the fuel cell, is needed 
to produce 23.739 kJ of useful work under 
standard conditions?
 
 (a) 2.27 L 
(b) 4.54 L
 
 (c) 1.13 L 
(d) 2.24 L
\end{problembox}

\begin{problembox}
\textbf{Problem 236} \hfill \textbf{[Neeraj Kumar [NRJ] EXERCISE II Q16]}
\label{q:ext-nrj_all_questions-236}

Suppose the concentration of hydroxide ion in the 
cell is doubled at 298 K. The cell voltage will be
 
 (a) reduced by half
 
 (b) increased by a factor of 2
 
 (c) increased by a factor of 4
 
 (d) unchanged
\end{problembox}

\begin{problembox}
\textbf{Problem 237} \hfill \textbf{[Neeraj Kumar [NRJ] EXERCISE II Q17]}
\label{q:ext-nrj_all_questions-237}

What is the approximate value of $/Delta$S$/circ$ for the fuel 
cell reaction at 25$/circ$C?
 
 (a) - 0.1624 JK-1 
(b) -162.4 JK-1
 
 (c) +162.4 JK-1 
(d) +0.1624 JK-1
\end{problembox}

\begin{problembox}
\textbf{Problem 238} \hfill \textbf{[Neeraj Kumar [NRJ] EXERCISE II Q18]}
\label{q:ext-nrj_all_questions-238}

The theoretical e   ciency of the fuel cell is given by
 
 (a) 83.06/% 
(b) 100/%
 
 (c) 67.53/% 
(d) 97.88/%
 
 
Visit www.jeebooks.in for more
\end{problembox}

\begin{problembox}
\textbf{Problem 239} \hfill \textbf{[Neeraj Kumar [NRJ] EXERCISE II Q19]}
\label{q:ext-nrj_all_questions-239}

Which of the following reactions may occur at the 
anode of the iron hull?
 
 (a) Fe2+(aq) $/to$ Fe3+(aq) + e-
 
 (b) Fe(s) $/to$ Fe2+(aq) + 2e-
 
 (c) O2(g) + 4H3O+(aq) $/to$ 6H2O(l) + 4e-
 
 (d) 2H2O(l) $/to$ 2OH - (aq) + H2(g) + 2e-
\end{problembox}

\begin{problembox}
\textbf{Problem 240} \hfill \textbf{[Neeraj Kumar [NRJ] EXERCISE II Q20]}
\label{q:ext-nrj_all_questions-240}

What is the standard voltage generated by the 
corrosion reaction on the iron hull of a ship?
 
 (a) -1.676 V 
(b) - 0.782 V
 
 (c) 1.782 V 
(d) 1.676 V
\end{problembox}

\begin{problembox}
\textbf{Problem 241} \hfill \textbf{[Neeraj Kumar [NRJ] EXERCISE II Q21]}
\label{q:ext-nrj_all_questions-241}

Which of the following materials could not serve 
as a sacri  cial anode for lead?
 
 (a) Copper 
(b) Iron
 
 (c) Magnesium 
(d) Manganese
\end{problembox}

\begin{problembox}
\textbf{Problem 242} \hfill \textbf{[Neeraj Kumar [NRJ] EXERCISE II Q22]}
\label{q:ext-nrj_all_questions-242}

An ammeter reading shows a current of 0.50 A 
running through the iron hull. How many grams 
of Fe(s) are lost in 1.0 h?
 
 (a) 0.261 g 
(b) 0.522 g
 
 (c) 1.044 g 
(d) 0.783 g
\end{problembox}

\begin{problembox}
\textbf{Problem 243} \hfill \textbf{[Neeraj Kumar [NRJ] EXERCISE II Q23]}
\label{q:ext-nrj_all_questions-243}

A ship that has sunk to the bottom of the ocean may 
still exhibit corrosion. One di  erence in the corrosion 
reaction between ship above water and one that has 
sunk of the bottom of the ocean may be that
 
 (a)  the oxidation reaction of the sunken ship does 
not include H2O as a  reactant.
 
 (b)  the oxidation reaction of the sunken ship does 
not include O2 as a  reactant.
 
 (c)  the reduction reaction of the sunken ship does 
not include H2O as a  reactant.
 
 (d)  the reduction reaction of the sunken ship does 
not include O2 as a  reactant.
 
 
Visit www.jeebooks.in for more
\end{problembox}

\begin{problembox}
\textbf{Problem 244} \hfill \textbf{[Neeraj Kumar [NRJ] EXERCISE II Q24]}
\label{q:ext-nrj_all_questions-244}

How much of nickel is plated on the cathode per 
hour?
 
 (a) 16.43 g 
(b) 32.85 g
 
 (c) 19.7 g 
(d) 9.85 g
\end{problembox}

\begin{problembox}
\textbf{Problem 245} \hfill \textbf{[Neeraj Kumar [NRJ] EXERCISE II Q25]}
\label{q:ext-nrj_all_questions-245}

What is the thickness of the plating if the cathode 
consists of a sheet of metal 4.0 cm2, which is to be 
coated on both sides?
 
 (a) 1.38 mm 
(b) 2.76 mm
 
 (c) 0.69 mm 
(d) 23.0 mm
\end{problembox}

\begin{problembox}
\textbf{Problem 246} \hfill \textbf{[Neeraj Kumar [NRJ] EXERCISE II Q26]}
\label{q:ext-nrj_all_questions-246}

What volume of H2 at 0$/circ$C and 1 atm is formed per 
hour?
 
 (a) 6.27 L 
(b) 3.76 L
 
 (c) 2.5 L 
(d) 5.01 L
\end{problembox}

\begin{problembox}
\textbf{Problem 247} \hfill \textbf{[Neeraj Kumar [NRJ] EXERCISE II Q27]}
\label{q:ext-nrj_all_questions-247}

At the end of the electrolysis, how many grams of 
the gaseous product appear at the anode?
 
 (a) 4.48 g 
(b) 1.79 g
 
 (c) 2.69 g 
(d) 7.46 g
Comprehension VIII
Tollen's reagent is used for the detection of aldehyde when a solution of AgNO3 is added to glucose with NH4OH then 
gluconic acid is formed.
 
 Ag+ + e- $/to$ Ag; E $/circ$red = 0.8 V
 
 C6H12O6 + H2O $/to$ C6H12O7 (Gluconic acid) + 2H+ + 2e; E $/circ$oxd = -0.05 V
 
 Ag(NH3)2
+ + e- $/to$ Ag(s) + 2 NH3; E $/circ$red = 0.337 V
[Use 2.303 RT/F = 0.0592 and F/RT = 38.92 at 298 K]
\end{problembox}

\begin{problembox}
\textbf{Problem 248} \hfill \textbf{[Neeraj Kumar [NRJ] EXERCISE II Q28]}
\label{q:ext-nrj_all_questions-248}

2Ag+ + C6H12O6 + H2O $/to$ 2Ag(s) + C6H12O7 + 
2H+. Find lnK of this reaction.
 
 (a) 66.13 
(b) 58.38
 
 (c) 28.30 
(d) 46.29
\end{problembox}

\begin{problembox}
\textbf{Problem 249} \hfill \textbf{[Neeraj Kumar [NRJ] EXERCISE II Q29]}
\label{q:ext-nrj_all_questions-249}

When ammonia is added to the solution, pH is 
raised to 11. Which half-cell reaction is a  ected by 
pH and by how much?
 
 (a)  Eoxd will increase by a factor of 0.65 from 
E $/circ$oxd
 
 (b)  Eoxd will decrease by a factor of 0.65 from 
E $/circ$oxd
 
 (c)  Ered will increase by a factor of 0.65 from E $/circ$red
 
 (d)  Ered will decrease by a factor of 0.65 from 
E $/circ$red
\end{problembox}

\begin{problembox}
\textbf{Problem 250} \hfill \textbf{[Neeraj Kumar [NRJ] EXERCISE II Q30]}
\label{q:ext-nrj_all_questions-250}

Ammonia is always added in this reaction. Which 
of the following must be incorrect?
 
 (a)  NH3 combines with Ag+ to form a complex
 
 (b)  Ag(NH3)2
+ is a stronger oxidizing agent than 
Ag+
 
 (c)  In absence of ammonia, silver salt of gluconic 
acid is formed
 
 (d)  NH3 has a  ected the standard reduction 
potential of glucose-gluconic acid electrode
 
 
Visit www.jeebooks.in for more
\end{problembox}

\begin{problembox}
\textbf{Problem 251} \hfill \textbf{[Neeraj Kumar [NRJ] EXERCISE II Q31]}
\label{q:ext-nrj_all_questions-251}

The overall cell reaction is
 
 (a)  H2(g) + 2AgCl(s) $/to$ 2HCl(aq) + 2Ag(s)
 
 (b) H2(g) + 2Ag+(aq) $/to$ 2H+(aq) + 2Ag(s)
 
 (c) Ag+(aq) + Cl-(aq) $/to$ AgCl(s)
 
 (d) H2(g) + Cl2(g) $/to$ 2HCl(aq)
\end{problembox}

\begin{problembox}
\textbf{Problem 252} \hfill \textbf{[Neeraj Kumar [NRJ] EXERCISE II Q32]}
\label{q:ext-nrj_all_questions-252}

$/Delta$H$/circ$ for the cell reaction per mol of AgCl is
 
 (a) -99,974 J 
(b) -49,987 J
 
 (c) +99,974 J 
(d) +49,987 J
\end{problembox}

\begin{problembox}
\textbf{Problem 253} \hfill \textbf{[Neeraj Kumar [NRJ] EXERCISE II Q33]}
\label{q:ext-nrj_all_questions-253}

$/Delta$S$/circ$ for the cell reaction per mol of AgCl is
 
 (a) -193 J/K 
(b) -96.5 J/K
 
 (c) +96.5 J/K 
(d) +193 J/K
\end{problembox}

\begin{problembox}
\textbf{Problem 254} \hfill \textbf{[Neeraj Kumar [NRJ] EXERCISE II Q34]}
\label{q:ext-nrj_all_questions-254}

The solubility of AgCl in water at 25$/circ$C is (2.303 
RT/F = 0.058)
 
 (a) 10-10 M 
(b) 10-12 M
 
 (c) 10-5 M 
(d) 10-6 M
Comprehension X
For the reaction: Ag+(aq) + Cl-(aq)   AgCl(s). Given
Species
$/Delta$G o
f (kJ / mol) at 25$/circ$C
Ag+(aq)
+77
Cl-(aq)
-129
AgCl (s)
-109
 
 Ag+ + e- $/to$ Ag; E $/circ$ = 0.80 V
 
 Zn2+ + 2e- $/to$ Zn; E $/circ$ = -0.76 V
\end{problembox}

\begin{problembox}
\textbf{Problem 255} \hfill \textbf{[Neeraj Kumar [NRJ] EXERCISE II Q35]}
\label{q:ext-nrj_all_questions-255}

The cell representations of the above reaction and 
Ecell
$/circ$ are
 
 (a)  Ag(s) | AgCl(s) | Cl-(aq) || Ag+(aq) | Ag(s); 
Ecell
$/circ$ = 0.59 V
 
 (b)  Ag(s) | Ag+(aq) || Cl- (aq) | AgCl(s) | Ag(s); 
Ecell
$/circ$ = 0.59 V
 
 (c)  Ag(s) | Cl-(aq) || Ag+(aq) | Ag; Ecell
$/circ$ = 0.59 V
 
 (d)  Ag(s) | Ag+(aq) || Cl- (aq) | AgCl(s) | Ag(s); 
Ecell
$/circ$ = 1.18 V
\end{problembox}

\begin{problembox}
\textbf{Problem 256} \hfill \textbf{[Neeraj Kumar [NRJ] EXERCISE II Q36]}
\label{q:ext-nrj_all_questions-256}

The solubility product of AgCl at 298 K is
 
 (a) 2 $/times$ 10-10 
(b) 10-10
 
 (c) e-10 
(d) 10-5
\end{problembox}

\begin{problembox}
\textbf{Problem 257} \hfill \textbf{[Neeraj Kumar [NRJ] EXERCISE II Q37]}
\label{q:ext-nrj_all_questions-257}

A quantity of 6.539 $/times$ 10-2 g of metallic 
zinc is added to 100 ml of saturated solution of 
AgCl. The value of log [
]
[
]
Zn
Ag
2
2
+
+
 is (Zn = 65.39)
 
 (a) 5.288 
(b) 52.88
 
 (c) 528.8 
(d) 26.44
\end{problembox}

\begin{problembox}
\textbf{Problem 258} \hfill \textbf{[Neeraj Kumar [NRJ] EXERCISE II Q38]}
\label{q:ext-nrj_all_questions-258}

How many moles of Ag will be precipitated in the 
above reaction?
 
 (a) 10-3 
(b) 2 $/times$ 10-3
 
 (c) 10-6 
(d) 10-5
 
 
Visit www.jeebooks.in for more
\end{problembox}

\begin{problembox}
\textbf{Problem 259} \hfill \textbf{[Neeraj Kumar [NRJ] EXERCISE II Q39]}
\label{q:ext-nrj_all_questions-259}

Cell constant of the conductivity cell is
 
 (a) 4 cm-1 
(b) 0.4 cm-1
 
 (c) 40 cm-1 
(d) 0.04 cm-1
\end{problembox}

\begin{problembox}
\textbf{Problem 260} \hfill \textbf{[Neeraj Kumar [NRJ] EXERCISE II Q40]}
\label{q:ext-nrj_all_questions-260}

Conductivity of water (in  $/Omega$-1 cm-1) is
 
 (a) 0.4 
(b) 0.04
 
 (c) 0.0004 
(d) 0.00004
\end{problembox}

\begin{problembox}
\textbf{Problem 261} \hfill \textbf{[Neeraj Kumar [NRJ] EXERCISE II Q41]}
\label{q:ext-nrj_all_questions-261}

Volume (in L) of water in the pool is
 
 (a) 1,25,000 
(b) 12,500
 
 (c) 1250 
(d) 125
Comprehension XII
Redox reactions play a pivoted role in chemistry and biology. The values of standard redox potential (E $/circ$) of two half-cell 
reactions decide which way the reaction is expected to proceed. A simple example is Daniel cell in which zinc goes into 
solution and copper gets deposited. Given below are a set of half-cell reactions (acidic medium) along with their E $/circ$ (V 
with respect to normal hydrogen electrode) values. Using this data, obtain the correct explanations to questions below:
 
 I2 + 2e- $/to$ 2I- 
E $/circ$ = 0.54
 
 Cl2 + 2e- $/to$ 2Cl- 
E $/circ$ = 1.36
 
 Mn3+ + e- $/to$ Mn2+ 
E $/circ$ = 1.50
 
 Fe3+ + e- $/to$ Fe2+ 
E $/circ$ = 0.77
 
 O2 + 4H+ + 4e- $/to$ 2H2O 
E $/circ$ = 1.23
\end{problembox}

\begin{problembox}
\textbf{Problem 262} \hfill \textbf{[Neeraj Kumar [NRJ] EXERCISE II Q42]}
\label{q:ext-nrj_all_questions-262}

Among the following, identify the correct 
statement?
 
 (a) Chloride ion is oxidized by O2
 
 (b) Fe2+ is oxidized by iodine
 
 (c) Iodide ion is oxidized by chlorine
 
 (d) Mn2+ is oxidized by chlorine
\end{problembox}

\begin{problembox}
\textbf{Problem 263} \hfill \textbf{[Neeraj Kumar [NRJ] EXERCISE II Q43]}
\label{q:ext-nrj_all_questions-263}

While Fe3+ is stable, Mn3+ is not stable in acid 
solution because
 
 (a) O2 oxidizes Mn2+ to Mn3+
 
 (b)  O2 oxidizes both Mn2+ to Mn3+ and Fe2+ to 
Fe3+
 
 (c) Fe3+ oxidizes H2O to O2
 
 (d) Mn3+ oxidizes H2O to O2
Comprehension XIII
The concentration of potassium ions inside a biological cell is at least 20 times higher than the  outside. The resulting 
potential di  erence across the cell is important in several processes such as transmission of nerve impulses and maintaining 
the ion balance. A simple model for such a concentration cell involving a metal 'M' is
M(s) | M+(aq; 0.05 molar) || M+(aq; 1 molar) | M(s)
For the above electrolytic cell, the magnitude of the cell potential, |Ecell| = 70 mV.
\end{problembox}

\begin{problembox}
\textbf{Problem 264} \hfill \textbf{[Neeraj Kumar [NRJ] EXERCISE II Q44]}
\label{q:ext-nrj_all_questions-264}

For the above cell
 
 (a) Ecell < 0; $/Delta$G > 0 
(b) Ecell > 0; $/Delta$G < 0
 
 (c) Ecell < 0; $/Delta$G$/circ$ > 0 
(d) Ecell > 0; $/Delta$G$/circ$< 0
\end{problembox}

\begin{problembox}
\textbf{Problem 265} \hfill \textbf{[Neeraj Kumar [NRJ] EXERCISE II Q45]}
\label{q:ext-nrj_all_questions-265}

If the 0.05 molar solution of M+ is replaced by a 
0.0025 molar M+ solution, then the magnitude of 
the cell potential would be
 
 (a) 35 mV 
(b) 70 mV
 
 (c) 140 mV 
(d) 700 mV
 
 
Visit www.jeebooks.in for more
\end{problembox}

\begin{problembox}
\textbf{Problem 266} \hfill \textbf{[Neeraj Kumar [NRJ] EXERCISE II Q46]}
\label{q:ext-nrj_all_questions-266}

The value of $/Delta$G (kJ mol-1) for the given cell is 
(take 1 F = 96500 C mol-1)
 
 (a) -5.7 
(b) 5.7
 
 (c) 11.4 
(d) -11.4
\end{problembox}

\begin{problembox}
\textbf{Problem 267} \hfill \textbf{[Neeraj Kumar [NRJ] EXERCISE II Q47]}
\label{q:ext-nrj_all_questions-267}

The solubility product (Ksp; mol3 dm-3) of MX2 
at 298 K based on the information available for 
the given concentration cell is (take 2.303 $/times$ R $/times$ 
298/F = 0.059)
 
 (a) 1 $/times$ 10-15 
(b) 4 $/times$ 10-15
 
 (c) 1 $/times$ 10-12 
(d) 4 $/times$ 10-12
Comprehension XV
Conductance measurements are frequently employed to   nd the end points of acid-base and other titrations. The 
principle involved is that electrical conductance of a solution depends upon the number and mobility of ions. In the 
conductometric titrations, the conductance of the resulting solution is measured at di  erent stages on adding some 
volume of the standard solution and then a graph is plotted from the experimental observations to get the end point. 
Conductometric titrations have several advantages. Coloured solutions, titration of weak acid and weak base, etc. 
cannot be titrated by normal methods, but can be titrated successfully by this method. Further, no special care is needed 
in titration because the end point is determined graphically.
\end{problembox}

\begin{problembox}
\textbf{Problem 268} \hfill \textbf{[Neeraj Kumar [NRJ] EXERCISE II Q48]}
\label{q:ext-nrj_all_questions-268}

Which of the following graph truly represents the 
titration of HCl solution against NaOH solution?
 
 (a) 
Cond.
VNaOH
 
 (b) 
Cond.
VNaOH
 
 (c) 
Cond.
VNaOH
 
 (d) 
Cond.
VNaOH
\end{problembox}

\begin{problembox}
\textbf{Problem 269} \hfill \textbf{[Neeraj Kumar [NRJ] EXERCISE II Q49]}
\label{q:ext-nrj_all_questions-269}

Which of the following graph truly represents the 
titration of CH3COOH solution against NaOH 
solution?
 
 (a) 
Cond.
VNaOH
 
 (b) 
Cond.
VNaOH
 
 (c) 
Cond.
VNaOH
 
 (d) 
Cond.
VNaOH
 
 
Visit www.jeebooks.in for more
\end{problembox}

\begin{problembox}
\textbf{Problem 270} \hfill \textbf{[Neeraj Kumar [NRJ] EXERCISE II Q50]}
\label{q:ext-nrj_all_questions-270}

Which of the following graph truly represents 
the titration of HCl solution against NH4OH 
solution?
 
 (a) 
Cond.
VNH4OH
 
 (b) 
Cond.
VNH4OH
 
 (c) 
Cond.
VNH4OH
 
 (d) 
Cond.
VNH4OH
\end{problembox}

\begin{problembox}
\textbf{Problem 271} \hfill \textbf{[Neeraj Kumar [NRJ] EXERCISE II Q51]}
\label{q:ext-nrj_all_questions-271}

Which of the following graph truly represents the 
titration of a solution containing a mixture of 
HCl and CH3COOH against NaOH solution?
 
 (a) 
Cond.
VNaOH
 
 (b) 
Cond.
VNaOH
 
 (c) 
Cond.
VNaOH
 
 (d) 
Cond.
VNaOH
\end{problembox}

\begin{problembox}
\textbf{Problem 272} \hfill \textbf{[Neeraj Kumar [NRJ] EXERCISE II Q52]}
\label{q:ext-nrj_all_questions-272}

Which of the following graph truly represents 
the titration of AgNO3 solution against KCl 
solution?
 
 (a) 
Cond.
VKCI
 
 (b) 
Cond.
VKCI
 
 (c) 
Cond.
VKCI
 
 (d) 
Cond.
VKCI
 
 
Visit www.jeebooks.in for more
\end{problembox}

\begin{problembox}
\textbf{Problem 273} \hfill \textbf{[Neeraj Kumar [NRJ] EXERCISE II Q53]}
\label{q:ext-nrj_all_questions-273}

Which of the following graph truly represents the 
titration of CH3COOH solution against NH4OH 
solution?
 
 (a) 
Cond.
VNH4OH
 
 (b) 
Cond.
VNH4OH
 
 (c) 
Cond.
VNH4OH
 
 (d) 
Cond.
VNH4OH
Section D (Assertion - Reason)
The following questions consist of two statements. Mark
 (a) If both statements are CORRECT, and Statement 
II is the CORRECT explanation of Statement I. 
 (b) If both statements are CORRECT, and Statement 
II is NOT the CORRECT explanation of 
Statement I. 
 (c) If Statement I is CORRECT, but  Statement II is 
INCORRECT. 
 (d) If Statement I is INCORRECT, but Statement II 
is CORRECT.
\end{problembox}

\begin{problembox}
\textbf{Problem 274} \hfill \textbf{[Neeraj Kumar [NRJ] EXERCISE II Q1]}
\label{q:ext-nrj_all_questions-274}

Statement I: Electrolysis of CuCl2(aq) gives 1 mole 
of Cu and 1 mole of Cl2 by the passage of suitable 
charge.
 
 Statement II: Equal equivalents of Cu and Cl2 are 
formed during the passage of same charge.
\end{problembox}

\begin{problembox}
\textbf{Problem 275} \hfill \textbf{[Neeraj Kumar [NRJ] EXERCISE II Q2]}
\label{q:ext-nrj_all_questions-275}

Statement I: Zinc metal can displace Ag metal from 
a solution containing the complex [Ag((CN)2]-.
 
 Statement II: EZn
Zn
2+
$/circ$
,
 is greater than EAg
Ag
+
$/circ$
,
.
\end{problembox}

\begin{problembox}
\textbf{Problem 276} \hfill \textbf{[Neeraj Kumar [NRJ] EXERCISE II Q3]}
\label{q:ext-nrj_all_questions-276}

Statement I: In electrolysis, the quantity of 
electricity needed for depositing 1 mole of silver is 
di  erent from that required for 1 mole copper.
 
 Statement II: The atomic masses of silver and 
copper are di  erent.
\end{problembox}

\begin{problembox}
\textbf{Problem 277} \hfill \textbf{[Neeraj Kumar [NRJ] EXERCISE II Q4]}
\label{q:ext-nrj_all_questions-277}

Statement I: In an electrochemical cell, anode and 
cathode are, respectively, negative and positive 
electrode.
 
 Statement II: At anode, oxidation takes place and 
at cathode reduction takes place.
\end{problembox}

\begin{problembox}
\textbf{Problem 278} \hfill \textbf{[Neeraj Kumar [NRJ] EXERCISE II Q5]}
\label{q:ext-nrj_all_questions-278}

Statement I: A metal having negative reduction 
potential when dipped in the solution of its own 
ions has a tendency to pass into the solution.
 
 Statement II: Metals having negative reduction 
potential have large hydration energy.
\end{problembox}

\begin{problembox}
\textbf{Problem 279} \hfill \textbf{[Neeraj Kumar [NRJ] EXERCISE II Q6]}
\label{q:ext-nrj_all_questions-279}

Statement I: Speci  c conductance decreases with 
dilution while molar conductance increases.
 
 Statement II: On dilution, number of ions per 
unit volume decreases but total number of ions 
increases considerably.
\end{problembox}

\begin{problembox}
\textbf{Problem 280} \hfill \textbf{[Neeraj Kumar [NRJ] EXERCISE II Q7]}
\label{q:ext-nrj_all_questions-280}

Statement I: When acidi  ed zinc sulphate solution 
is electrolysed between zinc electrodes, it is zinc 
that is deposited at the cathode and evolution of 
hydrogen gas does not take place.
 
 Statement II: The electrode potential of zinc 
becomes less negative than hydrogen as the 
overvoltage for the hydrogen evolution on zinc is 
quite large.
\end{problembox}

\begin{problembox}
\textbf{Problem 281} \hfill \textbf{[Neeraj Kumar [NRJ] EXERCISE II Q8]}
\label{q:ext-nrj_all_questions-281}

Statement I: The conductivity of solutions of 
di  erent electrolytes in the same solvent and at a 
given temperature is same.
 
 Statement II: The conductivity depends on the 
charge and size of the ions, the concentrations 
of ions and ease with which the ions move under 
potential gradient.
\end{problembox}

\begin{problembox}
\textbf{Problem 282} \hfill \textbf{[Neeraj Kumar [NRJ] EXERCISE II Q9]}
\label{q:ext-nrj_all_questions-282}

Statement I: Sodium ions are discharged in 
preference to hydrogen ions at a mercury cathode.
 
 Statement II: Na+ is a strong reducing agent in 
comparison to H+ ion.
  
 
 
Visit www.jeebooks.in for more
\end{problembox}

\begin{problembox}
\textbf{Problem 283} \hfill \textbf{[Neeraj Kumar [NRJ] EXERCISE II Q10]}
\label{q:ext-nrj_all_questions-283}

Statement I: The conductivity of metals decreases 
while that of electrolytic solution increases with 
increase in temperature.
 
 Statement II: Electrons in metals are very tightly 
held by the nucleus and are not free to move.
\end{problembox}

\begin{problembox}
\textbf{Problem 284} \hfill \textbf{[Neeraj Kumar [NRJ] EXERCISE II Q11]}
\label{q:ext-nrj_all_questions-284}

Statement I: An electrochemical cell can be set up 
only when the redox reaction is spontaneous.
 
 Statement II: A reaction is spontaneous if free 
energy change at constant temperature and 
pressure is negative.
\end{problembox}

\begin{problembox}
\textbf{Problem 285} \hfill \textbf{[Neeraj Kumar [NRJ] EXERCISE II Q12]}
\label{q:ext-nrj_all_questions-285}

Statement I: In the Daniel cell, if concentrations 
of Cu2+ and Zn2+ ions are doubled, the EMF of 
cell does not change.
 
 Statement II: If the concentration of ions in 
contact with the metal is doubled, the electrode 
potential will be doubled.
\end{problembox}

\begin{problembox}
\textbf{Problem 286} \hfill \textbf{[Neeraj Kumar [NRJ] EXERCISE II Q13]}
\label{q:ext-nrj_all_questions-286}

Statement I: KCl, NaCl, NH4Cl, etc., cannot be 
used in the salt bridge of a cell containing silver ion.
 
 Statement II: A salt bridge contains concentrated 
solution of an inert electrolyte like KCl, KNO3, 
K2SO4, etc., in agar-agar.
\end{problembox}

\begin{problembox}
\textbf{Problem 287} \hfill \textbf{[Neeraj Kumar [NRJ] EXERCISE II Q14]}
\label{q:ext-nrj_all_questions-287}

Statement I: During electrolysis of aqueous 
sodium acetate solution, the molar ratio of gases 
formed at cathode and anode is 1:3.
 
 Statement II: Acetate ion discharges at anode and 
H+ ion, at cathode.
\end{problembox}

\begin{problembox}
\textbf{Problem 288} \hfill \textbf{[Neeraj Kumar [NRJ] EXERCISE II Q15]}
\label{q:ext-nrj_all_questions-288}

Statement I: Electrode potential of any electrode 
will change on changing any of its intensive 
properties.
 
 Statement II: Any intensive property of system 
changes on changing any of its properties, whether 
intensive or extensive.
Section E (Column Match)
\end{problembox}

\begin{problembox}
\textbf{Problem 289} \hfill \textbf{[Neeraj Kumar [NRJ] EXERCISE II Q1]}
\label{q:ext-nrj_all_questions-289}

Match Column I with Column II
Column I
Column II 
(Electrolysis product 
using inert electrodes)
(A)  Dilute solution of 
HCl
(P)  O2 evolved at anode
(B)  Dilute solution of 
NaCl
(Q)  H2 evolved at cathode
(C)  Concentrated 
solution of NaCl
(R)  Cl2 evolved at anode
(D)  AgNO3 solution
(S)  Ag deposited at cathode
\end{problembox}

\begin{problembox}
\textbf{Problem 290} \hfill \textbf{[Neeraj Kumar [NRJ] EXERCISE II Q2]}
\label{q:ext-nrj_all_questions-290}

An aqueous solution of 'X' is added slowly to an 
aqueous solution of 'Y' as shown in List I. The 
variation in conductivity of these reactions is 
given in List II. Match List I with List II.
List I  (X + Y)
List II
(A)  (C2H5)3N + 
CH3COOH
(P)  Conductivity decreases 
and then increases
(B)  KI(0.1 M) + 
AgNO3(0.01 M)
(Q)  Conductivity decreases 
and then does not change 
much
(C)  CH3COOH + 
KOH
(R)  Conductivity increases 
and then does not change 
much
(D)  NaOH + HI
(S)  Conductivity does not 
change much and then 
increases
\end{problembox}

\begin{problembox}
\textbf{Problem 291} \hfill \textbf{[Neeraj Kumar [NRJ] EXERCISE II Q3]}
\label{q:ext-nrj_all_questions-291}

The standard reduction potential data at 25$/circ$C is 
given below:
 
 E $/circ$ (Fe3+, Fe2+) = +0.77 V
 
 E $/circ$ (Fe2+, Fe) = -0.44 V
 
 E $/circ$ (Cu2+, Cu) = +0.34 V
 
 E $/circ$ (Cu+, Cu) = +0.52 V
 
 E $/circ$ [O2(g) + 4H+ + 4e- $/to$ 2H2O] = +1.23 V
 
 E $/circ$ [O2(g) + 2H2O + 4e- $/to$ 4OH-] = +0.40 V
 
 E $/circ$ (Cr3+, Cr) = -0.74 V
 
 E $/circ$ (Cr2+, Cr) = -0.91 V
 
 Match E $/circ$ of the redox pair in List I with the values 
given in List II.
List I
List II
(A)  E $/circ$ (Fe3+, Fe)
(P)  -0.70 V
(B)  E $/circ$ (4H2O $/to$ 4H+ + 4 OH-)
(Q)  -0.4 V
(C)  E $/circ$ (Cu2+ + Cu $/to$ 2Cu+)
(R)  -0.04 V
(D)  E $/circ$ (Cr3+, Cr2+)
(S)  -0.83 V
\end{problembox}

\begin{problembox}
\textbf{Problem 292} \hfill \textbf{[Neeraj Kumar [NRJ] EXERCISE II Q4]}
\label{q:ext-nrj_all_questions-292}

Match the following
Column I
Column II
(A)  Concentration cell (P)  Ag | AgCl | Cl- || Ag+ | Ag
(B)  Spontaneous cell 
reaction
(Q)  Ag | AgCl | Cl- || Br- | 
AgBr | Ag
(C)  Non-spontaneous 
cell reaction
(R)  Ag | Ag+ (0.1 M) || Ag+ 
(1.0 M) | Ag
(S)  Ag | AgCl | Cl- (0.1 M) || Cl-
(1.0 M)|AgCl | Ag
 
 
Visit www.jeebooks.in for more
\end{problembox}

\begin{problembox}
\textbf{Problem 293} \hfill \textbf{[Neeraj Kumar [NRJ] EXERCISE II Q1]}
\label{q:ext-nrj_all_questions-293}

Percentage of aniline hydrochloride hydrolysed in 
its M/40 solution at 25$/circ$C is (Given: EC H NH
HCl H
6
5
2
2
.
|
= -0.18   V, 2.303 RT/F = 0.06)
\end{problembox}

\begin{problembox}
\textbf{Problem 294} \hfill \textbf{[Neeraj Kumar [NRJ] EXERCISE II Q2]}
\label{q:ext-nrj_all_questions-294}

The standard reduction potential for Cu2+|Cu 
electrode is +0.34 V. The solubility product of 
Cu(OH)2 is 1.0 $/times$ 10-19. The pH of solution at 
which the reduction potential for the above 
electrode becomes 0.31 V, is (2.303 RT/F = 0.06)
\end{problembox}

\begin{problembox}
\textbf{Problem 295} \hfill \textbf{[Neeraj Kumar [NRJ] EXERCISE II Q3]}
\label{q:ext-nrj_all_questions-295}

The reduction potential at 25$/circ$C for Fe3+|Fe2+ 
electrode is +0.718 V. If EFe
Fe
2+
$/circ$
|
= -0.44 V and 
EFe
Fe
3+
$/circ$
|
= -0.04 V, the ratio of molar concentrations 
of Fe2+ to Fe3+ ions in solution is (2.303 RT/F = 
0.06, log 5 = 0.7)
\end{problembox}

\begin{problembox}
\textbf{Problem 296} \hfill \textbf{[Neeraj Kumar [NRJ] EXERCISE II Q4]}
\label{q:ext-nrj_all_questions-296}

If [Fe3+] at equilibrium, when potassium iodide is 
added to a solution of Fe3+ init ially at 0.50 M until 
[I-] = 1.0 M, is x $/times$ 10-5 M, the value of x is (Given 
EFe
Fe
3
2
+
+
$/circ$
|
 
= 0.77 V, EI I
2 | -
$/circ$
 
= 0.53 V, 2.303 RT/F = 0.06)
\end{problembox}

\begin{problembox}
\textbf{Problem 297} \hfill \textbf{[Neeraj Kumar [NRJ] EXERCISE II Q5]}
\label{q:ext-nrj_all_questions-297}

If it is desired to construct the following galvanic 
cell:
 
   Ag(s) | Ag+ (saturated AgI) || Ag+ (saturated 
AgCl, x $/times$ 10-4 M Cl-) | Ag(s)
 
 to have ECell = 0.102 V, what should be the value of 
x to get [Cl-], which must be present in the cathodic 
half-cell to achieve the desired EMF. Given Ksp 
of AgCl and AgI are 1.8 $/times$ 10-10 and 8.1 $/times$ 10-17, 
respectively. (2.303 RT/F = 0.06, log 2 = 0.3)
\end{problembox}

\begin{problembox}
\textbf{Problem 298} \hfill \textbf{[Neeraj Kumar [NRJ] EXERCISE II Q6]}
\label{q:ext-nrj_all_questions-298}

An alloy of lead (valency = 2)-thallium (valency = 
1) containing 70/% Pb and 30/% Tl, by weight, can 
be electroplated onto a cathode from a perchloric 
acid solution. How many hours (approx.) would 
be required to deposit 5.0 g of this alloy at a 
current of 1.10 A? (Pb = 208, Tl = 204)
\end{problembox}

\begin{problembox}
\textbf{Problem 299} \hfill \textbf{[Neeraj Kumar [NRJ] EXERCISE II Q7]}
\label{q:ext-nrj_all_questions-299}

Iridium was plated from a solution containing 
IrCl6
y for 2.0 h with a current of 0.075 A. The 
Iridium deposited on the cathode weighed 0.36 g. 
If the oxidation state of Ir in IrCl6
y is x, then the 
value of (x + y) is (Ir = 192)
\end{problembox}

\begin{problembox}
\textbf{Problem 300} \hfill \textbf{[Neeraj Kumar [NRJ] EXERCISE II Q8]}
\label{q:ext-nrj_all_questions-300}

The electrolysis of cold sodium chloride solution 
produces sodium hypochlorite by reacting NaOH 
and Cl2 thoroughly. How long (in days) will a 
cell operate to produce 10 L of 7.45/% (by mass) 
solution of NaClO if the cell current is 2.5 A? 
Assume that the density of solution is 1.0 g/ml.
\end{problembox}

\begin{problembox}
\textbf{Problem 301} \hfill \textbf{[Neeraj Kumar [NRJ] EXERCISE II Q9]}
\label{q:ext-nrj_all_questions-301}

A volume of 500 ml of 0.1 M -- CuSO4 solution is 
electrolysed for 5 min at a current of 0.161 A. If 
Cu is produced at one electrode and oxygen at the 
other, the approximate pH of the   nal solution is
\end{problembox}

\begin{problembox}
\textbf{Problem 302} \hfill \textbf{[Neeraj Kumar [NRJ] EXERCISE II Q10]}
\label{q:ext-nrj_all_questions-302}

A test for complete removal of Cu2+ ions from a 
solution of Cu2+ is to add NH3(aq). A blue colour 
signi  es the  formation of complex [Cu(NH3)4]2+ 
having Kf = 1.1  $/times$  1013 and thus con  rms the 
presence of Cu2+ in solution. 250 ml of 0.1 M -- 
CuSO4 is electrolysed by passing a current of 5 A 
for 1351 s. After passage of this charge, su   cient 
quantity of NH3 is added to electrolysed solution 
maintaining [NH3] = 0.10 M. If [Cu(NH3)4] 2+ is 
detectable up to its concentration as low as 1 $/times$ 
10-5 M, would a blue colour be shown by the 
electrolysed solution on addition of NH3. Mark 
'1', if the answer is 'yes' and mark '2', if the 
answer is 'no'.
\end{problembox}

\begin{problembox}
\textbf{Problem 303} \hfill \textbf{[Neeraj Kumar [NRJ] EXERCISE II Q11]}
\label{q:ext-nrj_all_questions-303}

By passing a certain amount of charge through 
NaCl solution, 9.08 L of chlorine gas were 
liberated at STP. When the same amount of charge 
is passed through a nitrate solution of metal M, 
52.8 g of the metal was deposited. If the speci  c 
heat of metal is 0.032 Cal/ C-g, the valency of 
metal is
\end{problembox}

\begin{problembox}
\textbf{Problem 304} \hfill \textbf{[Neeraj Kumar [NRJ] EXERCISE II Q12]}
\label{q:ext-nrj_all_questions-304}

The electrode reactions for charging of a lead 
storage battery are:
 
  PbSO4 + 2e $/to$ Pb + SO4
2-
 
   PbSO4 + 2H2O $/to$ PbO2 + SO4
2- + 4H+ + 2e
 
 The electrolyte in the battery is an aqueous solution 
of sulphuric acid. Before charging, the speci  c 
gravity of the liquid was found to be 1.10 (16/% 
H2SO4 by wt.). After charging for 965
9 h, the 
speci  c gravity of the liquid was found to be 1.42 
(40/% H2SO4 by weight). If the battery contained 
2 L of the liquid and the volume remains constant 
during charge, the average current (in A) used for 
charging the battery is
 
 
Visit www.jeebooks.in for more
\end{problembox}

\begin{problembox}
\textbf{Problem 305} \hfill \textbf{[Neeraj Kumar [NRJ] EXERCISE II Q13]}
\label{q:ext-nrj_all_questions-305}

A dilute solution of KCl was placed between two 
platinum electrodes, 12 cm apart, across which a 
potential of 1.93 volts was applied. How far (in 
cm) would the K+ ion move in 20 h at 25$/circ$C? Ionic 
conductance of K+ ion at in  nite dilution at 25$/circ$C 
is 7.5 $/times$ 10-3 mho m2 mol-1.
\end{problembox}

\begin{problembox}
\textbf{Problem 306} \hfill \textbf{[Neeraj Kumar [NRJ] EXERCISE II Q14]}
\label{q:ext-nrj_all_questions-306}

In the re  ning of silver by electrolytic method, what 
will be the   nal mass (in gm) of 72.8 g silver anode 
(60/% pure, by weight), if 9.65 A current is passed for 
1 h? (Ag = 108)
\end{problembox}

\begin{problembox}
\textbf{Problem 307} \hfill \textbf{[Neeraj Kumar [NRJ] EXERCISE II Q15]}
\label{q:ext-nrj_all_questions-307}

The pH of 100 L of concentrated KCl solution 
after the electrolysis for 10 s using 9.65 A at 
298 K is
Four-digit Integer Type
\end{problembox}

\begin{problembox}
\textbf{Problem 308} \hfill \textbf{[Neeraj Kumar [NRJ] EXERCISE II Q1]}
\label{q:ext-nrj_all_questions-308}

The speci  c conductivity of a saturated solution 
of AgCl is 2.80 $/times$ 10-4 mho m-1 at 25$/circ$C. If $/lambda$Ag+
$/circ$
= 
6.19 $/times$ 10-3 mho m2 mol-1 and $/lambda$Cl-
$/circ$
= 7.81 $/times$ 10-3 
mho m2 mol-1, the solubility of silver chloride (in 
order of 10-5 g l-1) at 25$/circ$C, is
\end{problembox}

\begin{problembox}
\textbf{Problem 309} \hfill \textbf{[Neeraj Kumar [NRJ] EXERCISE II Q2]}
\label{q:ext-nrj_all_questions-309}

The electrode potential (in millivolts) of 2Ag(s) + 
S2-(aq) $/to$ Ag2S(s) + 2e- in a solution bu  ered at 
pH = 3 and which is also saturated with 0.1 M -- 
H2S, is (Given: for H2S, Ka1 = 10-8; Ka2 = 10-13, Ksp 
of Ag2S = 4 $/times$ 10-48, E $/circ$Ag
+
|Ag = 0.80 V, log 2 = 0.3, 
2.303 RT/F = 0.06)
\end{problembox}

\begin{problembox}
\textbf{Problem 310} \hfill \textbf{[Neeraj Kumar [NRJ] EXERCISE II Q3]}
\label{q:ext-nrj_all_questions-310}

An excess of liquid mercury is added to an 
acidi  ed solution of 10-3 M -- Fe3+. It is found 
that 10/% of Fe3+ remains at equilibrium at 25$/circ$C. 
The value of EHg
Hg
2
2+
$/circ$
|
(in mV), assuming that the 
only reaction that occurs is: 2Hg + 2Fe3+ $/to$ Hg2
2+ 
+ 2Fe2+, is (Given: EFe
Fe
3
2
+
+
$/circ$
|
= 0.7724 V, log2 = 0.3, 
log3 = 0.48, 2.303 RT/F = 0.06)
\end{problembox}

\begin{problembox}
\textbf{Problem 311} \hfill \textbf{[Neeraj Kumar [NRJ] EXERCISE II Q4]}
\label{q:ext-nrj_all_questions-311}

Determine potential (in mV) of the cell: Pt | Fe2+, 
Fe3+ || Cr2O7
2-, Cr3+, H+ | Pt in which [Fe2+] = 0.75 
M, [Fe3+] = 0.75 M, [Cr2O7
2-] = 2 M, [Cr3+] = 4 M 
and [H+] = 1 M. Given:
 
  Fe3+ + e- $/to$ Fe2+; E  = 0.77 V
 
   14 H+ + 6e + Cr2O7
2- $/to$  2Cr3+ + 7H2O; 
E  = 1.35 V
 
  2.303 RT/F = 0.06, log 2 = 0.3, log 3 = 0.48
\end{problembox}

\begin{problembox}
\textbf{Problem 312} \hfill \textbf{[Neeraj Kumar [NRJ] EXERCISE II Q5]}
\label{q:ext-nrj_all_questions-312}

An alloy weighing 2.70 mg of Pb-Ag was dissolved 
in desired amount of HNO3 and volume was made 
250 ml. A silver electrode was dipped in solution 
and Ecell of the cell:
 
  Pt, H2(1 bar) | H+(1 M) || Ag+ | Ag
 
 was 0.50 V at 298 K. The percentage of lead in 
alloy is (Given: E Ag+ | Ag = 0.80 V, Ag = 108, 2.303 
RT/F = 0.06)
\end{problembox}

\begin{problembox}
\textbf{Problem 313} \hfill \textbf{[Neeraj Kumar [NRJ] EXERCISE II Q6]}
\label{q:ext-nrj_all_questions-313}

Given E $/circ$ = 0.08 V for Fe3+(cyt b) | Fe2+ (cyt  b) 
couple and E $/circ$ = 0.20 V for Fe3+(cyb c1) | Fe2+(cyt 
c1) couple, where E $/circ$ represents the standard state 
reduction potentials at pH = 7.0 at 25$/circ$C and cyt is 
an abbreviation for cytochromes. The value of Keq 
for the reaction
 
   Fe3+ (cyt c1) + Fe2+ (cyt b)   Fe2+ (cyt c1) + Fe3+ 
(cyt b)
 
 is (2.303 RT/F = 0.06)
\end{problembox}

\begin{problembox}
\textbf{Problem 314} \hfill \textbf{[Neeraj Kumar [NRJ] EXERCISE II Q7]}
\label{q:ext-nrj_all_questions-314}

Calculate potential (in mV) of the cell: Cu | Mn(s) | 
MnCl2(0.001 M), HCl(0.01 M) | O2(0.25 bar) | Pt | 
Cu. Given, E  = -1.185 V for the Mn2+ | Mn couple 
and 1.229 V for the O2 | H2O, H+ couple. (2.303 
RT/F = 0.06, log 2 = 0.3)
\end{problembox}

\begin{problembox}
\textbf{Problem 315} \hfill \textbf{[Neeraj Kumar [NRJ] EXERCISE II Q8]}
\label{q:ext-nrj_all_questions-315}

The temperature coe   cient of the cell: Zn | ZnCl2 
| AgCl(s) | Ag is + 0.001 V/K. The entropy change 
(in J/K) accompanying this reaction, at 298 K is
\end{problembox}

\begin{problembox}
\textbf{Problem 316} \hfill \textbf{[Neeraj Kumar [NRJ] EXERCISE II Q9]}
\label{q:ext-nrj_all_questions-316}

For the reaction: 4Al(s) + 3O2(g) + 6H2O(l) + 
4OH-(aq)   4Al(OH)4
-(aq); Ecell
$/circ$ = 2.5 V. If $/Delta$fGo 
for H2O(l) and OH-(aq) are -280.0 and -156.25 
kJ/mol, respectively, the magnitude of $/Delta$fGo for 
Al(OH)4
-(aq) (in kJ/mol) is
\end{problembox}

\begin{problembox}
\textbf{Problem 317} \hfill \textbf{[Neeraj Kumar [NRJ] EXERCISE II Q10]}
\label{q:ext-nrj_all_questions-317}

In a Zn-MnO2 Cell, the anode is made up of 
Zn and cathode of carbon rod surrounded by a 
mixture of MnO2, Carbon, NH4Cl and ZnCl2 in 
aqueous base. If 8.7 g MnO2 is present in cathodic 
compartment, how many days the dry cell will 
continue to give a current of 3.99 $/times$ 10-3 A?
\end{problembox}

\begin{problembox}
\textbf{Problem 318} \hfill \textbf{[Neeraj Kumar [NRJ] EXERCISE II Q11]}
\label{q:ext-nrj_all_questions-318}

An amount of 19 g of molten SnCl2 is electrolysed 
for some time. Inert electrodes are used. 1.19 g of 
tin is deposited at the cathode. No substance is 
lost during electrolysis. If the ratio of the masses 
of SnCl2 and SnCl4 after electrolysis is x : 261, the 
value of x is (Sn = 119)
 
 
Visit www.jeebooks.in for more
\end{problembox}

\begin{problembox}
\textbf{Problem 319} \hfill \textbf{[Neeraj Kumar [NRJ] Exercise I Q12]}
\label{q:ext-nrj_all_questions-319}

An object whose surface area is 80 cm2 is to be 
plated with an even layer of gold 8.0 $/times$ 10-4 cm 
thick. The density of gold is 19.7  g/cm3. The 
object is placed in a solution of Au(NO3)3 and a 
current of 2.4 A is applied. The time (in seconds) 
required for the electroplating to be completed, 
assuming that the layer of gold builds up evenly, is 
(Au = 197)
\end{problembox}

\begin{problembox}
\textbf{Problem 320} \hfill \textbf{[Neeraj Kumar [NRJ] Exercise I Q13]}
\label{q:ext-nrj_all_questions-320}

We have taken a saturated solution of AgBr. Ksp 
of AgBr is 12 $/times$ 10-14. If 10-7 moles of AgNO3 is 
added to 1 litre of this solution,   nd conductivity 
(speci  c conductance) of this solution in terms of 
10-7 S$/cdot$m-1 units. Given $/lambda$o(Ag+) = 6 $/times$ 10-3 S m2 
mol-1, $/lambda$$/circ$(Br-) = 8 $/times$ 10-3 S m2 mol-1, $/lambda$$/circ$(NO3
-) = 
7 $/times$ 10-3 S$/cdot$m2 mol-1.
\end{problembox}

\begin{problembox}
\textbf{Problem 321} \hfill \textbf{[Neeraj Kumar [NRJ] Exercise I Q14]}
\label{q:ext-nrj_all_questions-321}

On passing electricity through nitrobenzene 
solution, it is converted into azobenzene. The 
mass of azobenzene (in mg) produced, if the 
same quantity of electricity produces oxygen just 
su   cient to burn 96 g of fullerene (C60), is
\end{problembox}

\begin{problembox}
\textbf{Problem 322} \hfill \textbf{[Neeraj Kumar [NRJ] Exercise I Q15]}
\label{q:ext-nrj_all_questions-322}

The conductivity of saturated solution of 
sparingly soluble salt, Ba3(PO4)2, is 1.2 $/times$ 10--5 
$/Omega$--1cm--1. The limiting equivalent conductances 
of BaCl2, K3PO4 and KCl are 160, 140 and 100 
$/Omega$--1cm2
 eq--1, respectively. The ksp of Ba3(PO4)2 
(in the order of 10--25) is
\end{problembox}

\begin{problembox}
\textbf{Problem 323} \hfill \textbf{[Neeraj Kumar [NRJ] Exercise I Q30]}
\label{q:ext-nrj_all_questions-323}

(a)
Galvanic Cell
\end{problembox}

\begin{problembox}
\textbf{Problem 324} \hfill \textbf{[Neeraj Kumar [NRJ] Exercise I Q65]}
\label{q:ext-nrj_all_questions-324}

(a)
Electrolysis
\end{problembox}

\begin{problembox}
\textbf{Problem 325} \hfill \textbf{[Neeraj Kumar [NRJ] Exercise I Q99]}
\label{q:ext-nrj_all_questions-325}

(b) 100. (c) 101. (a) 102. (a) 103. (b) 104. (b) 105. (b)
Conductance
\end{problembox}

\begin{problembox}
\textbf{Problem 326} \hfill \textbf{[Neeraj Kumar [NRJ] Exercise I Q106]}
\label{q:ext-nrj_all_questions-326}

(b) 107. (a) 108. (c) 109. (b) 110. (a) 111. (b) 112. (c) 113. (c) 114. (a) 115. (d)
\end{problembox}

\begin{problembox}
\textbf{Problem 327} \hfill \textbf{[Neeraj Kumar [NRJ] Exercise I Q116]}
\label{q:ext-nrj_all_questions-327}

(a) 117. (b) 118. (d) 119. (c) 120. (d) 121. (a) 122. (c) 123. (a) 124. (d) 125. (a)
ANSWER KEYS
 
 
Visit www.jeebooks.in for more
\end{problembox}

\begin{problembox}
\textbf{Problem 328} \hfill \textbf{[Neeraj Kumar [NRJ] Exercise II Q80]}
\label{q:ext-nrj_all_questions-328}

(b)
Section B (One or More than one Correct)
\end{problembox}

\begin{problembox}
\textbf{Problem 329} \hfill \textbf{[Neeraj Kumar [NRJ] Exercise II Q15]}
\label{q:ext-nrj_all_questions-329}

(a), (b), (c), (d)
Section C
Comprehension I
\end{problembox}

\begin{problembox}
\textbf{Problem 330} \hfill \textbf{[Neeraj Kumar [NRJ] Exercise II Q3]}
\label{q:ext-nrj_all_questions-330}

(a)
Comprehension II
\end{problembox}

\begin{problembox}
\textbf{Problem 331} \hfill \textbf{[Neeraj Kumar [NRJ] Exercise II Q6]}
\label{q:ext-nrj_all_questions-331}

(a)
Comprehension III
\end{problembox}

\begin{problembox}
\textbf{Problem 332} \hfill \textbf{[Neeraj Kumar [NRJ] Exercise II Q10]}
\label{q:ext-nrj_all_questions-332}

(b)
Comprehension IV
\end{problembox}

\begin{problembox}
\textbf{Problem 333} \hfill \textbf{[Neeraj Kumar [NRJ] Exercise II Q13]}
\label{q:ext-nrj_all_questions-333}

(c)
Comprehension V
\end{problembox}

\begin{problembox}
\textbf{Problem 334} \hfill \textbf{[Neeraj Kumar [NRJ] Exercise II Q18]}
\label{q:ext-nrj_all_questions-334}

(a)
Comprehension VI
\end{problembox}

\begin{problembox}
\textbf{Problem 335} \hfill \textbf{[Neeraj Kumar [NRJ] Exercise II Q23]}
\label{q:ext-nrj_all_questions-335}

(d)
Comprehension VII
\end{problembox}

\begin{problembox}
\textbf{Problem 336} \hfill \textbf{[Neeraj Kumar [NRJ] Exercise II Q27]}
\label{q:ext-nrj_all_questions-336}

(a)
Comprehension VIII
\end{problembox}

\begin{problembox}
\textbf{Problem 337} \hfill \textbf{[Neeraj Kumar [NRJ] Exercise II Q30]}
\label{q:ext-nrj_all_questions-337}

(d)
Comprehension IX
\end{problembox}

\begin{problembox}
\textbf{Problem 338} \hfill \textbf{[Neeraj Kumar [NRJ] Exercise II Q34]}
\label{q:ext-nrj_all_questions-338}

(c)
Comprehension X
\end{problembox}

\begin{problembox}
\textbf{Problem 339} \hfill \textbf{[Neeraj Kumar [NRJ] Exercise II Q38]}
\label{q:ext-nrj_all_questions-339}

(c)
Comprehension XI
\end{problembox}

\begin{problembox}
\textbf{Problem 340} \hfill \textbf{[Neeraj Kumar [NRJ] Exercise II Q41]}
\label{q:ext-nrj_all_questions-340}

(a)
Comprehension XII
\end{problembox}

\begin{problembox}
\textbf{Problem 341} \hfill \textbf{[Neeraj Kumar [NRJ] Exercise II Q43]}
\label{q:ext-nrj_all_questions-341}

(d)
Comprehension XIII
\end{problembox}

\begin{problembox}
\textbf{Problem 342} \hfill \textbf{[Neeraj Kumar [NRJ] Exercise II Q45]}
\label{q:ext-nrj_all_questions-342}

(c)
Comprehension XIV
\end{problembox}

\begin{problembox}
\textbf{Problem 343} \hfill \textbf{[Neeraj Kumar [NRJ] Exercise II Q47]}
\label{q:ext-nrj_all_questions-343}

(b)
Comprehension XV
\end{problembox}

\begin{problembox}
\textbf{Problem 344} \hfill \textbf{[Neeraj Kumar [NRJ] Exercise II Q53]}
\label{q:ext-nrj_all_questions-344}

(b)
ANSWER KEYS
 
 
Visit www.jeebooks.in for more
\end{problembox}

\begin{problembox}
\textbf{Problem 345} \hfill \textbf{[Neeraj Kumar [NRJ] Exercise II Q15]}
\label{q:ext-nrj_all_questions-345}

(c)
Section E (Column Match)
\end{problembox}

\begin{problembox}
\textbf{Problem 346} \hfill \textbf{[Neeraj Kumar [NRJ] Exercise II Q1]}
\label{q:ext-nrj_all_questions-346}

A $/to$ P, Q; B $/to$ P, Q; C $/to$ Q, R; D $/to$ P, S
\end{problembox}

\begin{problembox}
\textbf{Problem 347} \hfill \textbf{[Neeraj Kumar [NRJ] Exercise II Q2]}
\label{q:ext-nrj_all_questions-347}

A $/to$ R; B $/to$ S; C $/to$ Q; D $/to$ P
\end{problembox}

\begin{problembox}
\textbf{Problem 348} \hfill \textbf{[Neeraj Kumar [NRJ] Exercise II Q3]}
\label{q:ext-nrj_all_questions-348}

A $/to$ R; B $/to$ S; C $/to$ P; D $/to$ Q
\end{problembox}

\begin{problembox}
\textbf{Problem 349} \hfill \textbf{[Neeraj Kumar [NRJ] Exercise II Q4]}
\label{q:ext-nrj_all_questions-349}

A $/to$ R, S; B $/to$ P, R; C $/to$ Q, S
Section F (Subjective)
Single-digit Integer Type
\end{problembox}

\begin{problembox}
\textbf{Problem 350} \hfill \textbf{[Neeraj Kumar [NRJ] Exercise II Q15]}
\label{q:ext-nrj_all_questions-350}

(9)
Four-digit Integer Type
\end{problembox}

\begin{problembox}
\textbf{Problem 351} \hfill \textbf{[Neeraj Kumar [NRJ] Exercise II Q15]}
\label{q:ext-nrj_all_questions-351}

(0108)
ANSWER KEYS
 
 
Visit www.jeebooks.in for more
\end{problembox}

\begin{problembox}
\textbf{Problem 352} \hfill \textbf{[Neeraj Kumar [NRJ] EXERCISE I (JEE MAIN) Q1]}
\label{q:ext-nrj_all_questions-352}

Potential of electrode is independent from the 
area of electrode immersed in solution.
\end{problembox}

\begin{problembox}
\textbf{Problem 353} \hfill \textbf{[Neeraj Kumar [NRJ] EXERCISE I (JEE MAIN) Q3]}
\label{q:ext-nrj_all_questions-353}

Copper cannot reduce Al3+.
\end{problembox}

\begin{problembox}
\textbf{Problem 354} \hfill \textbf{[Neeraj Kumar [NRJ] EXERCISE I (JEE MAIN) Q4]}
\label{q:ext-nrj_all_questions-354}

Cu will reduce Ag+ and itself becomes Cu2+.
\end{problembox}

\begin{problembox}
\textbf{Problem 355} \hfill \textbf{[Neeraj Kumar [NRJ] EXERCISE I (JEE MAIN) Q5]}
\label{q:ext-nrj_all_questions-355}

K is more active metal than Al.
\end{problembox}

\begin{problembox}
\textbf{Problem 356} \hfill \textbf{[Neeraj Kumar [NRJ] EXERCISE I (JEE MAIN) Q6]}
\label{q:ext-nrj_all_questions-356}

Mg is better reducing agent.
\end{problembox}

\begin{problembox}
\textbf{Problem 357} \hfill \textbf{[Neeraj Kumar [NRJ] EXERCISE I (JEE MAIN) Q8]}
\label{q:ext-nrj_all_questions-357}

Lower the reduction potential, strong is the 
reducing agent.
\end{problembox}

\begin{problembox}
\textbf{Problem 358} \hfill \textbf{[Neeraj Kumar [NRJ] EXERCISE I (JEE MAIN) Q9]}
\label{q:ext-nrj_all_questions-358}

Y-- + X 
 Y2-- + X+ is possible.
\end{problembox}

\begin{problembox}
\textbf{Problem 359} \hfill \textbf{[Neeraj Kumar [NRJ] EXERCISE I (JEE MAIN) Q10]}
\label{q:ext-nrj_all_questions-359}

Mno4
- Will also oxidize HCl.
\end{problembox}

\begin{problembox}
\textbf{Problem 360} \hfill \textbf{[Neeraj Kumar [NRJ] EXERCISE I (JEE MAIN) Q11]}
\label{q:ext-nrj_all_questions-360}

Electrochemical series
\end{problembox}

\begin{problembox}
\textbf{Problem 361} \hfill \textbf{[Neeraj Kumar [NRJ] EXERCISE I (JEE MAIN) Q12]}
\label{q:ext-nrj_all_questions-361}

M
M(s);E
ve
n+
-
+
 $/to$
 
= -
= +
(aq)
ne
,
G
ve
$/Delta$
 
 Hence, reverse reaction is spontaneous.
\end{problembox}

\begin{problembox}
\textbf{Problem 362} \hfill \textbf{[Neeraj Kumar [NRJ] EXERCISE I (JEE MAIN) Q13]}
\label{q:ext-nrj_all_questions-362}

Hg Cl
s
aq
Cl
aq
2
2
2
2
2
( )
Hg
(
)
(
)

  


 



+
-
+
 
 
Hg
aq
e
Hg
2
2
2
2
+
-
+
(
)
(l)

  


 
\end{problembox}

\begin{problembox}
\textbf{Problem 363} \hfill \textbf{[Neeraj Kumar [NRJ] EXERCISE I (JEE MAIN) Q14]}
\label{q:ext-nrj_all_questions-363}

E
E
Mno
n
Br
Br
$/circ$
>
$/circ$
-
+
-
4
2
2
/M
/
\end{problembox}

\begin{problembox}
\textbf{Problem 364} \hfill \textbf{[Neeraj Kumar [NRJ] EXERCISE I (JEE MAIN) Q15]}
\label{q:ext-nrj_all_questions-364}

Electrolyte used must have nearly same ionic 
mobility.
\end{problembox}

\begin{problembox}
\textbf{Problem 365} \hfill \textbf{[Neeraj Kumar [NRJ] EXERCISE I (JEE MAIN) Q16]}
\label{q:ext-nrj_all_questions-365}

MnO
H
e
Mn
H O
2
4
2
8
5
4
-
+
-
+
+
+
+

  


 



 
 
$/Delta$E
H
= -
 
$/times$
= -
 
+
0 059
5
1
1
0 059
5
1
0 01
8
8
.
log
[
]
.
log
( .
)
 
 = -- 0.1888 V
\end{problembox}

\begin{problembox}
\textbf{Problem 366} \hfill \textbf{[Neeraj Kumar [NRJ] EXERCISE I (JEE MAIN) Q17]}
\label{q:ext-nrj_all_questions-366}

Cu
e
Cu
2+ +
-
2

  


 



 
 
E
E
E
=
$/circ$ -
 
=
$/circ$ -
+
0 059
2
1
0 059
2
1
0 1
2
.
log
[Cu
]
.
log .
 
      = Eo -- 0.0295 V
\end{problembox}

\begin{problembox}
\textbf{Problem 367} \hfill \textbf{[Neeraj Kumar [NRJ] EXERCISE I (JEE MAIN) Q18]}
\label{q:ext-nrj_all_questions-367}

O
H O(l)
2
2
(g)
H (aq)
e
+
+
+
-
4
4
2

  


 



 
 
E
E
V
Basic
acidic
$/circ$
=
$/circ$
-
 
=
-
0 059
4
1
10
0 404
14 4
.
log
(
)
.
\end{problembox}

\begin{problembox}
\textbf{Problem 368} \hfill \textbf{[Neeraj Kumar [NRJ] EXERCISE I (JEE MAIN) Q19]}
\label{q:ext-nrj_all_questions-368}

E
E
E
V
Cu /Cu
Cu
/cu
Cu
Cu
2+
2+
+
+
=
$/times$
-$/times$
-
=
2
1
2
1
0 521
/
.
\end{problembox}

\begin{problembox}
\textbf{Problem 369} \hfill \textbf{[Neeraj Kumar [NRJ] EXERCISE I (JEE MAIN) Q20]}
\label{q:ext-nrj_all_questions-369}

2
2
2
H
H
+(aq)
e
(g)
+
-
  


 



 
 
E = E$/circ$ -
 
=
-
 
+
-
0 059
2
1
0
0 059
2
1
10
2
7 2
.
log
(H )
.
log
(
)
 
 = -- 0.0413
\end{problembox}

\begin{problembox}
\textbf{Problem 370} \hfill \textbf{[Neeraj Kumar [NRJ] EXERCISE I (JEE MAIN) Q22]}
\label{q:ext-nrj_all_questions-370}

Ionic mobility of ions in salt bridge should be 
same.
\end{problembox}

\begin{problembox}
\textbf{Problem 371} \hfill \textbf{[Neeraj Kumar [NRJ] EXERCISE I (JEE MAIN) Q23]}
\label{q:ext-nrj_all_questions-371}

Higher reduction potential, stronger is the 
oxidizing agent.
\end{problembox}

\begin{problembox}
\textbf{Problem 372} \hfill \textbf{[Neeraj Kumar [NRJ] EXERCISE I (JEE MAIN) Q24]}
\label{q:ext-nrj_all_questions-372}

Lower the reduction potential, stronger is the 
reducing agent.
\end{problembox}

\begin{problembox}
\textbf{Problem 373} \hfill \textbf{[Neeraj Kumar [NRJ] EXERCISE I (JEE MAIN) Q25]}
\label{q:ext-nrj_all_questions-373}

E
E
E
V
Fe
Fe
Fe
Fe
Fe
Fe
$/circ$
=
$/times$
$/circ$
-
$/times$
$/circ$
-
= +
+
+
+
+
3
2
3
2
3
2
3
2
0 772
/
/
/
.
\end{problembox}

\begin{problembox}
\textbf{Problem 374} \hfill \textbf{[Neeraj Kumar [NRJ] EXERCISE I (JEE MAIN) Q26]}
\label{q:ext-nrj_all_questions-374}

H
g
H
e
2
2
2
( )
(aq)

  


 



+
+
-
E
E
E
(H )
P
H /CH COOH
H /H
H /H
2
H
2
3
2
+
2
2
=
=
$/circ$
-
 
+
+
0 06
2
.
log
 
  
              
=
-
 
$/times$
(
)
= -
 
$/times$
$/times$
-
0
0 06
2
0 06
2
1 8 10
0 5
1
2
5
2
.
log
.
log .
.
Ka
C
PH
 
                       = + 0.1512 V
 
 
Visit www.jeebooks.in for more
\end{problembox}

\begin{problembox}
\textbf{Problem 375} \hfill \textbf{[Neeraj Kumar [NRJ] EXERCISE I (JEE MAIN) Q27]}
\label{q:ext-nrj_all_questions-375}

E
E
K
Cl /CuCl/Cu
Cu
Cu
sp
+
$/circ$
=
$/circ$
-
-
/
.
log
0 06
1
1
 
 or, 1 28
0 06
1
1
2
10 7
.
.
log
/
=
$/circ$
-
 
$/times$
-
E Cu
Cu
+
 
 
$/implies$ E
V
Cu
Cu
+
$/circ$
=
/
.1 682
\end{problembox}

\begin{problembox}
\textbf{Problem 376} \hfill \textbf{[Neeraj Kumar [NRJ] EXERCISE I (JEE MAIN) Q28]}
\label{q:ext-nrj_all_questions-376}

E
E
K
OH
Pb(OH)
Pb
Pb
Pb
sp
P
$/circ$
=
$/circ$
-
 
-/
/
/
.
log
2
2
0 06
2
1
 
 or, (--0.55) = (--0.13) -- 0 06
2
1
.
log
 
Ksp
 
   Ksp = 1.0 $/times$ 10--14
\end{problembox}

\begin{problembox}
\textbf{Problem 377} \hfill \textbf{[Neeraj Kumar [NRJ] EXERCISE I (JEE MAIN) Q29]}
\label{q:ext-nrj_all_questions-377}

E
E
E
MnO
MnO
MnO
Mn
MnO
Mn
2
4
4
2
2
2
5
2
5
2
-
-
+
+
=
$/times$
-
$/times$
-
/
/
/
 
                      =
$/times$
-
$/times$
=
5 1 51
2
1 23
3
1 697
.
.
.
V
\end{problembox}

\begin{problembox}
\textbf{Problem 378} \hfill \textbf{[Neeraj Kumar [NRJ] EXERCISE I (JEE MAIN) Q30]}
\label{q:ext-nrj_all_questions-378}

For stronger oxidizing agent, reduction potential 
should be high.
Galvanic Cell
\end{problembox}

\begin{problembox}
\textbf{Problem 379} \hfill \textbf{[Neeraj Kumar [NRJ] EXERCISE I (JEE MAIN) Q31]}
\label{q:ext-nrj_all_questions-379}

Reduction occur at right electrode.
\end{problembox}

\begin{problembox}
\textbf{Problem 380} \hfill \textbf{[Neeraj Kumar [NRJ] EXERCISE I (JEE MAIN) Q32]}
\label{q:ext-nrj_all_questions-380}

Circuit becomes incomplete.
\end{problembox}

\begin{problembox}
\textbf{Problem 381} \hfill \textbf{[Neeraj Kumar [NRJ] EXERCISE I (JEE MAIN) Q33]}
\label{q:ext-nrj_all_questions-381}

ECell = (ERP)R -- (ERP)L = 0
\end{problembox}

\begin{problembox}
\textbf{Problem 382} \hfill \textbf{[Neeraj Kumar [NRJ] EXERCISE I (JEE MAIN) Q34]}
\label{q:ext-nrj_all_questions-382}

Net cell reaction is Cl2 (P2 atm) 
 Cl2 (P1 
atm) for spontaneous reaction: P2 > P1
\end{problembox}

\begin{problembox}
\textbf{Problem 383} \hfill \textbf{[Neeraj Kumar [NRJ] EXERCISE I (JEE MAIN) Q38]}
\label{q:ext-nrj_all_questions-383}

E
E
RT
F
Zn
Cu
cell
cell
2+
=
$/circ$
-
 
+
2
2
ln [
]
[
]
\end{problembox}

\begin{problembox}
\textbf{Problem 384} \hfill \textbf{[Neeraj Kumar [NRJ] EXERCISE I (JEE MAIN) Q39]}
\label{q:ext-nrj_all_questions-384}

E cell
$/circ$
 is independent from size or amount of 
electrode as well as PH of solution.
\end{problembox}

\begin{problembox}
\textbf{Problem 385} \hfill \textbf{[Neeraj Kumar [NRJ] EXERCISE I (JEE MAIN) Q42]}
\label{q:ext-nrj_all_questions-385}

Ca and Cu are far away in the electrochemical 
series.
\end{problembox}

\begin{problembox}
\textbf{Problem 386} \hfill \textbf{[Neeraj Kumar [NRJ] EXERCISE I (JEE MAIN) Q43]}
\label{q:ext-nrj_all_questions-386}

Reaction during charging of cell is 2Pb SO4 (s) + 
2H2O (l) 
 Pb(s) + PbO2 (s) + 2H2 SO4 (aq)
\end{problembox}

\begin{problembox}
\textbf{Problem 387} \hfill \textbf{[Neeraj Kumar [NRJ] EXERCISE I (JEE MAIN) Q45]}
\label{q:ext-nrj_all_questions-387}

Depolarizer is the term used for oxidizing agent in 
the cell.
\end{problembox}

\begin{problembox}
\textbf{Problem 388} \hfill \textbf{[Neeraj Kumar [NRJ] EXERCISE I (JEE MAIN) Q47]}
\label{q:ext-nrj_all_questions-388}

Pb(s) + PbO2 (s) + 2H2SO4 (aq) 
 2PbSO4 
(aq) + 2H2O(l)
\end{problembox}

\begin{problembox}
\textbf{Problem 389} \hfill \textbf{[Neeraj Kumar [NRJ] EXERCISE I (JEE MAIN) Q48]}
\label{q:ext-nrj_all_questions-389}

E
n
K
cell
eq
$/circ$
=
 
0 059
.
log
 $/implies$ 0 295
0 059
2
.
.
log
=
 
Keq
 
   Keq = 1010
\end{problembox}

\begin{problembox}
\textbf{Problem 390} \hfill \textbf{[Neeraj Kumar [NRJ] EXERCISE I (JEE MAIN) Q49]}
\label{q:ext-nrj_all_questions-390}

L E : Ag (s)
Ag
c
K
M
e
+
SP
 
=
(
)+
-

  


 



1
 
 
R E : Ag
(C
0.1M)
e
Ag(s)
2
 
=
+
+
-
  


 



 
   Net reaction: Ag
C
Ag
+
+
(
)
(C )
2
1

  


 



 
 
E
C
C
cell = -
 
0 059
1
1
2
.
log
 
 
 $/implies$ 0 413
0 059
1
0 1
.
.
log
.
= -
Ksp
 
   Ksp = 1.0 $/times$ 10--16
\end{problembox}

\begin{problembox}
\textbf{Problem 391} \hfill \textbf{[Neeraj Kumar [NRJ] EXERCISE I (JEE MAIN) Q50]}
\label{q:ext-nrj_all_questions-391}

Net cell reaction = H2 (P1) 
 H2 (P2)
 
 
E
V
cell = -
 
= +
0 06
2
420
600
0 0045
.
log
.
\end{problembox}

\begin{problembox}
\textbf{Problem 392} \hfill \textbf{[Neeraj Kumar [NRJ] EXERCISE I (JEE MAIN) Q51]}
\label{q:ext-nrj_all_questions-392}

E cell
$/circ$
=
-
=
0 44
0 33
0 11
.
.
.
V
\end{problembox}

\begin{problembox}
\textbf{Problem 393} \hfill \textbf{[Neeraj Kumar [NRJ] EXERCISE I (JEE MAIN) Q52]}
\label{q:ext-nrj_all_questions-393}

Cell reaction: Zn
Fe
Zn
Fe
2+
+
+
+
2 
  


 



 
 
E
E
Zn
[Fe
cell
cell
=
$/circ$
-
 
[
]
+
+
0 059
2
2
2
.
log
]
 
 or, 0 2905
0 059
2
0 059
2
0 01
0 001
.
.
log
.
log
.
.
=
 
-
 
Keq
 
  Keq =
 
  
 
  
10
0 32
0 0295
.
.
 
 
Visit www.jeebooks.in for more
\end{problembox}

\begin{problembox}
\textbf{Problem 394} \hfill \textbf{[Neeraj Kumar [NRJ] EXERCISE I (JEE MAIN) Q53]}
\label{q:ext-nrj_all_questions-394}

E
K
cell
eq
$/circ$
=
 
0 06
.
log
n
 $/implies$ 0 18
0 06
4
.
.
log
=
 
Keq
 
   log Kc = 12
\end{problembox}

\begin{problembox}
\textbf{Problem 395} \hfill \textbf{[Neeraj Kumar [NRJ] EXERCISE I (JEE MAIN) Q54]}
\label{q:ext-nrj_all_questions-395}

E
E
RT
F
Zn
Cu
cell
cell
=
$/circ$
-
 
+
+
2
2
2
ln [
]
[
]
\end{problembox}

\begin{problembox}
\textbf{Problem 396} \hfill \textbf{[Neeraj Kumar [NRJ] EXERCISE I (JEE MAIN) Q55]}
\label{q:ext-nrj_all_questions-396}

E
E
RT
F
Ni
Cu
cell
cell
=
$/circ$
-
 
+
+
2
2
2
ln [
]
[
]
\end{problembox}

\begin{problembox}
\textbf{Problem 397} \hfill \textbf{[Neeraj Kumar [NRJ] EXERCISE I (JEE MAIN) Q56]}
\label{q:ext-nrj_all_questions-397}

E
E
0.059
E
2
1
=
-
 
=
-
1
2
1 0
0 01
0 059
log
.
.
.
\end{problembox}

\begin{problembox}
\textbf{Problem 398} \hfill \textbf{[Neeraj Kumar [NRJ] EXERCISE I (JEE MAIN) Q57]}
\label{q:ext-nrj_all_questions-398}

E
0.059
V
cell
$/circ$
=
 
=
2
10
0 354
12
log
.
\end{problembox}

\begin{problembox}
\textbf{Problem 399} \hfill \textbf{[Neeraj Kumar [NRJ] EXERCISE I (JEE MAIN) Q58]}
\label{q:ext-nrj_all_questions-399}

[
]
.
HA
H
K
C
M
a
+
-
-
=
 
=
$/times$
=
10
0 1
10
5
3
\end{problembox}

\begin{problembox}
\textbf{Problem 400} \hfill \textbf{[Neeraj Kumar [NRJ] EXERCISE I (JEE MAIN) Q60]}
\label{q:ext-nrj_all_questions-400}

E
n
K
cell
eq
$/circ$
=
 
0 06
.
log
 
 
 $/implies$ ( .
.
)
.
log
0 77
0 53
0 06
2
-
=
 
Keq
 
   Keq = 108
\end{problembox}

\begin{problembox}
\textbf{Problem 401} \hfill \textbf{[Neeraj Kumar [NRJ] EXERCISE I (JEE MAIN) Q61]}
\label{q:ext-nrj_all_questions-401}

Current   ow from right electrode to left electrode 
in external circuit.
\end{problembox}

\begin{problembox}
\textbf{Problem 402} \hfill \textbf{[Neeraj Kumar [NRJ] EXERCISE I (JEE MAIN) Q62]}
\label{q:ext-nrj_all_questions-402}

$/Delta$
$/circ$ = -
 $/circ$
G
nF E cell $/implies$ --48250 = --2 $/times$ 96500 $/times$ E cell
$/circ$
 
   E cell
$/circ$
 = 0.25 V
 
 Now, E
E
Cl
P
cell
cell
H 2
=
$/circ$
-
 
+
-
0 059
2
2
2
.
log [H ] [
]
 
 
=
-
 
$/times$
=
0 25
0 059
2
0 1
0 1
1
0 368
2
2
.
.
log ( . )
( . )
.
V
\end{problembox}

\begin{problembox}
\textbf{Problem 403} \hfill \textbf{[Neeraj Kumar [NRJ] EXERCISE I (JEE MAIN) Q63]}
\label{q:ext-nrj_all_questions-403}

Net cell reaction: H+(C = ?) 
 H+ (10--6 M)
 
 
E
c
cell = -
 
=
-
0 059
1
10
0 118
6
.
log
.
 $/implies$ C = 10--4 M
\end{problembox}

\begin{problembox}
\textbf{Problem 404} \hfill \textbf{[Neeraj Kumar [NRJ] EXERCISE I (JEE MAIN) Q64]}
\label{q:ext-nrj_all_questions-404}

E
E
E
OCl
OCl
Cl
$/circ$
=
$/times$
$/circ$
-$/times$
$/circ$
-
=
$/times$
-$/times$
-
-
-
-
/Cl
/Cl
/Cl
.
.
2
2
2
1
2
1
2
0 94
1 1 36
1
 
                   = 0.52 V
\end{problembox}

\begin{problembox}
\textbf{Problem 405} \hfill \textbf{[Neeraj Kumar [NRJ] EXERCISE I (JEE MAIN) Q65]}
\label{q:ext-nrj_all_questions-405}

For (i)  and  (iii); E
V
cell
$/circ$
= 3 0.
Electrolysis
\end{problembox}

\begin{problembox}
\textbf{Problem 406} \hfill \textbf{[Neeraj Kumar [NRJ] EXERCISE I (JEE MAIN) Q66]}
\label{q:ext-nrj_all_questions-406}

Oxidation occur at anode and reduction at 
cathode, in both.
\end{problembox}

\begin{problembox}
\textbf{Problem 407} \hfill \textbf{[Neeraj Kumar [NRJ] EXERCISE I (JEE MAIN) Q69]}
\label{q:ext-nrj_all_questions-407}

Due to high over-voltage potential, reduction of 
H2O becomes less favorable.
\end{problembox}

\begin{problembox}
\textbf{Problem 408} \hfill \textbf{[Neeraj Kumar [NRJ] EXERCISE I (JEE MAIN) Q70]}
\label{q:ext-nrj_all_questions-408}

Electrochemical equivalent, Z = E
F and E is 
maximum for silver.
\end{problembox}

\begin{problembox}
\textbf{Problem 409} \hfill \textbf{[Neeraj Kumar [NRJ] EXERCISE I (JEE MAIN) Q72]}
\label{q:ext-nrj_all_questions-409}

W = Q
F
A
$/times$ n
 
 For smaller n, w will be higher.
\end{problembox}

\begin{problembox}
\textbf{Problem 410} \hfill \textbf{[Neeraj Kumar [NRJ] EXERCISE I (JEE MAIN) Q73]}
\label{q:ext-nrj_all_questions-410}

Cathode to anode in internal supply.
+
-- e--
e--
\end{problembox}

\begin{problembox}
\textbf{Problem 411} \hfill \textbf{[Neeraj Kumar [NRJ] EXERCISE I (JEE MAIN) Q76]}
\label{q:ext-nrj_all_questions-411}

Equivalent weight is maximum for Na.
\end{problembox}

\begin{problembox}
\textbf{Problem 412} \hfill \textbf{[Neeraj Kumar [NRJ] EXERCISE I (JEE MAIN) Q77]}
\label{q:ext-nrj_all_questions-412}

n
n
eq
eq
H
Cu
2 =
 $/implies$ 0 5
1
63 5
2
.
.
=
$/times$
w
 
 
 $/implies$ w = 15.875 gm
 
 
Visit www.jeebooks.in for more
\end{problembox}

\begin{problembox}
\textbf{Problem 413} \hfill \textbf{[Neeraj Kumar [NRJ] EXERCISE I (JEE MAIN) Q78]}
\label{q:ext-nrj_all_questions-413}

Ni will oxidize and Ni2+ will reduce resulting no 
change in Ni2+ Concentration.
\end{problembox}

\begin{problembox}
\textbf{Problem 414} \hfill \textbf{[Neeraj Kumar [NRJ] EXERCISE I (JEE MAIN) Q79]}
\label{q:ext-nrj_all_questions-414}

W
E
Q
F
=
 $/implies$ 9 8
63 5
2
3 10000
96500
.
. $/times$
=
$/times$
$/times$  $/implies$   = 0.99
\end{problembox}

\begin{problembox}
\textbf{Problem 415} \hfill \textbf{[Neeraj Kumar [NRJ] EXERCISE I (JEE MAIN) Q80]}
\label{q:ext-nrj_all_questions-415}

Cathode = H2, anode = O2 and 
W
W
H
O
2
2
= 1
8
\end{problembox}

\begin{problembox}
\textbf{Problem 416} \hfill \textbf{[Neeraj Kumar [NRJ] EXERCISE I (JEE MAIN) Q81]}
\label{q:ext-nrj_all_questions-416}

W = Z   i   t $/implies$ Z   i = constant
 
   
i
i
1
2
2
2
2
1
=
=
Z
Z
E
E
\end{problembox}

\begin{problembox}
\textbf{Problem 417} \hfill \textbf{[Neeraj Kumar [NRJ] EXERCISE I (JEE MAIN) Q83]}
\label{q:ext-nrj_all_questions-417}

Informative
8 4. At cathode: Cu
aq)
2e
Cu(s)
2+
-
+
 $/to$
 
(
 
 At anode: 2
2
2
H O(l)
O (g)
4H
aq)
4e
 $/to$
 
+
+
+
-
(
 
 
 Hence, solution becomes rich in H2SO4.
\end{problembox}

\begin{problembox}
\textbf{Problem 418} \hfill \textbf{[Neeraj Kumar [NRJ] EXERCISE I (JEE MAIN) Q85]}
\label{q:ext-nrj_all_questions-418}

Reduction occur   rst for the species having higher 
reduction potential.
\end{problembox}

\begin{problembox}
\textbf{Problem 419} \hfill \textbf{[Neeraj Kumar [NRJ] EXERCISE I (JEE MAIN) Q86]}
\label{q:ext-nrj_all_questions-419}

For very dilute NaCl solution,
 
 Cathode: 2H2O(l) + 2e-- $/to$ H2(g) + 2OH-- (aq); PH  
 
 Anode:   2H2O(l) $/to$ O2(g) + 4H+ (aq) + 4e--; PH  
\end{problembox}

\begin{problembox}
\textbf{Problem 420} \hfill \textbf{[Neeraj Kumar [NRJ] EXERCISE I (JEE MAIN) Q87]}
\label{q:ext-nrj_all_questions-420}

w
w
E
E
A
A
1
2
1
2
2
3
3
2
=
=
=
/
/
\end{problembox}

\begin{problembox}
\textbf{Problem 421} \hfill \textbf{[Neeraj Kumar [NRJ] EXERCISE I (JEE MAIN) Q88]}
\label{q:ext-nrj_all_questions-421}

neq of Zn =
$/times$
=
100
65 4
2
3 058
.
.
 
 
 neq of CuSO4 = 1.0 $/times$ 2 = 2.0 (L.R.)
 
 
 neq
Q
F
=
 $/implies$ 2
1 0
96500
=
$/times$
.
t 
 
 $/implies$ t = 193000 sec = 53.6 hr.
\end{problembox}

\begin{problembox}
\textbf{Problem 422} \hfill \textbf{[Neeraj Kumar [NRJ] EXERCISE I (JEE MAIN) Q89]}
\label{q:ext-nrj_all_questions-422}

neq =
$/times$
$/times$
=
-
6
10
6
10
10
20
23
3
\end{problembox}

\begin{problembox}
\textbf{Problem 423} \hfill \textbf{[Neeraj Kumar [NRJ] EXERCISE I (JEE MAIN) Q90]}
\label{q:ext-nrj_all_questions-423}

neq
Q
F
=
 $/implies$ w
208
2
1 100
3600
96500
$/times$
=
$/times$
$/times$
 
 
 $/implies$ w $/approx$ 388 gm
\end{problembox}

\begin{problembox}
\textbf{Problem 424} \hfill \textbf{[Neeraj Kumar [NRJ] EXERCISE I (JEE MAIN) Q91]}
\label{q:ext-nrj_all_questions-424}

neq Ag deposited = neqCu dissolved
 
 or, 3 24
108
1
64
2
.
$/times$ =
$/times$
w
 $/implies$ w = 0.96 gm
 
   Final mass of Cu anode = 60 -- 0.96 = 59.04 gm
\end{problembox}

\begin{problembox}
\textbf{Problem 425} \hfill \textbf{[Neeraj Kumar [NRJ] EXERCISE I (JEE MAIN) Q92]}
\label{q:ext-nrj_all_questions-425}

neq Fe2O3 reduced = n $/times$ n-factor = 1 $/times$ 6
 
   Charge needed = 6 F
\end{problembox}

\begin{problembox}
\textbf{Problem 426} \hfill \textbf{[Neeraj Kumar [NRJ] EXERCISE I (JEE MAIN) Q93]}
\label{q:ext-nrj_all_questions-426}

Number of moles, n = neq/n-factor
 
   nAl : nCu : nNa = 3
3
3
2
3
1
2 3 6
:
:
:
:
=
\end{problembox}

\begin{problembox}
\textbf{Problem 427} \hfill \textbf{[Neeraj Kumar [NRJ] EXERCISE I (JEE MAIN) Q94]}
\label{q:ext-nrj_all_questions-427}

2
2
2
4
2
8
2
HSO
S O
H
e
-
-
+
-
+
+

  


 



 
 Now, neq
Q
F
=
 $/implies$ 1
2
3600
0 75
96500
$/times$
=
$/times$
$/times$
i
.
 
 
   i = 71.48 A
\end{problembox}

\begin{problembox}
\textbf{Problem 428} \hfill \textbf{[Neeraj Kumar [NRJ] EXERCISE I (JEE MAIN) Q95]}
\label{q:ext-nrj_all_questions-428}

Number of Faradays required = neq = 1 $/times$ 2 = 2
\end{problembox}

\begin{problembox}
\textbf{Problem 429} \hfill \textbf{[Neeraj Kumar [NRJ] EXERCISE I (JEE MAIN) Q96]}
\label{q:ext-nrj_all_questions-429}

neq = n $/times$ n-factor = Q
F
 
 or, 
1
0 0821 546
4
96500
96500
$/times$
$/times$
$/times$
=
V
.
 $/implies$ V = 5.6 L
\end{problembox}

\begin{problembox}
\textbf{Problem 430} \hfill \textbf{[Neeraj Kumar [NRJ] EXERCISE I (JEE MAIN) Q97]}
\label{q:ext-nrj_all_questions-430}

neqAg = neqO2 $/implies$ w
108
1
1 6
32
4
$/times$ =
$/times$
.
 $/implies$ w = 21.6 gm
\end{problembox}

\begin{problembox}
\textbf{Problem 431} \hfill \textbf{[Neeraj Kumar [NRJ] EXERCISE I (JEE MAIN) Q98]}
\label{q:ext-nrj_all_questions-431}

neq
Q
F
=
 $/implies$ 
3
96 5
12000
96500
. $/times$
=
n
 $/implies$ n = 4
\end{problembox}

\begin{problembox}
\textbf{Problem 432} \hfill \textbf{[Neeraj Kumar [NRJ] EXERCISE I (JEE MAIN) Q99]}
\label{q:ext-nrj_all_questions-432}

Number  of electrons 
 
  
 
= Q
e =
$/times$
$/times$
=
$/times$
-
5
200
1 602
10
6 24
10
19
21
.
.
\end{problembox}

\begin{problembox}
\textbf{Problem 433} \hfill \textbf{[Neeraj Kumar [NRJ] EXERCISE I (JEE MAIN) Q100]}
\label{q:ext-nrj_all_questions-433}

neq
Q
F
=
 $/implies$ 3 0
9 65 10
60
96500
.
.
E =
$/times$
$/times$
 $/implies$ E = 50
\end{problembox}

\begin{problembox}
\textbf{Problem 434} \hfill \textbf{[Neeraj Kumar [NRJ] EXERCISE I (JEE MAIN) Q101]}
\label{q:ext-nrj_all_questions-434}

Cu will deposit at cathode.
 
 
neq
Q
F
=
$/implies$w
63 5
2
1 96 5
60
96500
.
.
$/times$
=
$/times$
$/times$
$/implies$ w = 1.905 gm
\end{problembox}

\begin{problembox}
\textbf{Problem 435} \hfill \textbf{[Neeraj Kumar [NRJ] EXERCISE I (JEE MAIN) Q102]}
\label{q:ext-nrj_all_questions-435}

n
n
n
eq
-factor
Q
F
=
$/times$
=
 
 or, 
2 8
22400
4
96500
.
$/times$
=
$/times$
i
1 $/implies$ i = 48.25 A/sec
\end{problembox}

\begin{problembox}
\textbf{Problem 436} \hfill \textbf{[Neeraj Kumar [NRJ] EXERCISE I (JEE MAIN) Q103]}
\label{q:ext-nrj_all_questions-436}

Mass of Na deposited 
 
  
 
=
-
$/times$
=
18 7
100
18 7
10
2 623
.
(
. )
.
gm
 
 Now, neq
Q
F
=
 $/implies$ 2 623
23
1
0 5
.
.
$/times$ =
$/times$
Q
F
 $/implies$ Q $/approx$ 0.2 F
\end{problembox}

\begin{problembox}
\textbf{Problem 437} \hfill \textbf{[Neeraj Kumar [NRJ] EXERCISE I (JEE MAIN) Q104]}
\label{q:ext-nrj_all_questions-437}

neq A = neq B
 
 or, 5 6
112
0 9
27
3
.
.
$/times$
=
$/times$
VA
 
   Valency of A, VA = 2
\end{problembox}

\begin{problembox}
\textbf{Problem 438} \hfill \textbf{[Neeraj Kumar [NRJ] EXERCISE I (JEE MAIN) Q105]}
\label{q:ext-nrj_all_questions-438}

neq Zn = neq Ni $/implies$ 22 89
32 7
20 55
.
.
.
= ENi
 
 
   ENi = 29.36
 
 
Visit www.jeebooks.in for more
\end{problembox}

\begin{problembox}
\textbf{Problem 439} \hfill \textbf{[Neeraj Kumar [NRJ] EXERCISE I (JEE MAIN) Q106]}
\label{q:ext-nrj_all_questions-439}

On dilution, conductance decreases and hence 
resistance increase.
\end{problembox}

\begin{problembox}
\textbf{Problem 440} \hfill \textbf{[Neeraj Kumar [NRJ] EXERCISE I (JEE MAIN) Q108]}
\label{q:ext-nrj_all_questions-440}

HCl is molecular.
\end{problembox}

\begin{problembox}
\textbf{Problem 441} \hfill \textbf{[Neeraj Kumar [NRJ] EXERCISE I (JEE MAIN) Q109]}
\label{q:ext-nrj_all_questions-441}

Radius of aq.ion : K+ < Na+ < Li+
\end{problembox}

\begin{problembox}
\textbf{Problem 442} \hfill \textbf{[Neeraj Kumar [NRJ] EXERCISE I (JEE MAIN) Q110]}
\label{q:ext-nrj_all_questions-442}

On dilution,  m increases.
\end{problembox}

\begin{problembox}
\textbf{Problem 443} \hfill \textbf{[Neeraj Kumar [NRJ] EXERCISE I (JEE MAIN) Q111]}
\label{q:ext-nrj_all_questions-443}

H2SO4 is strong electrolyte
\end{problembox}

\begin{problembox}
\textbf{Problem 444} \hfill \textbf{[Neeraj Kumar [NRJ] EXERCISE I (JEE MAIN) Q112]}
\label{q:ext-nrj_all_questions-444}

For AgCl,  $/circ$
=  
=
$/circ$
+ $/circ$
+
-
m
m
m(Ag
m Cl
$/lambda$
$/lambda$
)
(
)
 
 = 67.9 + 82.1 = 150 $/Omega$--1 cm2 mol--1
 
 Now,  
=
m
C
$/kappa$ $/implies$ 150
1 8 10 6
=
$/times$
-
.
s
 
   S = 1.2 $/times$ 10--8 mol cm--3 = 1.2 $/times$ 10--5 mol e--1
\end{problembox}

\begin{problembox}
\textbf{Problem 445} \hfill \textbf{[Neeraj Kumar [NRJ] EXERCISE I (JEE MAIN) Q113]}
\label{q:ext-nrj_all_questions-445}

 $/circ$
=  $/circ$
+  $/circ$
- $/circ$
eq(NH OH)
eq(NH Cl)
eq(NaOH)
eq(NaCl)
4
4
 
 
 = 129.8 + 217.4 -- 108.9 = 238.3 $/Omega$--1 cm2 mol--1
 
 Now, $/alpha$ =
=
=
=
$/circ$
$/lambda$
$/lambda$
eq
eq
9 532
238 3
0 04
4
.
.
.
/%
\end{problembox}

\begin{problembox}
\textbf{Problem 446} \hfill \textbf{[Neeraj Kumar [NRJ] EXERCISE I (JEE MAIN) Q114]}
\label{q:ext-nrj_all_questions-446}

 
=
=
$/times$
$/times$
=
$/times$
-
m
C
$/kappa$
1
250 125
1 10
5 10
3
4 $/Omega$--1 m2 mol--1
\end{problembox}

\begin{problembox}
\textbf{Problem 447} \hfill \textbf{[Neeraj Kumar [NRJ] EXERCISE I (JEE MAIN) Q115]}
\label{q:ext-nrj_all_questions-447}

MgSO4(aq) + Ba(OH)2 (added) 
 
 
Mg(OH)2(s) + BaSO4(s)
 
 First number of ions decreases and then increases.
\end{problembox}

\begin{problembox}
\textbf{Problem 448} \hfill \textbf{[Neeraj Kumar [NRJ] EXERCISE I (JEE MAIN) Q116]}
\label{q:ext-nrj_all_questions-448}

 
=
=
$/times$
$/times$
$/times$
=
$/times$
-
-
eq
C
$/kappa$
2 6 10
1 2
10
1 3 10
1
3
2
.
(
)
.
 $/Omega$--1 cm2 eq--1
\end{problembox}

\begin{problembox}
\textbf{Problem 449} \hfill \textbf{[Neeraj Kumar [NRJ] EXERCISE I (JEE MAIN) Q117]}
\label{q:ext-nrj_all_questions-449}

$/alpha$ =  
 $/circ$
=
+
=
m
m
or 
16 30
349 8
40 9
0 04172
4 172
.
(
.
. )
.
.
/%
\end{problembox}

\begin{problembox}
\textbf{Problem 450} \hfill \textbf{[Neeraj Kumar [NRJ] EXERCISE I (JEE MAIN) Q118]}
\label{q:ext-nrj_all_questions-450}

Cell constant, C
A
cm
m
=
=
=
=
-
-
l
2 5
5
0 5
50
1
1
.
.
\end{problembox}

\begin{problembox}
\textbf{Problem 451} \hfill \textbf{[Neeraj Kumar [NRJ] EXERCISE I (JEE MAIN) Q119]}
\label{q:ext-nrj_all_questions-451}

 $/circ$
=  $/circ$
- $/circ$
+  $/circ$
-
-
m(
m(
m(Cl
m(
NH OH)
NH Cl)
)
OH )
4
4
 
 = 150 -- 75 + 200 = 275 $/Omega$--1 cm2 mol--1
 
 Now, $/alpha$ =  
 $/circ$
=
=
m
m
22
275
0 08
.
\end{problembox}

\begin{problembox}
\textbf{Problem 452} \hfill \textbf{[Neeraj Kumar [NRJ] EXERCISE I (JEE MAIN) Q120]}
\label{q:ext-nrj_all_questions-452}

 $/circ$
=
$/circ$
+
$/circ$
=
+
+
-
m(
m(
m(OH
.
.
H O)
H )
)
2
0 035
0 020
$/lambda$
$/lambda$
 
 = 0.055 mho m2 mol--1
 
 Now,  
=
m
C
$/kappa$ $/implies$ 0 055
5 5 10 6
.
.
=
$/times$
-
C
 
 
   C = 10--4 mol m--3 = 10--7 M $/implies$ Kw = C2 = 10--14 M2
\end{problembox}

\begin{problembox}
\textbf{Problem 453} \hfill \textbf{[Neeraj Kumar [NRJ] EXERCISE I (JEE MAIN) Q121]}
\label{q:ext-nrj_all_questions-453}

$/alpha$ =  
 $/circ$
=
$/times$
$/times$
=
-
-
m
m
7 814
10
3 907
10
0 02
4
2
.
.
.
 
 Now, K
C
C
a =
 
-
 
=
$/times$
=
$/times$
-
$/alpha$
$/alpha$
$/alpha$
2
2
2
5
1
0 02
0 05
2
10

( .
)
.
\end{problembox}

\begin{problembox}
\textbf{Problem 454} \hfill \textbf{[Neeraj Kumar [NRJ] EXERCISE I (JEE MAIN) Q122]}
\label{q:ext-nrj_all_questions-454}

n
C
m =
=
$/times$
$/times$
=
-
-
$/kappa$
6 3 10
0 015
2
10
84
4
3
.
.
 S cm2 mol--1
\end{problembox}

\begin{problembox}
\textbf{Problem 455} \hfill \textbf{[Neeraj Kumar [NRJ] EXERCISE I (JEE MAIN) Q123]}
\label{q:ext-nrj_all_questions-455}

 
=  $/circ$
-
 
m
m
A
C 
(for strong electrolyte)
 
 or, 260
0 25
=  $/circ$
-
$/times$
m
A
.
 
(1)
 
 
250
1 00
=  $/circ$
-
$/times$
m
A
.
 
(2)
 
    $/circ$
=
m
270 $/Omega$--1 cm2 mol--1
\end{problembox}

\begin{problembox}
\textbf{Problem 456} \hfill \textbf{[Neeraj Kumar [NRJ] EXERCISE I (JEE MAIN) Q124]}
\label{q:ext-nrj_all_questions-456}

$/kappa$
$/kappa$
A
B
G and 
G
=
$/times$
=
$/times$
 
 
1
50
1
100
 
 Now, G
G
$/Omega$
mix
A
B
=
+
 
  
 
  =
 
-
1
2
2
3
200
1
$/kappa$
$/kappa$
 
   R
G
$/Omega$
=
=
1
200
3
\end{problembox}

\begin{problembox}
\textbf{Problem 457} \hfill \textbf{[Neeraj Kumar [NRJ] EXERCISE I (JEE MAIN) Q125]}
\label{q:ext-nrj_all_questions-457}

Conductivity of solution will remain uncharged 
but the cell constant becomes double.
 
 
$/kappa$ =
$/times$
=
$/times$
=
$/times$
 
-
-
-
G G
$/Omega$ cm
1
50
1 5
5
3
5 10 2
1
1
.
 
    
=
=
$/times$
$/times$
=
-
-
-
-
eq
K
C
$/Omega$ cm eq
3
5
10
0 05 10
120
2
3
1
2
1
.
 
 
Visit www.jeebooks.in for more
\end{problembox}

\begin{problembox}
\textbf{Problem 458} \hfill \textbf{[Neeraj Kumar [NRJ] EXERCISE II (JEE ADVANCED) Q1]}
\label{q:ext-nrj_all_questions-458}

(I) Cu cannot reduce Pb
 
 (II) Pb can reduce Ag
 
 (III) Ag cannot reduce Cu.
 
 Hence, reducing power: Pb > Cu > Ag
\end{problembox}

\begin{problembox}
\textbf{Problem 459} \hfill \textbf{[Neeraj Kumar [NRJ] EXERCISE II (JEE ADVANCED) Q2]}
\label{q:ext-nrj_all_questions-459}

E
E
n
=
$/circ$ -
 
0 06
.
log [R]
[O]
 
 $/implies$ 0 24
0 36
0 06
1
.
.
.
log [
[
=
-
R]
O]
 
   [
[
O]
R] = 1
100
\end{problembox}

\begin{problembox}
\textbf{Problem 460} \hfill \textbf{[Neeraj Kumar [NRJ] EXERCISE II (JEE ADVANCED) Q3]}
\label{q:ext-nrj_all_questions-460}

For the complex ion to get oxidised, its reduction 
potential should be low.
\end{problembox}

\begin{problembox}
\textbf{Problem 461} \hfill \textbf{[Neeraj Kumar [NRJ] EXERCISE II (JEE ADVANCED) Q4]}
\label{q:ext-nrj_all_questions-461}

Ag
NH
Ag(NH)
E
+
+
+
$/circ$ =
-
=
2
0 79
0 37
0 42
3
3

  


 



;
.
.
.
V
 
 Now, E
n
Keq
$/circ$ -
 
0 06
.
log
 $/implies$ 0 42
0 06
1
.
.
log
=
 
Kf 
 
 $/implies$ Kf = 107
\end{problembox}

\begin{problembox}
\textbf{Problem 462} \hfill \textbf{[Neeraj Kumar [NRJ] EXERCISE II (JEE ADVANCED) Q5]}
\label{q:ext-nrj_all_questions-462}

Given: Co3+ + e-- 
 Co2+; E$/circ$ = 1.81 V; 
 
 
$/Delta$
$/circ$ = -$/times$
$/times$
G
F 1.81
1
1
 
 
Co(CN)
e
Co(CN)
6
3
6
4
+
-
-
+
 $/to$
 
 E$/circ$ = --0.83 V; 
 
 
$/Delta$
$/circ$ = -$/times$
$/times$ -
G
F
2
1
0 83
(
.
)
 
 
Co
6CN
Co(CN)
2
6
4
+
-
-
+
 $/to$
 
; Kf = 1019; 
 
 
$/Delta$
$/circ$ = -
$/times$
G
RT
ln10
3
19
 
 Required Co
CN
Co(CN)
3
6
3
6
+
-
-
+
 $/to$
 
 Kf = ?; 
 
 
$/Delta$
$/circ$ = -
$/times$
G
RT
lnKf
 
 Now, $/Delta$
$/circ$ = $/Delta$
$/circ$ -$/Delta$
$/circ$ + $/Delta$
$/circ$
G
G
G
G
1
2
3
 
 or, -
 
= -
-
+ -
RT lnK
F)
F
RTln10
f
(
.
.
(
)
1 81
0 83
19
 
 or, RT ln K
F
 
= -
10
2 64
19
f
.
 
 
 $/implies$ log
.
.
.
.
10
2 64
2 303
2 64
0 06
19
K
F
RT
f
= -
= -
 
   Kf = 1063
\end{problembox}

\begin{problembox}
\textbf{Problem 463} \hfill \textbf{[Neeraj Kumar [NRJ] EXERCISE II (JEE ADVANCED) Q6]}
\label{q:ext-nrj_all_questions-463}

E
E
Cu
E
Cu
Cu
Cu
Cu
Cu
Cu
Cu
2
2
2
0 06
2
1
0 03
2
2
+
+
+
=
$/circ$
-
 
=
$/circ$
+
+
+
|
|
|
.
log
[
]
.
log[
]
 
 = 0.34 + 0.03 $/times$ log(0.1) = 0.31 V
 
   ECu/Cu
2+ = - 0.31 V
\end{problembox}

\begin{problembox}
\textbf{Problem 464} \hfill \textbf{[Neeraj Kumar [NRJ] EXERCISE II (JEE ADVANCED) Q7]}
\label{q:ext-nrj_all_questions-464}

E
E
RT
F
Cd
Ag
cell
cell
=
$/circ$
-
 
+
+
2
2
2
ln [
]
[
]
 
 As CN-- will form complex with Ag+ ion in the 
cathodic compartment, Ecell will decrease.
\end{problembox}

\begin{problembox}
\textbf{Problem 465} \hfill \textbf{[Neeraj Kumar [NRJ] EXERCISE II (JEE ADVANCED) Q8]}
\label{q:ext-nrj_all_questions-465}

-$/Delta$
=
 
=
$/circ$
-
 
=
$/circ$
-
+
+
G
nF E
nF E
RT
nF
Zn
Cu
nFE
cell
cell
Cell
Cell
[
ln [
]
[
]
2
2
RT
C
C
 ln
1
2
\end{problembox}

\begin{problembox}
\textbf{Problem 466} \hfill \textbf{[Neeraj Kumar [NRJ] EXERCISE II (JEE ADVANCED) Q9]}
\label{q:ext-nrj_all_questions-466}

E
E
RT
F
K
X
Ag X|Ag
Ag |Ag
sp
$/circ$
=
$/circ$
-
 
-
+
|
ln 1
\end{problembox}

\begin{problembox}
\textbf{Problem 467} \hfill \textbf{[Neeraj Kumar [NRJ] EXERCISE II (JEE ADVANCED) Q10]}
\label{q:ext-nrj_all_questions-467}

Cell reaction: H+(cathode, PH = 3) 
 H+ 
(anode, PH = ?)
 
 
E
H
anode
H
P
P
cell
anode
H
catho
=
-
 
=
-
+
+
0
0 059
1
0 059
.
log [
]
[
]cathode
.
de
H
  
  
 
 or, 0 272
0 059
3
.
.
[P
]
=
-
anode
H
 $/implies$ Panode
H
 = 7.6
\end{problembox}

\begin{problembox}
\textbf{Problem 468} \hfill \textbf{[Neeraj Kumar [NRJ] EXERCISE II (JEE ADVANCED) Q11]}
\label{q:ext-nrj_all_questions-468}

Cl2 + H2O   Cl-- + ClO-- + 2H+; 
 
 
E
V
cell
$/circ$
=
-
= -
1 36
1 63
0 27
.
.
.
 
 Now, E
E
H
cell
cell
=
$/circ$
-
+
0 06
1
2
.
log[
]
 
 or, 0
0 27
0 06
1
2
2 25
= -
+
$/times$
=
.
.
.
PH
\end{problembox}

\begin{problembox}
\textbf{Problem 469} \hfill \textbf{[Neeraj Kumar [NRJ] EXERCISE II (JEE ADVANCED) Q12]}
\label{q:ext-nrj_all_questions-469}

E
E
E
cell
quinohydrone
calomel
=
-
 
 0.210 = Equinohydrone -- 0.279 
 
 $/implies$ Equinohydrone = 0.489 V
EXERCISE II (JEE ADVANCED)
 
 
Visit www.jeebooks.in for more
\end{problembox}

\begin{problembox}
\textbf{Problem 470} \hfill \textbf{[Neeraj Kumar [NRJ] EXERCISE II (JEE ADVANCED) Q13]}
\label{q:ext-nrj_all_questions-470}

Anode: Ag(s) 
 Ag+(aq) + e-- 1$/times$ 6
 
 Cathode: 
 
 
Cr O
aq
H
aq
Cr
aq
H O(l)
2
7
2
3
2
14
6
2
7
-
+
-
+
+
+
 $/to$
 
+
(
)
(
)
e
(
)
 
 Net: 
 
 
6
14
6
2
7
2
3
2
Ag(s)
Cr O
aq
H
aq
Ag
aq)
2Cr
aq)+7H O(l)
+
+
 $/to$
 
+
-
+
+
+
(
)
(
)
(
(
 
 
E
E
n
Ag
Cr
Cr O
H
cell
cell
=
$/circ$
-
 
+
+
-
+
0 06
6
3
2
2
7
2
14
.
log [
] [
]
[
][
]
 
 = (1.33 -- 0.80) -- 0 06
6
0 1
0 4
1 6
0 1
0 46
6
2
14
.
log ( . )
( . )
.
( . )
.
 
$/times$
$/times$
=
V
\end{problembox}

\begin{problembox}
\textbf{Problem 471} \hfill \textbf{[Neeraj Kumar [NRJ] EXERCISE II (JEE ADVANCED) Q14]}
\label{q:ext-nrj_all_questions-471}

Net cell reaction: Zn -- Hg(C1M) 
 
 
  
Zn -- Hg(C2M)
 
E
C
C
V
cell =
-
= -
=
0
0 059
2
0 059
2
1
10
0 0295
2
1
.
log
.
log
.
\end{problembox}

\begin{problembox}
\textbf{Problem 472} \hfill \textbf{[Neeraj Kumar [NRJ] EXERCISE II (JEE ADVANCED) Q15]}
\label{q:ext-nrj_all_questions-472}

E
K
cell
eq
$/circ$
=
 
0 06
2
.
log
 
 
 $/implies$ 0 75
1 50
3 1 68 1
3 1
0 03
.
.
.
.
log
-
$/times$
-
$/times$
-
=
Keq
 
   Keq = 10--22
\end{problembox}

\begin{problembox}
\textbf{Problem 473} \hfill \textbf{[Neeraj Kumar [NRJ] EXERCISE II (JEE ADVANCED) Q16]}
\label{q:ext-nrj_all_questions-473}

[
]
.
.
H
Ka C
C M
left
1
+
-
=
 
=
$/times$
$/times$
=
1 8 10
0 1
5
 
 
[
]
.
.
H
K
Kb C
C M
right
w
+
-
-
=
 
=
$/times$
$/times$
=
10
1 8 10
0 01
14
5
2
 
 Net cell reaction, assuming as concentration cell:
 
 H+(C2M) 
 H+ (C1M)
 
 
E
C
C
V
cell =
-
 
= -
0
0 06
1
0 465
1
2
.
log
.
\end{problembox}

\begin{problembox}
\textbf{Problem 474} \hfill \textbf{[Neeraj Kumar [NRJ] EXERCISE II (JEE ADVANCED) Q17]}
\label{q:ext-nrj_all_questions-474}

E
V
V
V
$/circ$
= $/times$
+ $/times$
-$/times$
=
+
+
2
3
1
0 616
1
0 439
1
0 799
0 256
|
.
.
.
.
 
   E
V
V
V
$/circ$
= -
+
+
3
2
0 256
|
.
\end{problembox}

\begin{problembox}
\textbf{Problem 475} \hfill \textbf{[Neeraj Kumar [NRJ] EXERCISE II (JEE ADVANCED) Q18]}
\label{q:ext-nrj_all_questions-475}

E
E
n
H
P
cell
cell
H
=
$/circ$
-
 
+
0 06
2
2
.
log [
]
 
 or, 0 70
0 28
0
0 06
2
1
2
.
( .
)
.
log [
]
=
-
-
 
+
H
 $/implies$ PH = 7.0
\end{problembox}

\begin{problembox}
\textbf{Problem 476} \hfill \textbf{[Neeraj Kumar [NRJ] EXERCISE II (JEE ADVANCED) Q19]}
\label{q:ext-nrj_all_questions-476}

Net cell reaction: 
 
 
Ag
C
M)
Ag
C
K
M
sp
+
+
=
 $/to$
 
=
 
  
 
  
 
  
 
  
(
.
/
1
2
1 3
0 1
2
4
 
 Now, E
C
C
cell =
-
 
0
0 06
1
2
1
.
log
 
 
$/implies$ 0.162 = -
 
0 06
2
0 1
1 3
.
log
(
)
.
/
Ksp
 
   Ksp = 4 $/times$ 10--12
\end{problembox}

\begin{problembox}
\textbf{Problem 477} \hfill \textbf{[Neeraj Kumar [NRJ] EXERCISE II (JEE ADVANCED) Q20]}
\label{q:ext-nrj_all_questions-477}

E
E
V
Tl
Tl
Pb
Pb
+
+
-
= -
|
|
.
2
0 444
 
 or, E
E
Pb
Tl
Tl
Tl
Pb
Pb
+
+
-
(
) -
 
= -
+
+
|
|
.
log [
]
[
]
.
2
0 06
2
0 444
2
2
 
 or, [(
.
)
(
.
)]
.
log
.
.
.
-
--
-
 
 
  
 
  
= -
0 336
0 126
0 03
0 1
0 1
0 444
2
Ksp
 
   Ksp = 4 $/times$ 10--6
\end{problembox}

\begin{problembox}
\textbf{Problem 478} \hfill \textbf{[Neeraj Kumar [NRJ] EXERCISE II (JEE ADVANCED) Q21]}
\label{q:ext-nrj_all_questions-478}

Ecell = ECu -- EZn = (ECu -- Ecalomel) -- (EZn -- Ecalomel)
 
 From question Ecalomel -- EZn = 1.083 V
 
 and Ecalomel -- ECu = -- 0.018 V
\end{problembox}

\begin{problembox}
\textbf{Problem 479} \hfill \textbf{[Neeraj Kumar [NRJ] EXERCISE II (JEE ADVANCED) Q22]}
\label{q:ext-nrj_all_questions-479}

The potential of hydrogen electrode at H2(1 bar) 
may be expressed as E = -- 0.059 PH
 
 Now, E
PKa
1
0 059
= -
+
 
  
 
  
.
log x
y 
 
 and E
PKa
2
0 059
= -
+
 
  
 
  
.
log y
x
 
   P
E
E )
Ka
1
2
= -
+
(
.0 118
 
 
Visit www.jeebooks.in for more
\end{problembox}

\begin{problembox}
\textbf{Problem 480} \hfill \textbf{[Neeraj Kumar [NRJ] EXERCISE II (JEE ADVANCED) Q23]}
\label{q:ext-nrj_all_questions-480}

For the given reaction, E
E
E
cell
given
hydrogen
$/circ$
=
$/circ$
-
$/circ$
 
   (--0.84) -- 0 =
 
0 06
1
.
logKeq $/implies$ Keq = 10--14
\end{problembox}

\begin{problembox}
\textbf{Problem 481} \hfill \textbf{[Neeraj Kumar [NRJ] EXERCISE II (JEE ADVANCED) Q24]}
\label{q:ext-nrj_all_questions-481}

Net cell reaction may be written as 
 
 H2 + Zn2+   2H+ + Zn
 
 
E
E
n
H
P
cell
cell
H
=
$/circ$
-
 
+
+
0 06
2
2
2
.
log
[
]
[Zn
]
 
 or, (
.
)
(
.
)
.
log [
]
.
-
= -
-
 
$/times$
+
0 61
0 76
0 06
2
1
0 4
2
H
 
   [H+] = 2 $/times$ 10--3 M
 
 Now, K
H
SO
HSO
a2
3
2
3
3
2
2
10
6 4
10
0 4
=
=
$/times$
$/times$
$/times$
+
-
-
-
-
[
][
]
[
]
(
)
( .
)
.
 
= 3.2 $/times$ 10--4
\end{problembox}

\begin{problembox}
\textbf{Problem 482} \hfill \textbf{[Neeraj Kumar [NRJ] EXERCISE II (JEE ADVANCED) Q25]}
\label{q:ext-nrj_all_questions-482}

Cu
NH
Cu(NH
Kf
2
3
3 4
2
12
4
10
+
+
+
=

  


 



)
,
 
             1.0 M   excess
 
 100/%        0                             1.0 M
 
 Equ.          x      2.0 M            1.0 M
 
 
10
1 0
2 0
12
4
=
$/times$
.
( . )
x
 $/implies$ x =
=
$/times$
-
-
10
16
6 25 10
12
14
.
 
 Now, cell reaction: Zn + Cu2+   Zn2+ + Cu
 
 
E
E
Zn
Cu
cell
cell
=
$/circ$
-
 
+
+
0 06
2
2
2
.
log [
]
[
]
 
 
=
-
-
 
$/times$
=
-
[ .
( .
)]
.
log
.
.
.
0 34
0 76
0 06
2
1 0
6 25 10
0 704
14
V
\end{problembox}

\begin{problembox}
\textbf{Problem 483} \hfill \textbf{[Neeraj Kumar [NRJ] EXERCISE II (JEE ADVANCED) Q26]}
\label{q:ext-nrj_all_questions-483}

Net cell reaction: Zn + 2H+   Zn2+ + H2
 
 
E
E
Zn
P
H
cell
cell
H
=
$/circ$
-
 
+
+
0 06
2
2
2
2
.
log
[
]
[
]
 
 or, 0 70
0
0 76
0 06
2
0 01 1
2
.
[
(
.
)]
.
log .
[
]
=
--
-
 
$/times$
+
H
 
 
 $/implies$ [H+] = 0.01 M
 
 Moles of HCl in RHS = 500
0 01
1000
5 10 3
$/times$
=
$/times$
-
.
 
   Mass of NaOH needed = 5 $/times$ 10--3 $/times$ 40 = 0.2 gm
\end{problembox}

\begin{problembox}
\textbf{Problem 484} \hfill \textbf{[Neeraj Kumar [NRJ] EXERCISE II (JEE ADVANCED) Q27]}
\label{q:ext-nrj_all_questions-484}

On assuming concentration cell, the net cell 
reaction is Ag+ (C1M, Right) 
  Ag+ (C2M, 
left)
 
 Now, C
M
1
0 1
40
100
0 04
=
$/times$
=
.
.
 
 and C
K
Cl
K
K
M
sp
sp
sp
2
0 1
50
100
0 05
=
=
$/times$
=
-
[
]
.
.
 
 Now, E
C
C
cell =
-
0
0 06
1
2
1
.
log
 
 or, 0 42
0 06
1
0 05
0 04
.
.
log
/ .
.
= -
Ksp
 $/implies$ Ksp = 2 $/times$ 10--10
\end{problembox}

\begin{problembox}
\textbf{Problem 485} \hfill \textbf{[Neeraj Kumar [NRJ] EXERCISE II (JEE ADVANCED) Q28]}
\label{q:ext-nrj_all_questions-485}

$/Delta$Gcell = -- nFEcell $/implies$ -- 965 $/times$ 3 $/times$ 103 =  --12 $/times$ 96500 
$/times$ Ecell
 
   Ecell = 2.5 V
\end{problembox}

\begin{problembox}
\textbf{Problem 486} \hfill \textbf{[Neeraj Kumar [NRJ] EXERCISE II (JEE ADVANCED) Q29]}
\label{q:ext-nrj_all_questions-486}

Theoretical e   ciency = -$/Delta$
$/circ$
-$/Delta$
$/circ$
G
H $/implies$ 0 84
285
.
= -$/Delta$
$/circ$
G
 
   $/Delta$G$/circ$ = -- 0.84 $/times$ 285 KJ = --nF   E$/circ$cell
 
 or, 0.84 $/times$ 285 $/times$ 103 = 2 $/times$ 96500 $/times$ E$/circ$cell
 
 
 $/implies$ E$/circ$cell = 1.24 V
\end{problembox}

\begin{problembox}
\textbf{Problem 487} \hfill \textbf{[Neeraj Kumar [NRJ] EXERCISE II (JEE ADVANCED) Q30]}
\label{q:ext-nrj_all_questions-487}

Net cell reaction: Ag(s) + H+ + Cl--   AgCl(s) + 
1
2 H2(g) But for Ecell calculation, reaction may be 
written as
 
 Ag(s) + H+   Ag+ + 1
2 H2
 
 Now, E
E
Ag
P
H
cell
cell
H
=
$/circ$
-
 
 
+
+
0 06
1
2
1 2
.
log
[
]
[
]
/
 
 
=
-
-
 
 
  
 
  $/times$
= -
-
[
.
]
.
log
.
.
.
/
0
0 80
0 06
1
10
0 1
1
0 1
0 32
10
1 2
V
 
 It means that actual reaction is in reverse direction 
and Ecell = 0.32 V
\end{problembox}

\begin{problembox}
\textbf{Problem 488} \hfill \textbf{[Neeraj Kumar [NRJ] EXERCISE II (JEE ADVANCED) Q32]}
\label{q:ext-nrj_all_questions-488}

x
y
y
x
 Ag
 NH
Ag(NH
+
+
+
3
3

  


 



)
 
               aM        bM                          0            b >> a
 
 Eqn.        ?           bM                      a
x M
 
 
K
/
Ag
f
x
y
a x
b
=
 
+
[
]
 $/implies$ [Ag ]
/
+ =
 
 
 
  
 
  
a
x b
f
y
x
K
1
 
 Cell reaction: Ag+ (Right) 
 Ag+ (Left)
 
 
Visit www.jeebooks.in for more
\end{problembox}

\begin{problembox}
\textbf{Problem 489} \hfill \textbf{[Neeraj Kumar [NRJ] EXERCISE II (JEE ADVANCED) Q33]}
\label{q:ext-nrj_all_questions-489}

E
K
V
cell
eq
$/circ$
=
 
=
$/times$
= -
0 06
1
0 06
1
1 667
10
0 3732
6
.
log
.
log
.
.
 
 Now, E
E
V
Cu
Cu
Cu
Cu
$/circ$
-
$/circ$
= -
+
+
+
2
0 3732
|
|
.
 
(1)
 
 and E
E
E
V
Cu
Cu
Cu
Cu
Cu
Cu
$/circ$
=
$/times$
$/circ$
+ $/times$
$/circ$
+
=
+
+
+
+
2
2
1
1
1 1
0 3376
|
|
|
.
 
 or, E
E
Cu
Cu
Cu
Cu
$/circ$
+
$/circ$
=
+
+
+
2
0 6752
|
|
.
V 
(2)
 
 From (1) and (2), E
V
Cu
Cu
$/circ$
=
+|
.0 5242
\end{problembox}

\begin{problembox}
\textbf{Problem 490} \hfill \textbf{[Neeraj Kumar [NRJ] EXERCISE II (JEE ADVANCED) Q34]}
\label{q:ext-nrj_all_questions-490}

Zn + Ni2+   Zn2+ + Ni
 
 
E
K
cell
eq
$/circ$
= 0 06
2
.
log
 
 
 $/implies$ (--0.24) -- ( --0.75) = 0.03 log Keq
 
    Keq = 1017 $/implies$ It means that Ni2+ will react 
almost completely and [Zn2+] $/approx$ 1.0 M
 
 Now, 10
1 0
17
2
=
+
.
[Ni
]
 $/implies$ [Ni2+] = 10--17 M
\end{problembox}

\begin{problembox}
\textbf{Problem 491} \hfill \textbf{[Neeraj Kumar [NRJ] EXERCISE II (JEE ADVANCED) Q35]}
\label{q:ext-nrj_all_questions-491}

Assuming the cell as concentration cell, the cell 
reaction may be written as
 
 
Ag
C
M
Ag
C
M
+
-
-
+
-
-
=
$/times$
=
$/times$
 
  
 
  
 $/to$
 
=
$/times$
=
1
13
10
2
10
9
4
10
0 001
4
10
2
10
0 2
10
.
.
 
  
 
  
 
 Now, E
V
cell =
-
 
$/times$
= -
-
-
0
0 06
1
10
4
10
0 024
9
10
.
log
.
\end{problembox}

\begin{problembox}
\textbf{Problem 492} \hfill \textbf{[Neeraj Kumar [NRJ] EXERCISE II (JEE ADVANCED) Q36]}
\label{q:ext-nrj_all_questions-492}

0 03
0 06
2
0 3
0 5
2
2
2
.
.
log
[
]
[
]
.
log
.
[
]
lower
=
 
=
 
+
+
+
Cu
Cu
Cu
higher
lower
 
   [Cu2+]lower = 0.05 M
\end{problembox}

\begin{problembox}
\textbf{Problem 493} \hfill \textbf{[Neeraj Kumar [NRJ] EXERCISE II (JEE ADVANCED) Q37]}
\label{q:ext-nrj_all_questions-493}

Cell reaction: H+(C1, HA1) 
 H+(C2, HA2)
 
 
E
C
C
Ka
Ka
cell =
-
 
 
  
 
  =
 
0
0 059
1
0 059
2
1
2
1
.
log
.
log
 
 
=
-
=
0 059
2
0 059
1
2
.
(
)
.
V
P
P
K
K
a
a
\end{problembox}

\begin{problembox}
\textbf{Problem 494} \hfill \textbf{[Neeraj Kumar [NRJ] EXERCISE II (JEE ADVANCED) Q38]}
\label{q:ext-nrj_all_questions-494}

Au+ + 2CN--   Au (
)
CN 2
- , $/Delta$G$/circ$1 = --RT   ln x
 
 O2 + 2H2  + 4e--   40H-- ; $/Delta$G$/circ$2 = --4 $/times$ F $/times$ 0.41
 
 Au3+ + 3e--   Au; $/Delta$G$/circ$3 = --3 $/times$ F $/times$ 1.50
 
 Au3+ + 2e--   Au+; $/Delta$G$/circ$4 = --2 $/times$ F $/times$ 1.40
 
 From $/Delta$
$/circ$ +
$/Delta$
$/circ$ -$/Delta$
$/circ$ + $/Delta$
$/circ$
$/Delta$
$/circ$
= -
+
G
G
G
G
G
RT
 
F
required
1
2
3
4
1
4
1 29
,
ln
.
x
\end{problembox}

\begin{problembox}
\textbf{Problem 495} \hfill \textbf{[Neeraj Kumar [NRJ] EXERCISE II (JEE ADVANCED) Q40]}
\label{q:ext-nrj_all_questions-495}

PH of right electrode will increase due to formation 
of OH-- ion.
\end{problembox}

\begin{problembox}
\textbf{Problem 496} \hfill \textbf{[Neeraj Kumar [NRJ] EXERCISE II (JEE ADVANCED) Q43]}
\label{q:ext-nrj_all_questions-496}

At low [Cl--], O2 becomes anode product
\end{problembox}

\begin{problembox}
\textbf{Problem 497} \hfill \textbf{[Neeraj Kumar [NRJ] EXERCISE II (JEE ADVANCED) Q46]}
\label{q:ext-nrj_all_questions-497}

In the electrolysis of aq. KNO3, neither K+ are 
NO3
- participate in electrode reaction.
\end{problembox}

\begin{problembox}
\textbf{Problem 498} \hfill \textbf{[Neeraj Kumar [NRJ] EXERCISE II (JEE ADVANCED) Q47]}
\label{q:ext-nrj_all_questions-498}

Na will react with water. S will not conduct 
electricity.
\end{problembox}

\begin{problembox}
\textbf{Problem 499} \hfill \textbf{[Neeraj Kumar [NRJ] EXERCISE II (JEE ADVANCED) Q48]}
\label{q:ext-nrj_all_questions-499}

w
E
Q
F
=
 $/implies$ 0 635
63 5
2
0 965 3600
96500
.
.
.
$/times$
= $/times$
$/times$
i
 
 
 $/implies$ i = 2
3 6. A
 
   /%
.
.
.
/%
error =
-
=
2
3 6
0 5
2
3 6
10
\end{problembox}

\begin{problembox}
\textbf{Problem 500} \hfill \textbf{[Neeraj Kumar [NRJ] EXERCISE II (JEE ADVANCED) Q49]}
\label{q:ext-nrj_all_questions-500}

Detonating gas is mixture of H2 and O2
 
 neq Sn = neq H2 = neq O2 $/implies$ 1 2
120
2
2
4
2
2
.
$/times$
=
$/times$
=
$/times$
n
n
H
O
 
 
   n
and n
H
O
2
1
0 01
0 005
=
=
.
.
 
 
Visit www.jeebooks.in for more
\end{problembox}

\begin{problembox}
\textbf{Problem 501} \hfill \textbf{[Neeraj Kumar [NRJ] EXERCISE II (JEE ADVANCED) Q50]}
\label{q:ext-nrj_all_questions-501}

neq of Li OH = Q
F $/implies$ w
24
1
2 5
4825
0 8
96500
$/times$ =
$/times$
$/times$
.
.
\end{problembox}

\begin{problembox}
\textbf{Problem 502} \hfill \textbf{[Neeraj Kumar [NRJ] EXERCISE II (JEE ADVANCED) Q51]}
\label{q:ext-nrj_all_questions-502}

neq Cu deposited at cathode = Q
F 
 
 or, w
63 5
2
12 4
4825
96500
.
.
$/times$
=
$/times$
 $/implies$ w = 19.685 gm
 
 But the increase in mass of cathode is only 19.05 
gm It represents that 20 gm of sample contains 
only 19.05 gm Cu.
 
   /% of Cu = 19 05
20
100
95 25
.
.
/%
$/times$
=
 
 Now, Q
F
n Cu
n Fe
eq
eq
=
+
 (oxidised at anode)
 
 or, 12 4
4825
96500
19 05
63 5
2
56
2
.
.
.
$/times$
=
$/times$
+
$/times$
w
 $/implies$ w = 0.56 gm
 
   Percentage of Fe = 0 56
20
100
2 8
.
. /%
$/times$
=
\end{problembox}

\begin{problembox}
\textbf{Problem 503} \hfill \textbf{[Neeraj Kumar [NRJ] EXERCISE II (JEE ADVANCED) Q52]}
\label{q:ext-nrj_all_questions-503}

Theoretical neq of NaOH formed = neqCu 
 
  
 
 
 
= 3 18
63 6
2
0 1
.
.
.
$/times$
=
 
 Actual neq of NaOH formed =
$/times$
=
60
1
1000
0 06
.
 
   Percentage yeild = 
0 06
0 1 100
60
.
.
/%
$/times$
=
\end{problembox}

\begin{problembox}
\textbf{Problem 504} \hfill \textbf{[Neeraj Kumar [NRJ] EXERCISE II (JEE ADVANCED) Q53]}
\label{q:ext-nrj_all_questions-504}

Cell reaction during discharge:
 
 Pb + PbO2 + 2H2SO4 
 2PbSO4 + 2H2O
 
 nPb taken = 200
208 and nPbO taken
2
200
240
=
 
 
 Hence, PbO2 is L.R.
 
 Now, neqPbO
Q
F
2 =
 $/implies$ 200
240
2
10
96500
$/times$
=
$/times$t 
 
 $/implies$ t = 16083.33 sec
\end{problembox}

\begin{problembox}
\textbf{Problem 505} \hfill \textbf{[Neeraj Kumar [NRJ] EXERCISE II (JEE ADVANCED) Q55]}
\label{q:ext-nrj_all_questions-505}

Q
F
of I
Na S O
eq
eq
2
2
3
=
=
-
n
n
 
 or, i $/times$
$/times$
=
$/times$
$/times$
2
3600
96500
72
1 0
1000
1
.
 $/implies$ i = 0.965 A
\end{problembox}

\begin{problembox}
\textbf{Problem 506} \hfill \textbf{[Neeraj Kumar [NRJ] EXERCISE II (JEE ADVANCED) Q56]}
\label{q:ext-nrj_all_questions-506}

C14H10 + 2H2O 
 C14H8O2 + 6H+ + 6e--
 
 neq C14H8O2 = Q
F $/implies$ w
208
6
1
40
60
0 965
96500
$/times$
=
$/times$
$/times$
$/times$ .
 
   w = 0.832 gm
\end{problembox}

\begin{problembox}
\textbf{Problem 507} \hfill \textbf{[Neeraj Kumar [NRJ] EXERCISE II (JEE ADVANCED) Q57]}
\label{q:ext-nrj_all_questions-507}

Number of coulombs required =
1000
0 00033
.
 per Kg Cu
 
 Energy required =
$/times$
=
1000
0 00033
0 33
106
.
.
J
 
   Cost of electricity 
 
  
 
=
$/times$
$/times$
=
4
10
3600
10
1 11
3
6
.
Rupee
\end{problembox}

\begin{problembox}
\textbf{Problem 508} \hfill \textbf{[Neeraj Kumar [NRJ] EXERCISE II (JEE ADVANCED) Q58]}
\label{q:ext-nrj_all_questions-508}

neq CH3Coo-- oxidised = Q
F
 
 or, n $/times$ =
$/times$
$/times$
$/times$
1
0 5
482 5
60
0 8
96500
.
.
. $/implies$ n = 0.12
 
   Moles of (C2H6 + CO2) produced 
 
  
 
 = 3
2
0 12
0 18
$/times$
=
.
.
 
 and total volume = 0.18 $/times$ 22.4 = 4.032 L
\end{problembox}

\begin{problembox}
\textbf{Problem 509} \hfill \textbf{[Neeraj Kumar [NRJ] EXERCISE II (JEE ADVANCED) Q59]}
\label{q:ext-nrj_all_questions-509}

$/Delta$ $/circ$ =
 
=
E
V
0 059
2
10
0 177
6
.
log
.
\end{problembox}

\begin{problembox}
\textbf{Problem 510} \hfill \textbf{[Neeraj Kumar [NRJ] EXERCISE II (JEE ADVANCED) Q60]}
\label{q:ext-nrj_all_questions-510}

Back EMF =
 
=
$/times$
-
0 06
2
0 12
0 08
5 4
10 3
.
log .
.
.
V
\end{problembox}

\begin{problembox}
\textbf{Problem 511} \hfill \textbf{[Neeraj Kumar [NRJ] EXERCISE II (JEE ADVANCED) Q61]}
\label{q:ext-nrj_all_questions-511}

Equivalent of charge used 
 
 
= 0 4825 10
3600
96500
0 18
.
.
$/times$
$/times$
=
F
 
 Cell reaction during charge:
 
   Cu
Zn
Cu
Zn
+
 $/to$
 
+
+
+
2
2
 
                  
100
1
1000
0 1
$/times$
= . mole 
100
1
1000
0 1
$/times$
= . mole
 
                  = 0.2 eq                    = 0.2 eq
 
 Final           0.2 -- 0.1 8             = 0.2 + 0.18
 
                  = 0.02 eq                  = 0.38 eq
 
                  = 0.01 mole              = 0.19 mole
 
   Final [
]
.
.
Zn
M
2
0 01
100
1000
0 1
+ =
$/times$
=
 
 and [
]
.
.
Cu
M
2
0 19
100
1000
1 9
+ +
$/times$
=
 
 Now, cell reaction as galvanic cell:
 
 
Visit www.jeebooks.in for more
\end{problembox}

\begin{problembox}
\textbf{Problem 512} \hfill \textbf{[Neeraj Kumar [NRJ] EXERCISE II (JEE ADVANCED) Q62]}
\label{q:ext-nrj_all_questions-512}

Reactions involved are
 
 
2
2
2
2
H O
H
O
2
electrolysis
 
$/to$
    
+
 
 N2 + 3H2 
 2NH3
 
 NH3 + 2O2 
 HNO3 + H2O
 
 NH3 + HNO3 
 NH4NO3
 
 For 1 mole NH4NO3, 3 moles of H2O should be 
electrolyzed. Hence for 1152 Kg NH4NO3, moles 
of H2 needed =
$/times$
$/times$
=
$/times$
3
1152
10
80
4 32
10
3
4
.
 
 Now, neqH
Q
F
2 =
 $/implies$ 4.32 $/times$ 104 $/times$ 2 = i $/times$
$/times$
24
3600
96500
 
   i = 96500 A/day
\end{problembox}

\begin{problembox}
\textbf{Problem 513} \hfill \textbf{[Neeraj Kumar [NRJ] EXERCISE II (JEE ADVANCED) Q63]}
\label{q:ext-nrj_all_questions-513}

neq HC3H5O3 = nOH
Q
F
-=
 
 or, w
90
1
50
10
1158
96500
3
$/times$ =
$/times$
$/times$
-
 $/implies$ w = 0.054 gm
 
   /% of lactice acid = 0 054
1
100
5 4
.
. /%
$/times$
=
\end{problembox}

\begin{problembox}
\textbf{Problem 514} \hfill \textbf{[Neeraj Kumar [NRJ] EXERCISE II (JEE ADVANCED) Q64]}
\label{q:ext-nrj_all_questions-514}

n
n
n
H O
NH
S O
formed
2
2
4 2
2
8
=
=
(
)
 
 and neq NH
S O
Q
F
(
)
4 2
2
8 =
 
 
 $/implies$ n
i
$/times$
=
$/times$
=
$/times$
$/times$
2
102
34
2
3600
0 5
96500
.
 
 
 
[
]
2
2
4
2
2
8
2
SO
S O
e
-
-
-
 $/to$
 
+
  
$/implies$ i = 321.67 A
\end{problembox}

\begin{problembox}
\textbf{Problem 515} \hfill \textbf{[Neeraj Kumar [NRJ] EXERCISE II (JEE ADVANCED) Q65]}
\label{q:ext-nrj_all_questions-515}

n
n
n
As
H AsO
=
=
3
3
 and neq H3AsO3 = neqI2 = Q
F
 
 or, w
75
2
1 68 10
96 5
96500
3
$/times$
=
$/times$
$/times$
-
.
. $/implies$ w = 6.3 $/times$ 10--5 gm
\end{problembox}

\begin{problembox}
\textbf{Problem 516} \hfill \textbf{[Neeraj Kumar [NRJ] EXERCISE II (JEE ADVANCED) Q66]}
\label{q:ext-nrj_all_questions-516}

w
E
Q
F
=
 $/implies$ 52 2
87
2
19 3
2
3600
96500
.
.
$/times$
=
$/times$
$/times$
$/times$  
 
 $/implies$   = 0.8333 or 83.33/%
\end{problembox}

\begin{problembox}
\textbf{Problem 517} \hfill \textbf{[Neeraj Kumar [NRJ] EXERCISE II (JEE ADVANCED) Q67]}
\label{q:ext-nrj_all_questions-517}

neqCu2+ reduced = Q
F $/implies$ n $/times$ 2 = 2
10
19 3
60
96500
3
$/times$
$/times$
$/times$
-
.
 
   n = 1.2 $/times$ 10--5
  [
]
( .
)
.
CuSO
M
4 0
5
5
1 2
10
2
250
1000
9 6
10
=
$/times$
$/times$
$/times$
=
$/times$
-
-
\end{problembox}

\begin{problembox}
\textbf{Problem 518} \hfill \textbf{[Neeraj Kumar [NRJ] EXERCISE II (JEE ADVANCED) Q68]}
\label{q:ext-nrj_all_questions-518}

w
E
Q
F
=
 $/implies$ 12 3
123
6
. $/times$
=
$/times$
Q
0.5
F
 $/implies$ Q = 1.2 F
+ 6H+ + 6e--
+ 2H2O
NH2
NO2
 
 and Energy consumed = 1.2 F $/times$ 3.0 V = 347.4 KJ
\end{problembox}

\begin{problembox}
\textbf{Problem 519} \hfill \textbf{[Neeraj Kumar [NRJ] EXERCISE II (JEE ADVANCED) Q69]}
\label{q:ext-nrj_all_questions-519}

neq H2 (at cathode) = neq O2 + neq H2S2O8 
(at anode)
 
 or, 9 08
22 7
2
2 27
22 7
4
194
2
.
.
.
.
$/times$
=
$/times$
+
$/times$
w
 
 
   w = 38.8 gm
\end{problembox}

\begin{problembox}
\textbf{Problem 520} \hfill \textbf{[Neeraj Kumar [NRJ] EXERCISE II (JEE ADVANCED) Q70]}
\label{q:ext-nrj_all_questions-520}

Initial mass of H2SO4, 
 
 w1 = 3600 $/times$ 1.5 $/times$ 40
100
2160
=
gm
 
 Final mass of H2SO4, 
 
 w2 = 3600 $/times$ 1.1 $/times$ 10
100
396
=
gm
 
   Moles of H2SO4 consumed = w
w
1
2
98
18
-
=
 
 Now, neq H2SO4 = Q
F $/implies$ 18 $/times$ 1 = (
)
amp-hr $/times$3600
96500
 
   Number of ampere-hr = 482.5
\end{problembox}

\begin{problembox}
\textbf{Problem 521} \hfill \textbf{[Neeraj Kumar [NRJ] EXERCISE II (JEE ADVANCED) Q71]}
\label{q:ext-nrj_all_questions-521}

 
 
=
=
 
 
 
 
 
m
m
NaCl
KCl
NaCl
NaCl
KCl
NaCl
KCl
/C)
/C)
R
G
C
R
G
(
)
(
)
(
(
$/kappa$
$/kappa$
1
1
1
 
=
 
 
1
C
R C)
R C)
KCl
KCl
NaCl
(
(
 
 or,  
=
$/times$
$/times$
m NaCl)
(
.
.
120
200
0 1
6400
0 003 
 
 $/implies$  m(NaCl) = 125 $/Omega$--1 cm--1 mol--1
2
2
4
2
2
8
2
SO
S O
e
-
-
-
$/to$
+
  
  
 
 
Visit www.jeebooks.in for more
\end{problembox}

\begin{problembox}
\textbf{Problem 522} \hfill \textbf{[Neeraj Kumar [NRJ] EXERCISE II (JEE ADVANCED) Q72]}
\label{q:ext-nrj_all_questions-522}

 $/circ$
=
$/times$
$/circ$
+
+
-
m
m
m
NH
CrO
NH
CrO
[(
)
]
(
)
(
)
4 2
4
4
4
2
2
$/lambda$
$/lambda$
 
 = (2 $/times$ 6.6 $/times$ 10--8 + 5.4 $/times$ 10--8) $/times$ 96500 
 
 = 0.01795 $/Omega$--1 m2 mol--1
\end{problembox}

\begin{problembox}
\textbf{Problem 523} \hfill \textbf{[Neeraj Kumar [NRJ] EXERCISE II (JEE ADVANCED) Q73]}
\label{q:ext-nrj_all_questions-523}

 
=
  $/circ$
=
m
m
C
$/alpha$
$/kappa$
 
 or, 0.9 $/times$ 4.25 $/times$ 10--2 = 382 5
3
.
C
10
$/times$
 $/implies$ C = 0.1 M
\end{problembox}

\begin{problembox}
\textbf{Problem 524} \hfill \textbf{[Neeraj Kumar [NRJ] EXERCISE II (JEE ADVANCED) Q74]}
\label{q:ext-nrj_all_questions-524}

 
=
m
K
C $/implies$ 1.5 $/times$ 10--2 = 3 06
10
2 56
10
3
3
.
.
$/times$
-
$/times$
-
-
C
 
   C
mol m
mol l
V
=
=
$/times$
 
=
-
-
-
1
30
1
30 10
585 58 5
3
3
1
/
.
 
   V = 3 $/times$ 105 L
\end{problembox}

\begin{problembox}
\textbf{Problem 525} \hfill \textbf{[Neeraj Kumar [NRJ] EXERCISE II (JEE ADVANCED) Q75]}
\label{q:ext-nrj_all_questions-525}

 
=
=
 
 
m
C
G G
C
$/kappa$
$/implies$ 100
1
0 5
1 5
0 1 10 3
=
$/times$
$/times$
-
R
.
.
.
 
 
 $/implies$ R
$/Omega$
= 100
3
 
 Now, V = IR $/implies$ I
V
R
A
=
=
=
5
100 3
0 15
/
.
\end{problembox}

\begin{problembox}
\textbf{Problem 526} \hfill \textbf{[Neeraj Kumar [NRJ] EXERCISE II (JEE ADVANCED) Q76]}
\label{q:ext-nrj_all_questions-526}

Ionic mobility, $/mu$
$/lambda$
$/circ$ =
=
$/circ$
speed of ion
Pot. gradient
F
m
 
 or, speed
19 3
5
50
96500
.
 
  
 
  
=
 $/implies$ speed = 2 $/times$ 10--3 cm/s
\end{problembox}

\begin{problembox}
\textbf{Problem 527} \hfill \textbf{[Neeraj Kumar [NRJ] EXERCISE II (JEE ADVANCED) Q77]}
\label{q:ext-nrj_all_questions-527}

 eq = C
$/kappa$ $/implies$ 150
3 4
10
1 6
10
5
6
6
$/times$
$/times$
-
$/times$
-
-
.
.
 
   S = 1.2 $/times$ 10--8 mol cm--3 = 1.2 $/times$ 10--5 M
\end{problembox}

\begin{problembox}
\textbf{Problem 528} \hfill \textbf{[Neeraj Kumar [NRJ] EXERCISE II (JEE ADVANCED) Q78]}
\label{q:ext-nrj_all_questions-528}

K = G   G* $/implies$ $/kappa$
1 4
1 280
1 50
.
/
/
=
 
 
 $/implies$ $/kappa$ = 0.25 s m--1 for 0.5 M
 
 Now,  
=
=
$/times$
=
$/times$
-
-
m
C
s m mol
$/kappa$
0 25
0 5 10
5 10
3
4
2
1
.
.
\end{problembox}

\begin{problembox}
\textbf{Problem 529} \hfill \textbf{[Neeraj Kumar [NRJ] EXERCISE II (JEE ADVANCED) Q79]}
\label{q:ext-nrj_all_questions-529}

Ag A S
Ag
A
( )
(
)M

  


 



+
+
-
+
x
y
xM
 
 
 
Ag B S
Ag
B
M
( )
(
)M

  


 



+
+
-
+
x
y
y
 
 Now, (x + y)   x = 3 $/times$ 10--14 and 
 
           (x + y)   y = 1 $/times$ 10--14
 
   [Ag ]
,
[A ]
.
;
[B ]
.
+
-
-
-
-
-
=
+
=
$/times$
=
=
$/times$
=
=
$/times$
x
y
x
y
2
10
1 5 10
0 5 10
7
7
7
M
M
M
 
 
 Now, 
$/kappa$solution =
$/times$
=
$/times$
$/times$
$/times$
+
$/times$
$/times$
$/times$
+
$/times$
-
-
-
-
-
3 75 10
2 10
10
60
1 5 10
10
80
0 5 10
8
7
3
7
3
.
.
.
-
-
$/circ$
$/times$
7
3
10
$/lambda$ B
 
 
   $/lambda$B
$/circ$
-
-
= 135 $/Omega$
cm mol
1
2
1
\end{problembox}

\begin{problembox}
\textbf{Problem 530} \hfill \textbf{[Neeraj Kumar [NRJ] EXERCISE II (JEE ADVANCED) Q80]}
\label{q:ext-nrj_all_questions-530}

 $/circ$
=  $/circ$
+  $/circ$
- $/circ$
eq
eq
eq
eq
[Be (Po ) ]
[BeCl ]
[K Po ]
[K
]
3
4 2
2
3
4
Cl
 
 
=
+
-
=
-
-
160
140
100
200
1
2
1
$/Omega$
cm eq
 
 Now, neq =
$/implies$
=
$/times$
-
$/kappa$
C
C
200
1 2 10 5
.
 
 
 $/implies$ C = 6 $/times$ 10--8 eq cm--3 
 
           = 6 $/times$ 10--5 N
 
           = 10--5 M
 
 Now, Ksp =
=
$/times$
=
$/times$
-
-
108
108
10
1 08 10
5
5 5
23
S
(
)
.
 
Section B (One or More than one Correct)
\end{problembox}

\begin{problembox}
\textbf{Problem 531} \hfill \textbf{[Neeraj Kumar [NRJ] EXERCISE II (JEE ADVANCED) Q1]}
\label{q:ext-nrj_all_questions-531}

Net cell reaction is spontaneous in electrochemical 
cell but non-spontaneous in electrolytic cell. 
Cathode is +ve in electrochemical cell but --ve in 
electrolytic cell.
\end{problembox}

\begin{problembox}
\textbf{Problem 532} \hfill \textbf{[Neeraj Kumar [NRJ] EXERCISE II (JEE ADVANCED) Q2]}
\label{q:ext-nrj_all_questions-532}

For the cell: Ag(s) | Ag Cl (s) | Cl- || Ag+ | Ag(s), 
the net cell reaction is Ag+ + Cl-   Ag Cl (s).
\end{problembox}

\begin{problembox}
\textbf{Problem 533} \hfill \textbf{[Neeraj Kumar [NRJ] EXERCISE II (JEE ADVANCED) Q3]}
\label{q:ext-nrj_all_questions-533}

Left electrode: Ag(s) + Cl- (aq) $/to$ Ag Cl (s) + e- 1 $/times$ 2
Right electrode : Hg2Cl2(s) + 2e- $/to$ 2Hg(l) + 2 Cl- (aq)
Net reaction: 2Ag(s) + Hg2Cl2(s) $/to$ 2Ag Cl(s) + 2Hg (l)
 
 
Visit www.jeebooks.in for more
\end{problembox}

\begin{problembox}
\textbf{Problem 534} \hfill \textbf{[Neeraj Kumar [NRJ] EXERCISE II (JEE ADVANCED) Q5]}
\label{q:ext-nrj_all_questions-534}

n
n
n
n
eq
eq
eq
eq
Cu
Mg
Na
Al
=
$/times$
=
=
$/times$
=
=
$/times$ =
=
63 5
63 5
2
2
24
24
2
2
11 5
23
1
0 5
.
.
;
.
. ;
9
27
3
1
$/times$
=
\end{problembox}

\begin{problembox}
\textbf{Problem 535} \hfill \textbf{[Neeraj Kumar [NRJ] EXERCISE II (JEE ADVANCED) Q7]}
\label{q:ext-nrj_all_questions-535}

 $/circ$eq = 60 + 80 = 140 $/Omega$--1 cm2 eq--1
 
  $/circ$m = 140 $/times$ 6 = 840 $/Omega$--1 cm2 eq--1
\end{problembox}

\begin{problembox}
\textbf{Problem 536} \hfill \textbf{[Neeraj Kumar [NRJ] EXERCISE II (JEE ADVANCED) Q8]}
\label{q:ext-nrj_all_questions-536}

Salt bridge does not change standard potential of 
any electrode.
\end{problembox}

\begin{problembox}
\textbf{Problem 537} \hfill \textbf{[Neeraj Kumar [NRJ] EXERCISE II (JEE ADVANCED) Q9]}
\label{q:ext-nrj_all_questions-537}

Moles of electron involved 
 
 = 0 25
9 65
3600
96500
0 09
.
.
.
$/times$
$/times$
=
 
   Mass of Zn involved = 0 09
2
65 4
2 943
.
.
.
$/times$
=
gm
 
 Mass of MnO2 involved = 0.09 $/times$ 87 = 7.83 gm
 
 Mass of NH4
+ involved = 0.09 $/times$ 18 = 1.62 gm
\end{problembox}

\begin{problembox}
\textbf{Problem 538} \hfill \textbf{[Neeraj Kumar [NRJ] EXERCISE II (JEE ADVANCED) Q10]}
\label{q:ext-nrj_all_questions-538}

Net cell reaction is
 
 Cd(s) + 2AgCl(s)   2Ag(s) + Cd2+ (aq) + 2Cl-- 
(aq)
 
 
$/Delta$
$/circ$
= -
$/circ$ = -$/times$
$/times$
= -
$/circ$
G 50
2
96500
0 6
115800
C
nFE
J
.
 
 
 
$/Delta$
$/circ$
= -
$/circ$ = -$/times$
$/times$
= -
$/circ$
G 0
2
96500
0 7
135100
C
nFE
J
.
 
 
$/Delta$ $/circ$ =
 
$/circ$ -
$/circ$
-
 
  
 
  =
$/times$
$/times$
-
= -
S
T
T
nF
E
E
J/K
2
1
2
1
2
96500
0 6
0 7
50
386
.
.
 
 $/Delta$H$/circ$ = $/Delta$G$/circ$ + T   $/Delta$S$/circ$
 
  = (--135100) + 273 $/times$ (--386) = --240478 J
\end{problembox}

\begin{problembox}
\textbf{Problem 539} \hfill \textbf{[Neeraj Kumar [NRJ] EXERCISE II (JEE ADVANCED) Q11]}
\label{q:ext-nrj_all_questions-539}

Anode: 2H2O(l) 
 O2(g) + 4H+(aq) + 4e--
 
 Cathode: 2H2O(l) + 2e-- 
 H2(g) +2OH--(aq)
 
 neq H+ produced = neq OH-- produced = Q
F
 
 or, n
n
H
OH
+
-
$/times$ =
=
$/times$
$/times$
=
1
1 25
965
60
96500
0 75
.
.
 
 Anode: HPO
H
H PO
4
2
2
4
-
+
-
+

  


 



 
                 1.0 M   0.75 M       1.0 M
 
 Final     0.25 M                    1.75 M
 
   P
P
HPO
H PO
H
K
=
+
=
+
=
-
-
a
log [
]
[
]
.
log .
.
.
4
2
0
2
4
2 15
0 25
0 75
1 30
 
 Cathode: H PO
OH
HPO
H O
2
4
4
2
2
-
-
-
+
+

  


 



 
                    1.0 M    0.75 M       1.0 M
 
 Final        0.25 M         0            1.75 M
 
   P
P
HPO
H PO
H
K
=
+
=
+
=
-
-
a
log [
]
[
]
.
log .
.
.
4
2
0
2
4
2 15
1 75
0 25
3 0
\end{problembox}

\begin{problembox}
\textbf{Problem 540} \hfill \textbf{[Neeraj Kumar [NRJ] EXERCISE II (JEE ADVANCED) Q12]}
\label{q:ext-nrj_all_questions-540}

V, Fe and Hg will be oxidised by NO3
-
\end{problembox}

\begin{problembox}
\textbf{Problem 541} \hfill \textbf{[Neeraj Kumar [NRJ] EXERCISE II (JEE ADVANCED) Q13]}
\label{q:ext-nrj_all_questions-541}

Resistance and heat capacity depends on quantity.
\end{problembox}

\begin{problembox}
\textbf{Problem 542} \hfill \textbf{[Neeraj Kumar [NRJ] EXERCISE II (JEE ADVANCED) Q15]}
\label{q:ext-nrj_all_questions-542}

Net cell reaction of discharge is
 
 Pb + PbO2 + 2H2SO4 
 2PbSO4 + 2H2O
 
                          x mole                               x mole
 
 Initial mass of H2SO4,
 
  
 w1 = 1000 $/times$ 1.26 $/times$ 40
100
504
=
gm
 
 Final mass of H2SO4,
 
  
 w2 = (1260 -- 98x + 18x) $/times$ 28
100
 
  
      = (352.8 -- 22.4x) gm
 
 From reaction, 504 -- 98x = 352.8 -- 22.4 x $/implies$ x = 2
Section C (Comprehensions)
Comprehension I
\end{problembox}

\begin{problembox}
\textbf{Problem 543} \hfill \textbf{[Neeraj Kumar [NRJ] EXERCISE II (JEE ADVANCED) Q1]}
\label{q:ext-nrj_all_questions-543}

E$/circ$Cell = E
E
H
H
Zn
Zn
$/circ$
-
$/circ$
+
+
/
/
2
2
 because E H
H
$/circ$
+ /
2 was 
higher
 
 or, 0.76 = 1.00 -- E Zn
Zn
$/circ$
+
2 /
 $/implies$ E Zn
Zn
$/circ$
+
2 /
 = 0.24 V
\end{problembox}

\begin{problembox}
\textbf{Problem 544} \hfill \textbf{[Neeraj Kumar [NRJ] EXERCISE II (JEE ADVANCED) Q2]}
\label{q:ext-nrj_all_questions-544}

E$/circ$Cell = E
E
Cu
Cu
H
H
$/circ$
-
$/circ$
+
+
2
2
/
/
 because E Cu
Cu
$/circ$
+
2 /
 
was higher
 
 or, 0.34 = E Cu
Cu
$/circ$
+
2 /
 -- 1.00 $/implies$ E Cu
Cu
$/circ$
+
2 /
 = 1.34 V
\end{problembox}

\begin{problembox}
\textbf{Problem 545} \hfill \textbf{[Neeraj Kumar [NRJ] EXERCISE II (JEE ADVANCED) Q3]}
\label{q:ext-nrj_all_questions-545}

E$/circ$Cell = E
E
Cu
Cu
Zn
Zn
$/circ$
-
$/circ$
+
+
2
2
/
/
= 1.34 -- 0.24 = 1.10 V
 
 
Visit www.jeebooks.in for more
\end{problembox}

\begin{problembox}
\textbf{Problem 546} \hfill \textbf{[Neeraj Kumar [NRJ] EXERCISE II (JEE ADVANCED) Q4]}
\label{q:ext-nrj_all_questions-546}

E
E
Ag
K
Ag
Ag
Ag
Ag
sp
+
+
=
$/circ$
-
 
=
-
+
/
/
.
log
[
]
.
.
.log
0 06
1
1
0 80
0 06
1
1
 
                = + 0.314 V
\end{problembox}

\begin{problembox}
\textbf{Problem 547} \hfill \textbf{[Neeraj Kumar [NRJ] EXERCISE II (JEE ADVANCED) Q5]}
\label{q:ext-nrj_all_questions-547}

E
E
K
V
I |AgI|Ag
Ag
Ag
sp
$/circ$
=
$/circ$
-
 
= -
-
+ /
.
log
.
0 06
1
0 172
\end{problembox}

\begin{problembox}
\textbf{Problem 548} \hfill \textbf{[Neeraj Kumar [NRJ] EXERCISE II (JEE ADVANCED) Q6]}
\label{q:ext-nrj_all_questions-548}

E
E
I
I |AgI|Ag
I /AgI Ag
-
-
=
$/circ$
-
 
-
/
.
log[
]
0 06
1
 
                   = --0.172 -- 0.06   log 0.04 = --0.088 V
Comprehension III
\end{problembox}

\begin{problembox}
\textbf{Problem 549} \hfill \textbf{[Neeraj Kumar [NRJ] EXERCISE II (JEE ADVANCED) Q7]}
\label{q:ext-nrj_all_questions-549}

As(0.40) > (--0.87), reduction of Ni2O3(s) will 
occur.
\end{problembox}

\begin{problembox}
\textbf{Problem 550} \hfill \textbf{[Neeraj Kumar [NRJ] EXERCISE II (JEE ADVANCED) Q8]}
\label{q:ext-nrj_all_questions-550}

E$/circ$cell = (0.40) -- (--0.87) = 1.27 V
\end{problembox}

\begin{problembox}
\textbf{Problem 551} \hfill \textbf{[Neeraj Kumar [NRJ] EXERCISE II (JEE ADVANCED) Q9]}
\label{q:ext-nrj_all_questions-551}

Net cell reaction is independent from OH-- (aq)
\end{problembox}

\begin{problembox}
\textbf{Problem 552} \hfill \textbf{[Neeraj Kumar [NRJ] EXERCISE II (JEE ADVANCED) Q10]}
\label{q:ext-nrj_all_questions-552}

--$/Delta$G$/circ$ = nFE$/circ$cell = 2 $/times$ 96500 $/times$ 1.27 = 245110 J
Comprehension IV
\end{problembox}

\begin{problembox}
\textbf{Problem 553} \hfill \textbf{[Neeraj Kumar [NRJ] EXERCISE II (JEE ADVANCED) Q11]}
\label{q:ext-nrj_all_questions-553}

H3O+ is only needed in balancing cathode reaction:
 
 
NO
H O
e
HNO
H O
3
3
2
2
3
2
4
-
+
-
+
+
 $/to$
 
+
\end{problembox}

\begin{problembox}
\textbf{Problem 554} \hfill \textbf{[Neeraj Kumar [NRJ] EXERCISE II (JEE ADVANCED) Q12]}
\label{q:ext-nrj_all_questions-554}

Moles of electron needed = 2 $/times$ moles of HNO2 
formed
\end{problembox}

\begin{problembox}
\textbf{Problem 555} \hfill \textbf{[Neeraj Kumar [NRJ] EXERCISE II (JEE ADVANCED) Q13]}
\label{q:ext-nrj_all_questions-555}

neqHNO
Q
F
2 =
 $/implies$ 0 1
2
10
96500
. $/times$
=
$/times$ t $/implies$ t = 1930 sec
Comprehension V
\end{problembox}

\begin{problembox}
\textbf{Problem 556} \hfill \textbf{[Neeraj Kumar [NRJ] EXERCISE II (JEE ADVANCED) Q14]}
\label{q:ext-nrj_all_questions-556}

$/Delta$G$/circ$ = --nFE$/circ$ $/implies$ --237.39 $/times$ 103 = --2 $/times$ 96500 $/times$ E$/circ$cell
 
   E$/circ$cell = 1.23 V
\end{problembox}

\begin{problembox}
\textbf{Problem 557} \hfill \textbf{[Neeraj Kumar [NRJ] EXERCISE II (JEE ADVANCED) Q15]}
\label{q:ext-nrj_all_questions-557}

Moles of H2 needed = 23 739
237 39
0 1
.
.
.
=
 
   Volume of H2 needed = 0.1 $/times$ 22.7 = 2.27 L
\end{problembox}

\begin{problembox}
\textbf{Problem 558} \hfill \textbf{[Neeraj Kumar [NRJ] EXERCISE II (JEE ADVANCED) Q16]}
\label{q:ext-nrj_all_questions-558}

Ecell is independent from [OH--]
\end{problembox}

\begin{problembox}
\textbf{Problem 559} \hfill \textbf{[Neeraj Kumar [NRJ] EXERCISE II (JEE ADVANCED) Q17]}
\label{q:ext-nrj_all_questions-559}

$/Delta$ $/circ$ = $/Delta$
$/circ$ -$/Delta$
$/circ$
= -
$/times$
--
$/times$
S
H
G
T
(
.
)
(
.
)
285 8 10
237 39 10
298
3
3
 
         = --162.4 J/K--1
\end{problembox}

\begin{problembox}
\textbf{Problem 560} \hfill \textbf{[Neeraj Kumar [NRJ] EXERCISE II (JEE ADVANCED) Q18]}
\label{q:ext-nrj_all_questions-560}

  = -$/Delta$
$/circ$
-$/Delta$
$/circ$ =
=
(
)
(
)
.
.
.
G
H
237 39
285 8
0 8306 or 83.06/%
Comprehension VI
\end{problembox}

\begin{problembox}
\textbf{Problem 561} \hfill \textbf{[Neeraj Kumar [NRJ] EXERCISE II (JEE ADVANCED) Q19]}
\label{q:ext-nrj_all_questions-561}

Refer theory given is passage.
\end{problembox}

\begin{problembox}
\textbf{Problem 562} \hfill \textbf{[Neeraj Kumar [NRJ] EXERCISE II (JEE ADVANCED) Q20]}
\label{q:ext-nrj_all_questions-562}

E
E
E
V
O
H O H
Fe
Fe
$/circ$ =
$/circ$
-
$/circ$
=
--
=
+
+
2
2
2
1 229
0 447
1 676
/
,
/
.
(
.
)
.
\end{problembox}

\begin{problembox}
\textbf{Problem 563} \hfill \textbf{[Neeraj Kumar [NRJ] EXERCISE II (JEE ADVANCED) Q21]}
\label{q:ext-nrj_all_questions-563}

E
E
Cu
Cu
Pb
Pb
$/circ$
>
$/circ$
+
+
2
2
/
/
 and hence Cu will not 
oxidise easily.
\end{problembox}

\begin{problembox}
\textbf{Problem 564} \hfill \textbf{[Neeraj Kumar [NRJ] EXERCISE II (JEE ADVANCED) Q22]}
\label{q:ext-nrj_all_questions-564}

neq Fe
Q
F
=
 $/implies$ w
56
2
0 5 1 0
3600
96500
$/times$
=
$/times$
$/times$
.
.
 
 
 $/implies$ w = 0.522 gm
\end{problembox}

\begin{problembox}
\textbf{Problem 565} \hfill \textbf{[Neeraj Kumar [NRJ] EXERCISE II (JEE ADVANCED) Q23]}
\label{q:ext-nrj_all_questions-565}

Theoretical
Comprehension VII
\end{problembox}

\begin{problembox}
\textbf{Problem 566} \hfill \textbf{[Neeraj Kumar [NRJ] EXERCISE II (JEE ADVANCED) Q24]}
\label{q:ext-nrj_all_questions-566}

neq Ni
Q
F
=
 $/implies$ w
58 7
2
15
3600
0 6
96500
.
.
$/times$
=
$/times$
$/times$
 
 
 $/implies$ w = 9.85 gm
\end{problembox}

\begin{problembox}
\textbf{Problem 567} \hfill \textbf{[Neeraj Kumar [NRJ] EXERCISE II (JEE ADVANCED) Q25]}
\label{q:ext-nrj_all_questions-567}

V = A $/times$ t $/implies$ 9 85
8 9
4 0
2
.
.
( .
)
=
$/times$
$/times$t $/implies$ t = 0.138 cm
 
 
Visit www.jeebooks.in for more
\end{problembox}

\begin{problembox}
\textbf{Problem 568} \hfill \textbf{[Neeraj Kumar [NRJ] EXERCISE II (JEE ADVANCED) Q26]}
\label{q:ext-nrj_all_questions-568}

neq H
Q
F
2 =
 $/implies$ 
VH2
22 4
2
15
3600
0 4
96500
.
.
$/times$
=
$/times$
$/times$
 
 
 $/implies$ V
L
H2
2 5
= .
\end{problembox}

\begin{problembox}
\textbf{Problem 569} \hfill \textbf{[Neeraj Kumar [NRJ] EXERCISE II (JEE ADVANCED) Q27]}
\label{q:ext-nrj_all_questions-569}

Anode produced is only O2 gas.
 
 
neq O
Q
F
2 =
 $/implies$ w
8
15
3600
96500
=
$/times$
 $/implies$ w = 4.477 gm
Comprehension VIII
\end{problembox}

\begin{problembox}
\textbf{Problem 570} \hfill \textbf{[Neeraj Kumar [NRJ] EXERCISE II (JEE ADVANCED) Q28]}
\label{q:ext-nrj_all_questions-570}

E$/circ$cell = 0.8 -- 0.05 = 0.75 V
 
 Now, E
RT
nF
K
cell
$/circ$
=
 ln
 $/implies$ 0 75
1
2
38 92
.
.
ln
=
$/times$
 
K 
$/implies$ ln K = 58.38
\end{problembox}

\begin{problembox}
\textbf{Problem 571} \hfill \textbf{[Neeraj Kumar [NRJ] EXERCISE II (JEE ADVANCED) Q29]}
\label{q:ext-nrj_all_questions-571}

The oxidation reaction of glucose contains H+ 
and
 
 
E
E
H
=
$/circ$ -
 
+
0 0592
2
2
.
log[
] 
 
 $/implies$ E -- E$/circ$ = 0.0592 PH = 0.6512 V
\end{problembox}

\begin{problembox}
\textbf{Problem 572} \hfill \textbf{[Neeraj Kumar [NRJ] EXERCISE II (JEE ADVANCED) Q30]}
\label{q:ext-nrj_all_questions-572}

Standard potential is independent from ammonia 
concentration.
Comprehension IX
\end{problembox}

\begin{problembox}
\textbf{Problem 573} \hfill \textbf{[Neeraj Kumar [NRJ] EXERCISE II (JEE ADVANCED) Q31]}
\label{q:ext-nrj_all_questions-573}

Left electron: H2(g) 
 2H+(aq) + 2e--
 
 Right electrode: 2AgCl(s) + 2e-- 
 2Ag(s) + 
2Cl-- (aq)
 
    Net reaction: H2(g) + 2AgCl(s) 
 2Ag(s) 
+ 2H+(aq) + 2Cl--(aq)
\end{problembox}

\begin{problembox}
\textbf{Problem 574} \hfill \textbf{[Neeraj Kumar [NRJ] EXERCISE II (JEE ADVANCED) Q32]}
\label{q:ext-nrj_all_questions-574}

$/Delta$ $/circ$ =
 
-
-
 
  
 
  =
$/times$
$/times$
-
= -
S
nF
E
E
T
T
J/K
2
1
2
1
2
96500
0 21
0 23
20
193
.
.
 
 $/Delta$G$/circ$ = --nFE$/circ$ = --2 $/times$ 96500 $/times$ 0.23 = -- 44390 J at 
15$/circ$C
 
 Now, $/Delta$H$/circ$ = $/Delta$G$/circ$ + T   $/Delta$S$/circ$ = (--44390) + 288 $/times$ 
(--193)
 
 = --99974 J = --49987 J/mole AgCl
\end{problembox}

\begin{problembox}
\textbf{Problem 575} \hfill \textbf{[Neeraj Kumar [NRJ] EXERCISE II (JEE ADVANCED) Q33]}
\label{q:ext-nrj_all_questions-575}

$/Delta$S$/circ$ = --193 J/K = --96.5 J/K per mole AgCl
\end{problembox}

\begin{problembox}
\textbf{Problem 576} \hfill \textbf{[Neeraj Kumar [NRJ] EXERCISE II (JEE ADVANCED) Q34]}
\label{q:ext-nrj_all_questions-576}

$/Delta$G$/circ$298 = $/Delta$H$/circ$ -- T   $/Delta$S$/circ$ = (--49987) -- 298 $/times$ (--96.5)
 
 or, --1 $/times$ 96500 $/times$ E$/circ$ = -- 21230
 
  $/implies$ E$/circ$ = 0.22 V = E Cl
AgCl/Ag
$/circ$
-/
 
 or, E
E
K
Cl
AgCl/Ag
Ag
Ag
sp
$/circ$
=
$/circ$
+
-
+
/
/
.
log
0 058
1
 
 or, 0 22
0 80
0 058
1
.
.
.
log
=
+
Ksp $/implies$ Ksp = 1 $/times$ 10--10
 
 Hence, Solubility, S = 
K
M
sp =
-
10 5
Comprehension X
\end{problembox}

\begin{problembox}
\textbf{Problem 577} \hfill \textbf{[Neeraj Kumar [NRJ] EXERCISE II (JEE ADVANCED) Q35]}
\label{q:ext-nrj_all_questions-577}

Add Ag(s) in both sides of given reaction to get 
cell reaction. Now, for E$/circ$cell
 
 $/Delta$G$/circ$ = --nFE$/circ$ $/implies$ [--109] -- [77 + (--129)] $/times$ 103
 
 
           = --1 $/times$ 96500 $/times$ E$/circ$cell
 
   E$/circ$ cell = 0.59 V
\end{problembox}

\begin{problembox}
\textbf{Problem 578} \hfill \textbf{[Neeraj Kumar [NRJ] EXERCISE II (JEE ADVANCED) Q36]}
\label{q:ext-nrj_all_questions-578}

E
K
cell
eq
$/circ$
= 0 059
.
log
n
 $/implies$ 0 59
0 059
1
1
.
.
log
=
 
Ksp
 
   Ksp = 10--10
\end{problembox}

\begin{problembox}
\textbf{Problem 579} \hfill \textbf{[Neeraj Kumar [NRJ] EXERCISE II (JEE ADVANCED) Q37]}
\label{q:ext-nrj_all_questions-579}

Zn(s) + 2Ag+(aq)   Zn2+(aq) + 3Ag(s); E$/circ$ = 0.80 -- 
(--0.76)
 
 Now, E
Keq
$/circ$ =
 
0 059
.
log
n
 
 
 $/implies$ 1 56
0 059
2
2
2
.
.
log [
]
[
]
=
 
+
+
Zn
Ag
 
   log [
]
[
]
.
Zn
Ag
2
2
52 88
+
+
=
\end{problembox}

\begin{problembox}
\textbf{Problem 580} \hfill \textbf{[Neeraj Kumar [NRJ] EXERCISE II (JEE ADVANCED) Q38]}
\label{q:ext-nrj_all_questions-580}

Moles of Zn added = 6 539 10
65 39
10
2
3
.
.
$/times$
=
-
-
 
 Moles of Ag+ present = 10
1000
100
10
5
6
-
-
$/times$
=
(L.R.)
 
   Moles of Ag precipitated = 10--6
 
 
Visit www.jeebooks.in for more
\end{problembox}

\begin{problembox}
\textbf{Problem 581} \hfill \textbf{[Neeraj Kumar [NRJ] EXERCISE II (JEE ADVANCED) Q39]}
\label{q:ext-nrj_all_questions-581}

For KCl:  
=
=
 
 
m
C
G G
C
$/kappa$
 $/implies$ G* =  m   C/G
 
 or, G* =  m   C   R = 200 $/times$ (0.02 $/times$ 10--3) $/times$ 100 = 0.4 
cm--1
\end{problembox}

\begin{problembox}
\textbf{Problem 582} \hfill \textbf{[Neeraj Kumar [NRJ] EXERCISE II (JEE ADVANCED) Q40]}
\label{q:ext-nrj_all_questions-582}

$/kappa$water
G G
cm
=
 
=
$/times$
=
$/times$
$/Omega$
 
-
-
-
1
10000
0 4
4 10 5
1
1
.
\end{problembox}

\begin{problembox}
\textbf{Problem 583} \hfill \textbf{[Neeraj Kumar [NRJ] EXERCISE II (JEE ADVANCED) Q41]}
\label{q:ext-nrj_all_questions-583}

For NaCl:  
=
=
 
 
m
C
G G
C
$/kappa$
 
 
 or, 125
1
8000
1
10000
0 4
585 58 5
=
-
 
  
 
  $/times$
 
  
 
  
.
/
.
V
 
 
 $/implies$ V = 1.25 $/times$ 108 cm3
 
 = 1.25 $/times$ 105 L
Comprehension XII
\end{problembox}

\begin{problembox}
\textbf{Problem 584} \hfill \textbf{[Neeraj Kumar [NRJ] EXERCISE II (JEE ADVANCED) Q42]}
\label{q:ext-nrj_all_questions-584}

E
E
Cl
Cl
I
$/circ$
>
$/circ$
-
-
2
2
/
/I
\end{problembox}

\begin{problembox}
\textbf{Problem 585} \hfill \textbf{[Neeraj Kumar [NRJ] EXERCISE II (JEE ADVANCED) Q43]}
\label{q:ext-nrj_all_questions-585}

E
E
Mn
Mn
O
H O, H
$/circ$
>
$/circ$
+
+
+
3
2
2
2
/
/
Comprehension XIII
\end{problembox}

\begin{problembox}
\textbf{Problem 586} \hfill \textbf{[Neeraj Kumar [NRJ] EXERCISE II (JEE ADVANCED) Q44]}
\label{q:ext-nrj_all_questions-586}

Net reaction: M+(1M) 
 M+ (0.05 M); 
 
                       Higher Conc.      Lower Conc.
 
 Ecell > 0, $/Delta$Gcell < 0
\end{problembox}

\begin{problembox}
\textbf{Problem 587} \hfill \textbf{[Neeraj Kumar [NRJ] EXERCISE II (JEE ADVANCED) Q45]}
\label{q:ext-nrj_all_questions-587}

70 mV = E
RT
F
$/circ$ -
 ln .0 05
1
 
(1)
 
 
E
E
RT
F
E
RT
F
req =
$/circ$ -
 
=
$/circ$ -
 
ln .
ln ( .
)
0 0025
1
0 05
1
2
 
 
(2)
 
 and E$/circ$ = 0.  Hence, Ereq = 140 mV
Comprehension XIV
\end{problembox}

\begin{problembox}
\textbf{Problem 588} \hfill \textbf{[Neeraj Kumar [NRJ] EXERCISE II (JEE ADVANCED) Q46]}
\label{q:ext-nrj_all_questions-588}

$/Delta$Gcell = --nFEcell = --2 $/times$ 96500 $/times$ 0.059 = --11387 J
\end{problembox}

\begin{problembox}
\textbf{Problem 589} \hfill \textbf{[Neeraj Kumar [NRJ] EXERCISE II (JEE ADVANCED) Q47]}
\label{q:ext-nrj_all_questions-589}

E
E
M
M
cell
cell
Left
Right
=
$/circ$
-
 
+
+
0 059
2
2
2
.
log [
]
[
]
 
 or, 0 059
0
0 059
2
4
0 001
1 3
.
.
log
(
/ )
.
/
=
-
 
Ksp
 
 
 $/implies$ Ksp = 4 $/times$ 10--15
Comprehension XV
\end{problembox}

\begin{problembox}
\textbf{Problem 590} \hfill \textbf{[Neeraj Kumar [NRJ] EXERCISE II (JEE ADVANCED) Q48]}
\label{q:ext-nrj_all_questions-590}

H+ + Cl-- + (NaOH) 
 Na+ + Cl-- + H2O
 
 Conductance   rst decreases and H+ ions are 
replaced by Na+ ions. After equivalent point, 
conductance increase due to increase in number 
of ions.
\end{problembox}

\begin{problembox}
\textbf{Problem 591} \hfill \textbf{[Neeraj Kumar [NRJ] EXERCISE II (JEE ADVANCED) Q49]}
\label{q:ext-nrj_all_questions-591}

CH3COOH + (NaOH added) 
 CH3COO-- 
+ Na+ + H2O 
 
 As number of ions increases, conductance 
increases. slight decrease initially was due to some 
dissociated CH3COOH. After equivalence point, 
in place of CH3COO-- ion, number of OH-- ions 
increases and hence slope becomes greater.
\end{problembox}

\begin{problembox}
\textbf{Problem 592} \hfill \textbf{[Neeraj Kumar [NRJ] EXERCISE II (JEE ADVANCED) Q50]}
\label{q:ext-nrj_all_questions-592}

H+ + Cl-- + (NH4OH added) 
 
  NH
Cl
H O
4
2
+
-
+
+
 
 As H+ ions are replaced by 
NH4
+ ions, 
conductance decreases. After equivalent point, 
it become almost constant as the dissociation of 
added NH4OH will be suppressed in presence of 
NH4
+ ions.
\end{problembox}

\begin{problembox}
\textbf{Problem 593} \hfill \textbf{[Neeraj Kumar [NRJ] EXERCISE II (JEE ADVANCED) Q51]}
\label{q:ext-nrj_all_questions-593}

HCl will neutralize   rst followed by CH3COOH.
\end{problembox}

\begin{problembox}
\textbf{Problem 594} \hfill \textbf{[Neeraj Kumar [NRJ] EXERCISE II (JEE ADVANCED) Q52]}
\label{q:ext-nrj_all_questions-594}

Ionic mobilities of Ag+ and K+ ions do not di  er 
largely
\end{problembox}

\begin{problembox}
\textbf{Problem 595} \hfill \textbf{[Neeraj Kumar [NRJ] EXERCISE II (JEE ADVANCED) Q53]}
\label{q:ext-nrj_all_questions-595}

First number of ions increases and then remain 
atmost constant.
added
 
 
Visit www.jeebooks.in for more
\end{problembox}

\begin{problembox}
\textbf{Problem 596} \hfill \textbf{[Neeraj Kumar [NRJ] EXERCISE II (JEE ADVANCED) Q1]}
\label{q:ext-nrj_all_questions-596}

CuCl2 
 Cu + Cl2
\end{problembox}

\begin{problembox}
\textbf{Problem 597} \hfill \textbf{[Neeraj Kumar [NRJ] EXERCISE II (JEE ADVANCED) Q2]}
\label{q:ext-nrj_all_questions-597}

E
E
informative
Zn
Zn
Ag
Ag
$/circ$
<
$/circ$
+
+
2 /
/
(
)
\end{problembox}

\begin{problembox}
\textbf{Problem 598} \hfill \textbf{[Neeraj Kumar [NRJ] EXERCISE II (JEE ADVANCED) Q3]}
\label{q:ext-nrj_all_questions-598}

Reason is di  erent charges on ions.
\end{problembox}

\begin{problembox}
\textbf{Problem 599} \hfill \textbf{[Neeraj Kumar [NRJ] EXERCISE II (JEE ADVANCED) Q5]}
\label{q:ext-nrj_all_questions-599}

Negative reduction potential means greater 
tendency to get oxidised.
\end{problembox}

\begin{problembox}
\textbf{Problem 600} \hfill \textbf{[Neeraj Kumar [NRJ] EXERCISE II (JEE ADVANCED) Q6]}
\label{q:ext-nrj_all_questions-600}

Number of ions increases considerably only for 
weak electrolytes.
\end{problembox}

\begin{problembox}
\textbf{Problem 601} \hfill \textbf{[Neeraj Kumar [NRJ] EXERCISE II (JEE ADVANCED) Q9]}
\label{q:ext-nrj_all_questions-601}

It is due to very high over voltage potential of 
hydrogen at mercury cathode.
\end{problembox}

\begin{problembox}
\textbf{Problem 602} \hfill \textbf{[Neeraj Kumar [NRJ] EXERCISE II (JEE ADVANCED) Q12]}
\label{q:ext-nrj_all_questions-602}

E
E
RT
nF
Q
=
-
 
$/circ$
log
\end{problembox}

\begin{problembox}
\textbf{Problem 603} \hfill \textbf{[Neeraj Kumar [NRJ] EXERCISE II (JEE ADVANCED) Q13]}
\label{q:ext-nrj_all_questions-603}

Cl-- will combine with Ag+ to precipitate AgCl.
\end{problembox}

\begin{problembox}
\textbf{Problem 604} \hfill \textbf{[Neeraj Kumar [NRJ] EXERCISE II (JEE ADVANCED) Q14]}
\label{q:ext-nrj_all_questions-604}

At Cathode: 
 
 
2H O(l) + 2e
H (g) + 2OH
aq
2
2
-
-
 $/to$
 
(
) 
 
 At Anode: 
 
 
2
2
2
3
2
6
2
CH COO
aq
C H
CO
e
-
-
(
)
( )
( )
 $/to$
 
+
+
g
g
\end{problembox}

\begin{problembox}
\textbf{Problem 605} \hfill \textbf{[Neeraj Kumar [NRJ] EXERCISE II (JEE ADVANCED) Q15]}
\label{q:ext-nrj_all_questions-605}

Theoretical
Section E (Column Match)
\end{problembox}

\begin{problembox}
\textbf{Problem 606} \hfill \textbf{[Neeraj Kumar [NRJ] EXERCISE II (JEE ADVANCED) Q3]}
\label{q:ext-nrj_all_questions-606}

E
E
E
V
Fe
Fe
Fe
Fe
Fe
Fe
3
3
2
2
1
2
1
2
0 037
+
+
+
+
$/circ$
$/circ$
$/circ$
=
$/times$
+
$/times$
+
= -
|
,
|
.
 
 
 
E
V
H
( H O
OH )
.
.
.
4
4
4
2
0 40
1 23
0 83
$/to$
+
$/circ$
+
-=
-
= -
 
 
 
E
E
E
Cu
Cu
Cu
Cu
(Cu
Cu
)
| Cu
|
.
.
.
.
2
2
2
0 34
0 52
2
1
0 52
0
+
+
+
+
+
+
$/to$
$/circ$
$/circ$
$/circ$
=
-
=
-
-
-
= -
70 V
 
 
 
E
V
Cr
Cr
3
2
3
0 74
2
0 91
3
2
0 4
+
+
$/circ$
=
$/times$ -
-
$/times$ -
-
= -
,
(
.
)
(
.
)
.
\end{problembox}

\begin{problembox}
\textbf{Problem 607} \hfill \textbf{[Neeraj Kumar [NRJ] EXERCISE II (JEE ADVANCED) Q4]}
\label{q:ext-nrj_all_questions-607}

For concentration cell, both half cell must have 
same con  guration.
 
 P. Ag
Cl
AgCl sponteneous,
K
K
eq
sp
+
-
+
 $/to$
 
=
>>
1
1
 
 
 Q. Ag Br + Cl
AgCl
Br
Cl
Non-Spontene
eq
sp
sp
-
-
 $/to$
 
+
=
<<
;
(Ag Br)
(Ag
)
,
K
K
K
1
ous
 
 
 R. Ag
Ag
+
+
 $/to$
 
( . M)
( . M)
1 0
0 1
 
 
 Higher to lower concentration, spontaneous
 
 S. Cl
Cl
-
-
 $/to$
 
-
( . M)
( . M)
Non spontaneous
0 1
1 0
 
Section F (Subjective)
Single-digit Integer Type
\end{problembox}

\begin{problembox}
\textbf{Problem 608} \hfill \textbf{[Neeraj Kumar [NRJ] EXERCISE II (JEE ADVANCED) Q1]}
\label{q:ext-nrj_all_questions-608}

E = E
PH
$/circ$
 
+
$/to$
+
+
-
-0 06
2
2
2
2
2
.
log
[H ]
(assuming
H
e
H )
n
 
 
 
-
=
-
 
$/implies$
=
+
+
-
0 18
0
0 06
2
1
10
2
3
.
.
log
[H ]
[H ]
M 
 
 
C H NH
H O
C H NH
H O
M
6
5
3
2
6
5
2
3
+
-
+
+
+
(
)M
C x
x

  


 



 
 
 
h
x
c
=
=
=
-
10
0 04
4
3
140
.
/%
or
 
 
 
Visit www.jeebooks.in for more
\end{problembox}

\begin{problembox}
\textbf{Problem 609} \hfill \textbf{[Neeraj Kumar [NRJ] EXERCISE II (JEE ADVANCED) Q2]}
\label{q:ext-nrj_all_questions-609}

E
E
=
-
 
$/circ$
+
0 06
2
1
.
log
[Cu ]
 
 
 
0 31
0 34
0 06
2
1
0 1
2
2
.
.
.
log
[Cu
]
[Cu
]
. M
=
-
 
$/implies$
=
+
+
 
 
   [OH ]
.
-
+
-
-
=
=
=
$/implies$
=
Ksp
H
Cu
p
2
19
9
10
0 1
10
5
\end{problembox}

\begin{problembox}
\textbf{Problem 610} \hfill \textbf{[Neeraj Kumar [NRJ] EXERCISE II (JEE ADVANCED) Q3]}
\label{q:ext-nrj_all_questions-610}

E
E
E
Fe
Fe
Fe
Fe
Fe
Fe
3
2
3
2
3
2
3
2
3
0 04
2
0 44
+
+
+
+
$/circ$
$/circ$
$/circ$
=
$/times$
-$/times$
-
=
$/times$ -
-$/times$ -
|
|
|
(
.
)
(
.
)
1
0 76
= .
V
 
 
 Now, E
E
Fe
Fe
Fe
Fe
3
2
3
2
0 06
1
2
3
+
+
+
+
=
-
$/circ$
+
+
|
|
.
log [Fe
]
[Fe
]
 
 
 or, 0 7 8
0 76
0 06
5
2
3
2
3
.
.
.
log [Fe
]
[Fe
]
[Fe
]
[Fe
]
1 =
-
$/implies$
=
+
+
+
+
\end{problembox}

\begin{problembox}
\textbf{Problem 611} \hfill \textbf{[Neeraj Kumar [NRJ] EXERCISE II (JEE ADVANCED) Q4]}
\label{q:ext-nrj_all_questions-611}

2Fe3+    +    2I--   $/to$   2Fe2+   +    I2;  
 
                  0.5 M         excess
 
 100/%           0            1.0 M          0.5 M
 
 Equ.           CM         1.0 M         0.5 M
 
 
Ecell
V;
$/circ$
=
-
=
0 77
0 53
0 24
.
.
.
 
 
 Keq = 108
 
 Now,  10
0 5
1 0
5 10
8
2
2
2
5
=
$/times$
$/implies$
=
$/times$
-
( . )
( . )
C
C
M
\end{problembox}

\begin{problembox}
\textbf{Problem 612} \hfill \textbf{[Neeraj Kumar [NRJ] EXERCISE II (JEE ADVANCED) Q5]}
\label{q:ext-nrj_all_questions-612}

Assuming the cell as concentration cell, net cell 
reaction is
 
 
Ag
C
Ag
C
sp
sp(AgI)
+
-
+
=
 
  
 
   $/to$
 
=
C
K
K
1
2
(AgCl)
[
l ]
(
) 
 
 and E
C
C
cell
O
=
-0 06
1
2
1
.
log
 
 
 or, 0.102 
 
 = -- 0.06   log
.
.
[Cl ]
[Cl ]
8 1 10
1 8 10
4
10
17
10
4
$/times$
$/times$
 
  
 
  
$/implies$
=
$/times$
-
-
-
-
-M
\end{problembox}

\begin{problembox}
\textbf{Problem 613} \hfill \textbf{[Neeraj Kumar [NRJ] EXERCISE II (JEE ADVANCED) Q6]}
\label{q:ext-nrj_all_questions-613}

n
n
eq
eq
Pb
Tl =
+
Q
F 
 
 
$/implies$
$/times$
$/times$
+
$/times$
$/times$ =
$/times$
5 0
70
100
208
2
5 0
30
100
204
1
1 1
96500
.
.
.
t 
 
   t = 3597.4 sec ; 1 hr
\end{problembox}

\begin{problembox}
\textbf{Problem 614} \hfill \textbf{[Neeraj Kumar [NRJ] EXERCISE II (JEE ADVANCED) Q7]}
\label{q:ext-nrj_all_questions-614}

n
Q
F
x
x
eq
rI =
$/implies$
$/times$
=
$/times$
$/times$
$/implies$
=
0 36
192
0 075
2
3600
96500
3
.
.
 
 
 Now, x + 6(--1) = y $/implies$ y = --3
\end{problembox}

\begin{problembox}
\textbf{Problem 615} \hfill \textbf{[Neeraj Kumar [NRJ] EXERCISE II (JEE ADVANCED) Q8]}
\label{q:ext-nrj_all_questions-615}

2
2
2
2
2
NaCl + 2H O
NaOH
H
Cl
electrolysis
 
$/to$
    
+
+
 
 
 
2
2
2
NaOH + Cl
NaCl
NaClO
H O
 $/to$
 
+
+
 
 
 nNaClO formed = nCl2 produced from electrolysis
 
 or, 
(
. )
.
.
.
10
10
1 0
7 45
100
74 5
2
2 5
96500
3
$/times$
$/times$
$/times$
$/times$
=
$/times$ t 
 
 (n-factor of Cl2 in electrolysis)
 
   t = 7.72 $/times$ 105 sec = 8.93 days $/approx$ 9 days
\end{problembox}

\begin{problembox}
\textbf{Problem 616} \hfill \textbf{[Neeraj Kumar [NRJ] EXERCISE II (JEE ADVANCED) Q9]}
\label{q:ext-nrj_all_questions-616}

Cathode: Cu2+ + 2e-- 
 Cu
 
 Anode 2H2O 
O2 + 4H+ + 4e--
 
 Moles of e-- used = Q
F =
$/times$
$/times$
$/times$
-
0 161 5
60
96500
5 10 4
.

 
 Eq. of Cu2+ present = 500
0 1
1000
2
0 1
$/times$
$/times$
=
.
. (
)
excess
 
   Moles of H+ produced = Moles of e-- = 5 $/times$ 10--4
 
   [H+]  nal = 5 10
500
1000
10
4
3
$/times$
$/times$
=
-
-M $/implies$ PH = 3.0
\end{problembox}

\begin{problembox}
\textbf{Problem 617} \hfill \textbf{[Neeraj Kumar [NRJ] EXERCISE II (JEE ADVANCED) Q10]}
\label{q:ext-nrj_all_questions-617}

cathode:       Cu2+        +      2e-- 
 Cu
 
                 
250
0 1
1000
$/times$ .           5 1351
96500
$/times$
 
                = 0.025 mole      = 0.07 mole
 
 As Cu+ will not remain   nally in solution, no 
complex formation.
\end{problembox}

\begin{problembox}
\textbf{Problem 618} \hfill \textbf{[Neeraj Kumar [NRJ] EXERCISE II (JEE ADVANCED) Q11]}
\label{q:ext-nrj_all_questions-618}

neq metal = neq Cl2 $/implies$ 52 8
9 08
22 7
2
.
.
.
E
=
$/times$
 $/implies$ E = 66
 
 At. wt. (approx) = 6 4
0 032
200
.
.
=
 
   valency
At wt
Eq.wt
=
=
.
.
200
66
3

 
 
Visit www.jeebooks.in for more
\end{problembox}

\begin{problembox}
\textbf{Problem 619} \hfill \textbf{[Neeraj Kumar [NRJ] EXERCISE II (JEE ADVANCED) Q13]}
\label{q:ext-nrj_all_questions-619}

Initial mass of H2SO4, 
 
 w1 = 2000
1 1
16
100
352
$/times$
$/times$
=
.
gm 
 
 Final mass of H2SO4,
 
 w2 = 2000
1 42
40
100
1136
$/times$
$/times$
=
.
gm
 
 Now, neq H2SO4 produced = Q
F
 
 or, (
)
1136
352
98
1
965
9
3600
96500
-
$/times$ =
$/times$
$/times$
 
  
 
  
i
 $/implies$ i = 2A
\end{problembox}

\begin{problembox}
\textbf{Problem 620} \hfill \textbf{[Neeraj Kumar [NRJ] EXERCISE II (JEE ADVANCED) Q13]}
\label{q:ext-nrj_all_questions-620}

Ionic mobility, $/mu$
$/lambda$
=
$/circ$
=
m
F
Speed of ion
Potential gradient
 
 or, 7 5 10
96500
1 93 0 12
3
.
/
.
/ .
$/times$
=
$/times$
-
distance 20
3600 
 
 $/implies$ distance = 0.09 m = 9 cm
\end{problembox}

\begin{problembox}
\textbf{Problem 621} \hfill \textbf{[Neeraj Kumar [NRJ] EXERCISE II (JEE ADVANCED) Q14]}
\label{q:ext-nrj_all_questions-621}

neq Ag oxidised = Q
F $/implies$ w
108
1
9 65 1 3600
96500
$/times$ =
$/times$ $/times$
.
 
   w = 38.88 gm
 
   Mass of anode dissolved = 38 88
100
60
64 8
.
.
$/times$
$/times$
gm
 
   Final mass of anode = 72.8 -- 64.8 = 8 gm
\end{problembox}

\begin{problembox}
\textbf{Problem 622} \hfill \textbf{[Neeraj Kumar [NRJ] EXERCISE II (JEE ADVANCED) Q15]}
\label{q:ext-nrj_all_questions-622}

Cathode: 2H2O(l) + 2e-- 
 H2(g) + 2OH--(aq)
 
 Anode: 2ce--(aq) 
 Cl2(g) + 2e--
 
 neq OH-- produced = Q
F 
 
 $/implies$ nOH-$/times$ =
$/times$
=
-
1
9 65 10
96500
10 3
.
 
   [OH--]  nal = 10
100
10
3
5
-
-
=
M $/implies$ PH = 9.0
Four digit integer type
\end{problembox}

\begin{problembox}
\textbf{Problem 623} \hfill \textbf{[Neeraj Kumar [NRJ] EXERCISE II (JEE ADVANCED) Q1]}
\label{q:ext-nrj_all_questions-623}

$/Lambda$$/circ$
=
$/circ$
+ $/circ$
-
m
m
m
(AgCl)
(AgCl)
(Cl )
$/lambda$
$/lambda$
 
 
                 
=
$/times$
+
$/times$
=
$/times$
-
-
-
-
-
6 19 10
7 81 10
14 00 10
3
3
3
1
2
1
.
.
.
$/Omega$ m mol
 
 
 Now, $/Lambda$m
C
S
=
$/implies$
$/times$
=
$/times$
-
-
$/kappa$
14 10
2 8 10
3
4
.
 
 
 $/implies$ S = 0.02 mol m--3 
 
         = 2 $/times$ 10--5 M
 
         = 287 $/times$ 10--5 g/L
\end{problembox}

\begin{problembox}
\textbf{Problem 624} \hfill \textbf{[Neeraj Kumar [NRJ] EXERCISE II (JEE ADVANCED) Q2]}
\label{q:ext-nrj_all_questions-624}

[S
]
K
[H S]
[
]
.
(
)
2
1
2
2
2
8
13
3
16
10
10
0 1
10
10
-
+
-
-
-
-
=
 
 
=
$/times$
$/times$
=
a
a
K
H
M
 
 
   [Ag ]
[S
]
+
-
-
-
-
=
=
$/times$
=
$/times$
K
M
sp
2
48
16
16
4
10
10
2
10
 
 
 Now, EAg /Ag
+
=
-
$/times$
-
0 80
0 06
1
1
2
10 16
.
.
log
 = --0.142 V
 
   Required potential = 142 mV
\end{problembox}

\begin{problembox}
\textbf{Problem 625} \hfill \textbf{[Neeraj Kumar [NRJ] EXERCISE II (JEE ADVANCED) Q3]}
\label{q:ext-nrj_all_questions-625}

2Hg(l) + 2Fe3+             Hg
Fe
2
2
2
2
+
+
+
 
 
            excess    10--3 M
 
 Equ.                 10--3 -- x                x
2         x 
 
  
                  =
$/times$
-
10
100
10 3M 
 
   x = 9 $/times$ 10--4 
 
 Now, Keq =
$/times$
-
=
$/times$
$/times$
$/times$
$/times$
=
$/times$
-
-
-
-
-
x
x
x
2
10
9 10
2
9 10
1 10
9
10
2
2
3
2
4
4 2
4 2
3
4
(
)
(
)
(
)
 
 
 Now, E
E
E
K
cell
Fe
/Fe
Hg
/Hg
eq
3+
2
2+
$/circ$
=
$/circ$
-
$/circ$
= 0 06
2
.
log
 
 
 or, 0 7724
0 06
2
9
10
2
3
4
.
.
log
-
$/circ$
=
 
$/times$
-
E Hg
/Hg
2
2+
 
 
   E
V
Hg
/Hg
2
2+
$/circ$
= 0 815
.
\end{problembox}

\begin{problembox}
\textbf{Problem 626} \hfill \textbf{[Neeraj Kumar [NRJ] EXERCISE II (JEE ADVANCED) Q4]}
\label{q:ext-nrj_all_questions-626}

The cell reaction is 
 
 
6
14
6
2
7
3
3
Fe
Cr O
H
Fe
Cr
H O
2+
2
7
2
2
+
+
$/to$
+
+
-
+
+
+
 
 
 
E
E
n
cell
cell
=
$/circ$
-
 
+
+
+
+
+
0 06
3
6
3
2
2
6
2
7
2
8
.
log
[Fe
] [Cr
]
[Fe
] [Cr O
][H ]
 
 
         
=
-
-
$/times$
$/times$
$/times$
( .
.
)
.
log ( .
)
( )
( .
)
( )
1 35
0 77
0 06
6
0 75
4
0 75
2
1
6
2
6
8 
 
         = 0.571 V
\end{problembox}

\begin{problembox}
\textbf{Problem 627} \hfill \textbf{[Neeraj Kumar [NRJ] EXERCISE II (JEE ADVANCED) Q5]}
\label{q:ext-nrj_all_questions-627}

Cell reaction: H2(g) + 2Ag+ $/to$ 2H+ + 2Ag(s)
 
 
E
E
P
cell
cell
H2
=
$/circ$
-
 
+
+
0 06
2
2
2
.
log
[H ]
[Ag ]
 
 
 
Visit www.jeebooks.in for more
\end{problembox}

\begin{problembox}
\textbf{Problem 628} \hfill \textbf{[Neeraj Kumar [NRJ] EXERCISE II (JEE ADVANCED) Q6]}
\label{q:ext-nrj_all_questions-628}

E
K
cell
eq
$/circ$
= 0 06
.
log
n
 $/implies$ (0.2 -- 0.08) = 0 06
1
.
log Keq 
 
 $/implies$ Keq = 100
\end{problembox}

\begin{problembox}
\textbf{Problem 629} \hfill \textbf{[Neeraj Kumar [NRJ] EXERCISE II (JEE ADVANCED) Q7]}
\label{q:ext-nrj_all_questions-629}

Left electrode: Mn(s) $/to$ Mn2+ + 2e-- 
 
 Right electrode: 2H2O(l) $/to$ O2(g) + 4H+ + 4e-- 
 
    Net cell reaction: 2Mn (s) + 2H2O(l) $/to$ 2Mne+ 
+ O2(g) + 4H+ 
 
 Now, E
E
Mn
P
cell
cell
2+
O
=
$/circ$
-
 
 
+
0 06
4
1
2
4
2
.
log
[
]
[H ]
 
 
                       =
--
-
$/times$
$/times$
[ .
(
.
)]
.
log
( .
)
( .
)
.
1 229
1 185
0 06
4
0 001
0 01
0 25
1
2
4
 
 
                         = 2.633 V
\end{problembox}

\begin{problembox}
\textbf{Problem 630} \hfill \textbf{[Neeraj Kumar [NRJ] EXERCISE II (JEE ADVANCED) Q8]}
\label{q:ext-nrj_all_questions-630}

$/Delta$S
nF
E
T
P
=
  
 
 
  
 
  
=
$/times$
$/times$
=
2
96500
0 001
193
.
J/k
\end{problembox}

\begin{problembox}
\textbf{Problem 631} \hfill \textbf{[Neeraj Kumar [NRJ] EXERCISE II (JEE ADVANCED) Q9]}
\label{q:ext-nrj_all_questions-631}

$/Delta$
  $/Delta$
  $/Delta$
r
f
f
G
nFE
G
G
$/circ$ = -
$/circ$ =
$/circ$
-
$/circ$
Products
Reactants 
 
 or, -
$/times$
$/times$
=
$/times$
$/circ$
 
  
 
  
-
$/times$
+
$/times$
+
$/times$ -
+
$/times$
-
12
9600
2 5
1000
4
4
0
3
0
6
280
4
4
.
[
(
)
(OH)
$/Delta$fG Al
(
.
)]
-156 25
 
 
   $/Delta$fG$/circ$
= -
-
Al
KJ/mol
(OH)4
1300
\end{problembox}

\begin{problembox}
\textbf{Problem 632} \hfill \textbf{[Neeraj Kumar [NRJ] EXERCISE II (JEE ADVANCED) Q10]}
\label{q:ext-nrj_all_questions-632}

neq MnO2 = Q
F $/implies$ 8 7
87
1
3 99 10
96500
3
.
.
$/times$ =
$/times$
$/times$
-
t
 
   t = 2.418 $/times$ 106 sec $/approx$ 28 days
\end{problembox}

\begin{problembox}
\textbf{Problem 633} \hfill \textbf{[Neeraj Kumar [NRJ] EXERCISE II (JEE ADVANCED) Q11]}
\label{q:ext-nrj_all_questions-633}

SnCl2 
 Sn2+ + 2Cl--
 
 Cathode: Sn2+ + 2e-- 
 Sn
 
 Anode: 2Cl-- 
 Cl2 + 2e--
 
  
Cl2 + SnCl2 
 SnCl4
 
 Moles of SnCl2 taken = 19
190
0 1
= . 
 
 Moles of Sn produced = 1 19
119
0 01
.
.
=
 
                                   = moles of Cl2 produced
 
  
                           = moles of SnCl4 formed
 
 and moles of SnCl4 left = 0.1 -- (0.01 + 0.01) = 0.08
 
   
m
m
SnCl
SnCl
2
4
0 08 190
0 01 261
1520
261
=
$/times$
$/times$
=
.
.
\end{problembox}

\begin{problembox}
\textbf{Problem 634} \hfill \textbf{[Neeraj Kumar [NRJ] EXERCISE II (JEE ADVANCED) Q12]}
\label{q:ext-nrj_all_questions-634}

neq Au = Q
F $/implies$
$/times$
$/times$
$/times$
$/times$
=
$/times$
-
80 8 0 10
19 7 3
197
2 4
96 500
4
.
.
.
t 
 
   t = 772 sec
\end{problembox}

\begin{problembox}
\textbf{Problem 635} \hfill \textbf{[Neeraj Kumar [NRJ] EXERCISE II (JEE ADVANCED) Q13]}
\label{q:ext-nrj_all_questions-635}

AgBr(s)              Ag+      +                Br--
 
                                (10--7 + x)M              xM
 
 Now, (10--7 + x) x = 12 $/times$ 10--14 $/implies$ x = 3 $/times$ 10--7
 
 Final solution: [Ag+] = 4 $/times$ 10--7 M, [Br--] = 3 $/times$ 10--7 
M; [
]
NO
M
3
7
10
-
-
=
 
 Now, $/kappa$solution = $/lambda$$/circ$m (Ag+) $/times$ [Ag+] + $/lambda$$/circ$m(Br--) $/times$ 
[Br--] + $/lambda$$/circ$m (
)
[
]
NO
NO
3
3
-
-
$/times$
 
 
 =  6 $/times$ 10--3 $/times$ (4 $/times$ 10--7 $/times$ 103) + 8 $/times$10--3 $/times$ (3 $/times$ 10--7 $/times$ 
103) + 7 $/times$ 10--3 (10--7 $/times$ 10--3)
\end{problembox}

\begin{problembox}
\textbf{Problem 636} \hfill \textbf{[Neeraj Kumar [NRJ] EXERCISE II (JEE ADVANCED) Q14]}
\label{q:ext-nrj_all_questions-636}

C60 + 60O2 
 60 CO2
 
 Moles of O2 needed = 60
60
96
60
12
8
60
$/times$
=
$/times$
$/times$
=
nC
 
 
 neq O2 = neq azobenzene $/implies$ 8 $/times$ 4 = w
182
8
$/times$
 
 
 $/implies$ w = 728 gm
+ 8H+ + 8e--
+ 4H2O
N
N
NO2
2
\end{problembox}

\begin{problembox}
\textbf{Problem 637} \hfill \textbf{[Neeraj Kumar [NRJ] EXERCISE II (JEE ADVANCED) Q15]}
\label{q:ext-nrj_all_questions-637}

$/Lambda$$/circ$eq[Ba (PO ) ]
3
4 2 = 160 + 140 -- 100 
 
                          = 200 Ohm--1 cm2 eq--1
 
 Now,  $/circ$eq =  eq = $/kappa$
C $/implies$ 200
1 2
10
5
5
=
$/times$
-
.
 
   S = 6 $/times$ 10--8 eq/cm3 = 6 $/times$10--5 N = 10--5 M
 
   Ksp = 108 S5 = 108 $/times$ 10--25 M5
 
 
Visit www.jeebooks.in for more
\end{problembox}